% ============================================================
% M3 S4 拓展册（上）· 拓区收编（87 题）
% 依据：_tmpM3S4预备0913/S4波次方案.md §3.1（221 题次口径）＋canonical键名总表.md＋
%   定稿/定稿汇总-第2章.md（拓区账）＋_tmpM3S4题源闸0914/首跑钉码.md。
% 版式：练习线 qp-m3.sty（md5 7c3930362be8a0a2bdf21bbf8ac16573，本地挂载副本）；题后紧跟答案制；
%   全线括线模（\ansblockgrayfalse 导言一次置定——「>8 行→括线模」在全括线下自然满足）。
% 组织：节域序（canonical 01→…→17）→题型族组→组内席序；组名行只给 ≥2 成员族打。
% 跳号对号档：撤席（10-拓1／11-拓2／12-拓27／拓14-13／拓14-31／16-18）＝印面号空缺不重编号，
%   注记随席（tex 注释＋值台账抄件，不入印面）；章末 19 席缺料未收编（残余项，见汇编报告）。
% ============================================================
\documentclass[fontset=none]{ctexart}
\usepackage{qp-m3}
% —— 印面逐字符号路由补钉（承练习件母版；− 走 TNR、▱ NSC 子块；禁改 sty 保 md5） ——
\xeCJKDeclareCharClass{Default}{"2212}%
\setCJKfamilyfont{nscblk}[Path=C:/提示词/工作区/字替对照-0909/variantF/fonts/,BoldFont=NSC-w600.ttf]{NSC-w500.ttf}%
\xeCJKDeclareSubCJKBlock{nscblk}{"25B1}%
\newcommand{\pxparallelogram}{{\CJKfamily{nscblk}▱}}%
\xeCJKsetup{CJKglue={\hskip 0.10em plus 0.08em minus 0.05em}}%
\ansblockgrayfalse
\renewcommand{\qpzhangming}{第二章\quad 平面解析几何}%
\renewcommand{\qpjianming}{拓展册（上）}%
% —— 拓展册局部宏（组名行＝M2 拓展册档独立行；悬挂题号＝席号字面自适应宽，撤席跳号照印） ——
\newsavebox{\thbox}%
\newcommand{\tuhao}[1]{\par\glueguard{1}%
  \sbox{\thbox}{\textbf{#1.}}%
  \setlength{\lxhang}{\dimexpr\wd\thbox+1.4mm\relax}%
  \noindent\hangindent=\lxhang\hangafter=0
  \llap{\hbox to\lxhang{\textbf{#1.}\hss}}}%
\definecolor{gray102}{HTML}{666666}%
\newcommand{\zux}[1]{\par\glueguard{2}\addvspace{2.4mm}\noindent
  {\fontsize{10.5pt}{18pt}\selectfont\heibf 【#1】}\par\addvspace{1.2mm}}%

\begin{document}
%装配轮S6三本跨本连续页码（G7方案B）setcounter 注入位
\setcounter{page}{119}

\zhangtitle{拓展册（上）}

\begin{multicols}{2}
\emergencystretch=1em

\jietitle{2.1　坐标法（课时01·拓区）}
\zux{坐标法证明与存在性讨论（数轴绝对值距离不等式）}
\tuhao{T1} \tieside{难}在数轴上，对坐标分别为\(x_{1}\)和\(x_{2}\)的两点\(A\)和\(B\)，用绝对值定义两点间的距离，表示为 \(d(A,B)=|x_{1}-x_{2}|\)．
\((1)\) 在数轴上任意取三点\(A\)，\(B\)，\(C\)，证明 \(d(A,B)\leqslant d(A,C)+d(B,C)\)；
\((2)\) 设\(A\)和\(B\)两点的坐标分别为\(-3\)和\(2\)，分别找出\((1)\)中不等式等号成立和等号不成立时点\(C\)的范围．
\liubai[24mm]{解答书写区}
\begin{ansblock}[2章-拓-课时01-T1]
% ans:2章-拓-课时01-T1
\ansitem{T1}{\((1)\)见详解；\((2)\) 等号成立\(\Leftrightarrow C\)的坐标\(x\in [-3,2]\)；等号不成立\(\Leftrightarrow x<-3\)或\(x>2\)．}
\ansnote{详解}{\((1)\) 设\(C\)的坐标为\(x_{3}\)．因为 \(x_{1}-x_{2}=(x_{1}-x_{3})+(x_{3}-x_{2})\)，由绝对值三角不等式 \(d(A,B)=|x_{1}-x_{2}|=|(x_{1}-x_{3})+(x_{3}-x_{2})|\leqslant |x_{1}-x_{3}|+|x_{3}-x_{2}|=d(A,C)+d(B,C)\)，得证． \((2) |x_{1}-x_{2}|=|(x_{1}-x_{3})+(x_{3}-x_{2})|\)取等号\(\Leftrightarrow (x_{1}-x_{3})(x_{3}-x_{2})\geqslant 0\)，即\(x_{3}\)介于\(x_{1}\)与\(x_{2}\)之间（含端点）． 此处\(x_{1}=-3\)，\(x_{2}=2\)，故等号成立\(\Leftrightarrow -3\leqslant x\leqslant 2\)；等号不成立（即\(d(A,B)<d(A,C)+d(B,C)\)）\(\Leftrightarrow x<-3\)或\(x>2\)．}
\end{ansblock}

\zux{坐标法证明与存在性讨论（折线距离不等式与取等范围）}
\tuhao{T2} \tieside{难}城市的许多街道是相互垂直或平行的，往往不能沿直线行走到达目的地，只能按直角拐弯的方式行走．按照街道的垂直和平行方向建立平面直角坐标系，对两点\(A(x_{1},y_{1})\)和\(B(x_{2},y_{2})\)，定义两点间距离为 \(d(A,B)=|x_{1}-x_{2}|+|y_{1}-y_{2}|\)．
\((1)\) 在平面直角坐标系中任意取三点\(A\)，\(B\)，\(C\)，证明 \(d(A,B)\leqslant d(A,C)+d(B,C)\)；
\((2)\) 设\(A(x_{1},y_{1})\)，\(B(x_{2},y_{2})\)，分别找出\((1)\)中不等式等号成立和等号不成立时点\(C\)的范围．
\liubai[24mm]{解答书写区}
\begin{ansblock}[2章-拓-课时01-T2]
% ans:2章-拓-课时01-T2
\ansitem{T2}{\((1)\)见详解；\((2)\) 等号成立\(\Leftrightarrow C\)在以\(AB\)为对角线的矩形内（含边界），即\((x\_C-x_{1})(x\_C-x_{2})\leqslant 0\)且\((y\_C-y_{1})(y\_C-y_{2})\leqslant 0\)；等号不成立\(\Leftrightarrow C\)在该矩形外．}
\ansnote{详解}{\((1)\) 设\(C(x_{3},y_{3})\)．\(x_{1}-x_{2}=(x_{1}-x_{3})+(x_{3}-x_{2})\)，\(y_{1}-y_{2}=(y_{1}-y_{3})+(y_{3}-y_{2})\)，两次用绝对值三角不等式： \(d(A,B)=|x_{1}-x_{2}|+|y_{1}-y_{2}|\leqslant (|x_{1}-x_{3}|+|x_{3}-x_{2}|)+(|y_{1}-y_{3}|+|y_{3}-y_{2}|)=d(A,C)+d(B,C)\)，得证． \((2)\) 横、纵两个方向分别取等：第一个不等式取等\(\Leftrightarrow (x_{1}-x_{3})(x_{3}-x_{2})\geqslant 0\)，第二个取等\(\Leftrightarrow (y_{1}-y_{3})(y_{3}-y_{2})\geqslant 0\)， 故总等号成立\(\Leftrightarrow x_{3}\)介于\(x_{1}\)、\(x_{2}\)之间且\(y_{3}\)介于\(y_{1}\)、\(y_{2}\)之间，即\(C\)落在以\(AB\)为对角线（边平行于坐标轴）的矩形内部或边界上； 只要有一个方向不取等，就有\(d(A,B)<d(A,C)+d(B,C)\)，即\(C\)在该矩形外．}
\end{ansblock}

\jietitle{2.2　倾斜角与斜率（课时02·拓区）}
\zux{倾斜角与斜率（斜率随倾斜角变化说理）}
\tuhao{T1} \tieside{中}\((1)\) 如果直线的倾斜角\(\theta \in [0,\pi /2)\)，则当\(\theta \)增大时，直线的斜率将怎样变化？如果\(\theta \in (\pi /2,\pi )\)呢？
\((2)\) 能否说\("\)直线的倾斜角增大时斜率也增大\("\)？为什么？
\liubai[24mm]{解答书写区}
\begin{ansblock}[2章-拓-课时02-T1]
% ans:2章-拓-课时02-T1
\ansitem{T1}{\((1) \theta \in [0,\pi /2)\)时斜率由\(0\)递增并趋向\(+\infty \)；\(\theta \in (\pi /2,\pi )\)时斜率由\(-\infty \)递增趋向\(0\)；\((2)\) 不能，理由见详解．}
\ansnote{详解}{\((1) k=tan\theta \)在\([0,\pi /2)\)上单调递增，\(k\)从\(0\)增大到正无穷大；在\((\pi /2,\pi )\)上也单调递增，\(k\)从负无穷大增大到\(0\)． \((2)\) 不能．例如取\(\theta _{1}=45^{\circ}\)，\(\theta _{2}=120^{\circ}\)，\(\theta _{1}<\theta _{2}\)，但\(k_{1}=1>k_{2}=-\sqrt{3}\)，即倾斜角大者斜率反而小； 斜率关于倾斜角只在每一段\([0,\pi /2)\)与\((\pi /2,\pi )\)内分别递增，整体上不具有单调性．}
\end{ansblock}

\zux{倾斜角与斜率（三点斜率倾斜角与旋转范围）}
\tuhao{T2} \tieside{中}已知坐标平面内三点\(A(-1,1)\)，\(B(1,1)\)，\(C(2,\sqrt{3}+1)\)．
\((1)\) 求直线\(AB\)，\(BC\)，\(AC\)的斜率和倾斜角；
\((2)\) 若\(D\)为\(\triangle ABC\)的\(AB\)边上一动点，求直线\(CD\)的倾斜角的取值范围．
\liubai[24mm]{解答书写区}
\begin{ansblock}[2章-拓-课时02-T2]
% ans:2章-拓-课时02-T2
\ansitem{T2}{\((1) k\_AB=0\)（倾斜角\(0\)）；\(k\_BC=\sqrt{3}\)（倾斜角\(\pi /3\)）；\(k\_AC=\sqrt{3}/3\)（倾斜角\(\pi /6\)）；\((2) [\pi /6, \pi /3]\)．}
\ansnote{详解}{\((1) k\_AB=(1-1)/(1+1)=0\)，\(k\_BC=(\sqrt{3}+1-1)/(2-1)=\sqrt{3}\)，\(k\_AC=(\sqrt{3}+1-1)/(2+1)=\sqrt{3}/3\)． 因斜率等于倾斜角的正切值且倾斜角\(\alpha \in [0,\pi )\)，故\(AB\)、\(BC\)、\(AC\)的倾斜角依次为\(0\)、\(\pi /3\)、\(\pi /6\)． \((2)\) 当直线\(CD\)绕点\(C\)由\(CA\)逆时针旋转到\(CB\)时，\(CD\)与线段\(AB\)恒有交点（即\(D\)在\(AB\)上），此时\(k\_CD\)由\(k\_AC\)增大到\(k\_BC\)， \(k\_CD\in [\sqrt{3}/3, \sqrt{3}]\)，对应倾斜角\(\alpha \in [\pi /6, \pi /3]\)．}
\end{ansblock}

\zux{倾斜角与斜率（含参两点求倾斜角三问）}
\tuhao{T3} \tieside{中}已知\(a\)，\(b\)，\(c\)是两两不相等的实数，分别求经过下列两点的直线的倾斜角：
\((1) A(a,c)\)，\(B(b,c)\)；　\((2) C(a,b)\)，\(D(a,c)\)；　\((3) P(b,b+c)\)，\(Q(a,c+a)\)．
\liubai[16mm]{解答书写区}
\begin{ansblock}[2章-拓-课时02-T3]
% ans:2章-拓-课时02-T3
\ansitem{T3}{\((1) 0^{\circ}\)；\((2) 90^{\circ}\)；\((3) 45^{\circ}\)．}
\ansnote{详解}{\((1)\) 纵坐标同为\(c\)，直线与\(x\)轴平行，倾斜角为\(0^{\circ}\)（\(k=0\)）； \((2)\) 横坐标同为\(a\)，直线垂直于\(x\)轴，斜率不存在，倾斜角为\(90^{\circ}\)； \((3) k=[(c+a)-(b+c)]/(a-b)=(a-b)/(a-b)=1\)（\(a\neq b\)），\(tan\alpha =1\)，\(\alpha \in [0^{\circ},180^{\circ})\)，故倾斜角为\(45^{\circ}\)．}
\end{ansblock}

\zux{倾斜角与斜率（斜率几何意义求比值范围）}
\tuhao{T4} \tieside{难}点\(M(x,y)\)在函数\(y=-2x+8\)的图象上，当\(x\in [2,5]\)时，求\((y+1)/(x+1)\)的取值范围．
\liubai[16mm]{解答书写区}
\begin{ansblock}[2章-拓-课时02-T4]
% ans:2章-拓-课时02-T4
\ansitem{T4}{\([-1/6, 5/3]\)}
\ansnote{详解}{\((y+1)/(x+1)=[y-(-1)]/[x-(-1)]\) 的几何意义是过\(M(x,y)\)与定点\(N(-1,-1)\)两点的直线的斜率，其中点\(M\)在线段\(y=-2x+8\)（\(x\in [2,5]\)）上运动． 线段两端点为\((2,4)\)与\((5,-2)\)．记\(k(x)=(y+1)/(x+1)=(-2x+9)/(x+1)=-2+11/(x+1)\)，在\(x\in [2,5]\)上随\(x\)增大而单调递减， 故最大值在\(x=2\)处取得：\(k=(4+1)/(2+1)=5/3\)；最小值在\(x=5\)处取得：\(k=(-2+1)/(5+1)=-1/6\)． 所以\((y+1)/(x+1)\)的取值范围为\([-1/6, 5/3]\)．}
\end{ansblock}

\jietitle{2.2　点斜式与斜截式（课时04·拓区）}
\zux{直线方程形式与互化（含参恒过定点与面积最值·两问）}
\tuhao{T1} \tieside{中偏难}已知直线\(l\)：\(kx-y+4k+2=0\)（\(k\in R\)）．
\((1)\) 若直线\(l\)不经过第四象限，求\(k\)的取值范围；
\((2)\) 若直线\(l\)交\(x\)轴的负半轴于点\(A\)，交\(y\)轴的正半轴于点\(B\)，\(O\)为坐标原点，设\(\triangle AOB\)的面积为\(S\)，求\(S\)的最小值及此时直线\(l\)的方程．
\liubai[24mm]{解答书写区}
\begin{ansblock}[2章-拓-课时04-T1]
% ans:2章-拓-课时04-T1
\ansitem{T1}{\((1) k\geqslant 0\)；\((2) S\)的最小值为\(16\)，此时直线\(l\)的方程为\(x-2y+8=0\)}
\ansnote{详解}{\((1)\) 将方程整理为\(k(x+4)+(-y+2)=0\)，令\(x+4=0\)且\(y=2\)，得直线恒过定点\((-4,2)\)（在第二象限）． 斜率\(k=tan\theta \)．过\((-4,2)\)的直线不经过第四象限，当且仅当它\("\)向右上或水平张开\("\)：由\(l_{1}\)（过\((-4,2)\)与原点方向的临界位置）逆时针转到竖直位置均不进入第四象限，即倾斜角\(\theta \in [0^{\circ},90^{\circ}]\)，故\(k\geqslant 0\)（\(k\)不存在时直线\(x=-4\)也不过第四象限，但它不在\(y=kx+b\)的参数范围内——题设方程本身取遍\(k\in R\)时不含竖直线，而\(x=-4\)无法由该方程表示，故不另计）； \(k<0\)时直线向右下倾斜，自第二象限的点\((-4,2)\)出发必穿过第四象限，不合．所以\(k\geqslant 0\)． \((2)\) 由题意\(A\)、\(B\)存在且\(a<0\)、\(b>0\)（竖直线\(x=-4\)不与\(y\)轴相交，此时\(k=0\)给\(y=2\)与\(x\)轴不相交，均已在端点情形排除，直线可设为截距式\(x/a+y/b=1\)，\(a<0\)，\(b>0\)）． \((-4,2)\)在直线上：\(-4/a+2/b=1\)．因\(a<0\)，\(b>0\)，\(-4/a>0\)，\(2/b>0\)， \(1=-4/a+2/b\geqslant 2\sqrt{}((-4/a)\cdot (2/b))=2\sqrt{}(-8/(ab))\)，两边平方得\(1\geqslant -32/(ab)\)，即\(-ab\geqslant 32\)，\(ab\leqslant -32\)． \(S=(1/2)|OA|\cdot |OB|=(1/2)(-a)b=-ab/2\geqslant 16\)， 当且仅当\(-4/a=2/b=1/2\)，即\(a=-8\)，\(b=4\)时取等号． 此时直线为\(x/(-8)+y/4=1\)，化简得\(x-2y+8=0\)． 所以\(S\)的最小值为\(16\)，此时直线\(l\)的方程为\(x-2y+8=0\)．}
\end{ansblock}

\jietitle{2.2　两点式与一般式（课时05·拓区）}
\zux{直线方程形式与互化（正半轴截距与面积·两问）}
\tuhao{T6} \tieside{中}过点\(P(1,4)\)作直线\(l\)，直线\(l\)与\(x\)、\(y\)轴的正半轴分别交于\(A\)，\(B\)两点，\(O\)为原点．
\((1)\) 若\(\triangle ABO\)的面积为\(9\)，求直线\(l\)的方程；\((2)\) 若\(\triangle ABO\)的面积为\(S\)，求\(S\)的最小值，并求出此时直线\(l\)的方程．
\liubai[16mm]{解答书写区}
\begin{ansblock}[2章-拓-课时05-T6]
% ans:2章-拓-课时05-T6
\ansitem{T6}{\((1) 2x+y-6=0\)或\(8x+y-12=0\)；\((2) 8\)，\(4x+y-8=0\)}
\ansnote{详解}{设\(A(a,0)\)、\(B(0,b)\)，\(a>0\)，\(b>0\)，\(l\)：\(x/a+y/b=1\)，代入\(P(1,4)\)得\(1/a+4/b=1\)（\(*\)）． \((1)\) 又\((1/2)ab=9\)即\(ab=18\)；由\((*)\)去分母\(b+4a=ab=18\)，\(b=18-4a\)，\(a(18-4a)=18\)，\(2a^{2}-9a+9=0\)，\(a=3\)或\(3/2\)； 对应\(b=6\)或\(12\)，\(l\)为\(x/3+y/6=1\)或\(x/(3/2)+y/12=1\)，即\(2x+y-6=0\)或\(8x+y-12=0\)． \((2) S=(1/2)ab=(1/2)ab\cdot (1/a+4/b)^{2}=(1/2)(8+b/a+16a/b)\geqslant (1/2)(8+2\sqrt{b/a\cdot 16a/b})=8\)， 当且仅当\(b/a=16a/b\)即\(b=4a\)，代\((*)\)：\(1/a+1/a=1\)，\(a=2\)、\(b=8\)时取等，\(S\)最小值\(8\)，此时\(l\)：\(x/2+y/8=1\)即\(4x+y-8=0\)．}
\end{ansblock}

\zux{直线方程形式与互化（截距积与距离积最值·两问）}
\tuhao{T7} \tieside{中}过点\(P(2,1)\)作直线\(l\)分别交\(x\)轴、\(y\)轴的正半轴于\(A\)，\(B\)两点．
\((1)\) 当\(|OA|\cdot |OB|\)取最小值时，求出最小值及直线\(l\)的截距式方程；
\((2)\) 当\(|PA|\cdot |PB|\)取最小值时，求出最小值及直线\(l\)的截距式方程．
\liubai[24mm]{解答书写区}
\begin{ansblock}[2章-拓-课时05-T7]
% ans:2章-拓-课时05-T7
\ansitem{T7}{\((1)\) 最小值\(8\)，\(x/4+y/2=1\)；\((2)\) 最小值\(4\)，\(x/3+y/3=1\)}
\ansnote{详解}{\((1)\) 设\(l\)：\(x/a+y/b=1\)（\(a>0\)，\(b>0\)），过\(P\)得\(2/a+1/b=1\)． \(1=2/a+1/b\geqslant 2\sqrt{}(2/(ab))\)，得\(ab\geqslant 8\)，当且仅当\(2/a=1/b\)即\(a=4\)、\(b=2\)时取等； \(|OA|\cdot |OB|=ab\)的最小值为\(8\)，此时\(l\)：\(x/4+y/2=1\)． \((2)\) 由\((1)\)得\(b=a/(a-2)\)（\(a>2\)）．\(A(a,0)\)、\(B(0,b)\)， \(|PA|=\sqrt{}((a-2)^{2}+1)\)，\(|PB|=\sqrt{}(4+(b-1)^{2})\)，而\(b-1=a/(a-2)-1=2/(a-2)\)， \(|PA|\cdot |PB|=\sqrt{}((a-2)^{2}+1)\cdot \sqrt{}(4+4/(a-2)^{2})=2\sqrt{}((a-2)^{2}+1/(a-2)^{2}+2)\geqslant 2\sqrt{2+2}=4\)， 当且仅当\((a-2)^{2}=1/(a-2)^{2}\)即\(a-2=1\)、\(a=3\)时取等，此时\(b=3\)； \(|PA|\cdot |PB|\)的最小值为\(4\)，此时\(l\)：\(x/3+y/3=1\)．}
\end{ansblock}

\zux{直线方程形式与互化（三条件求直线·源误改定）}
\tuhao{T8} \tieside{中}求分别满足下列条件的直线\(l\)的方程：
\((1)\) 斜率是\(3/4\)，且与两坐标轴围成的三角形的面积是\(6\)；
\((2)\) 经过两点\(A(1,0)\)，\(B(m,1)\)；
\((3)\) 经过点\((4,-3)\)，且在两坐标轴上的截距的绝对值相等．
\liubai[32mm]{解答书写区}
\begin{ansblock}[2章-拓-课时05-T8]
% ans:2章-拓-课时05-T8
\ansitem{T8}{\((1) y=(3/4)x\pm 3\)；\((2) m\neq 1\)时\(y=(1/(m-1))(x-1)\)，\(m=1\)时\(x=1\)；\((3) x+y-1=0\)或\(x-y-7=0\)或\(y=-(3/4)x\)}
\ansnote{详解}{\((1)\) 设\(l\)：\(y=(3/4)x+b\)，则\(y\)轴截距\(b\)，\(x\)轴截距\(-4b/3\)，面积\((1/2)|b\cdot (-4b/3)|=2b^{2}/3=6\)，\(b^{2}=9\)，\(b=\pm 3\)，\(l\)：\(y=(3/4)x\pm 3\)；回代验证：\(b=3\) 代回，面积\((1/2)\cdot |3\cdot (-4)|=6\) ；\(b=-3\) 代回，面积\((1/2)\cdot |(-3)\cdot 4|=6\) ，双向验证成立； \((2) m\neq 1\)时两点横、纵坐标均不等，斜率\(k=(1-0)/(m-1)=1/(m-1)\)，点斜式\(y=(1/(m-1))(x-1)\)；\(m=1\)时\(A\)、\(B\)同在直线\(x=1\)上，\(l\)：\(x=1\)； \((3)\) 设截距\(a\)、\(b\)．①过原点（\(a=b=0\)）：\(l\)过\((0,0)\)、\((4,-3)\)，\(y=-(3/4)x\)； ②不过原点，\(x/a+y/b=1\)，\(|a|=|b|\neq 0\)且\(4/a-3/b=1\)：\(a=b\)时\(1/a=1\)、\(a=1\)（\(x+y-1=0\)）；\(a=-b\)时\(4/a+3/a=1\)、\(a=7\)（\(x-y-7=0\)）． 共\(3\)条． （按源卷核对注：源答案\((1)\)写作\("y=(4/3)x\pm 3"\)，与其题面\("\)斜率\(3/4"\)不符——本卷按题面取修正值 \(y=(3/4)x\pm 3\)；面积回代双向验证已并入\((1)\)详解正文 ，详见亲算账④\("\)源误更正\("\)。）}
\end{ansblock}

\jietitle{2.4　曲线与方程（课时10·拓区）}
% 【10-拓1｜撤位登记（案2令1：与G5·单1 同题双列撤席让位，唯一活位＝10-G1）】本席不设题——跳号留痕（拓2 起），注记随席（台账抄件）；防双收 B22。
\zux{定义法轨迹＋椭圆定直线（重心模型）}
% 席注（账面）：⚠前引观察（同G2）：详解用椭圆定义/方程与直线联立韦达（2.5.1后知识），S4装配建议课时11后启用；设x=my−3免斜率讨论技巧照源。
\tuhao{02} \tieside{难}已知\(B(-1,0)\)、\(C(1,0)\)为\(\triangle ABC\)的两个顶点，\(P\)为\(\triangle ABC\)的重心，边\(AC\)、\(AB\)上的两条中线长度之和为\(6\)．
\((1)\)求点\(P\)的轨迹\(T\)的方程；
\((2)\)已知点\(N(-3,0)\)、\(E(-2,0)\)、\(F(2,0)\)，直线\(PN\)与曲线\(T\)的另一个公共点为\(Q\)，直线\(EP\)与\(FQ\)交于点\(M\)，试问：当点\(P\)变化时，点\(M\)是否恒在一条定直线上？若是，请证明；若不是，请说明理由．
\liubai[24mm]{解答书写区}
\begin{ansblock}[2章-拓-课时10-02]
% ans:2章-拓-课时10-02
\ansitem{02}{\((1)x^{2}/4+y^{2}/3=1\)（\(x\neq \pm 2\)）；\((2)\)点\(M\)恒在定直线\(x=-4/3\)上}
\ansnote{详解}{\((1)\)重心\(P\)分中线为\(2:1\)，故 \(|PB|+|PC|=(2/3)\times 6=4>|BC|=2\)，\(P\)的轨迹是以\(B\)、\(C\)为焦点的椭圆（\(A\)与\(BC\)共线时退化，除去长轴端点），\(a=2\)、\(c=1\)、\(b=\sqrt{3}\)：\(x^{2}/4+y^{2}/3=1\)（\(x\neq \pm 2\)）。 \((2)\)设\(PQ\)：\(x=my-3\)，\(P(x_{1},y_{1})\)、\(Q(x_{2},y_{2})\)。联立 \((3m^{2}+4)y^{2}-18my+15=0\)，\(y_{1}+y_{2}=18m/(3m^{2}+4)\)、\(y_{1}y_{2}=15/(3m^{2}+4)\)，得 \(2my_{1}y_{2}=(5/3)(y_{1}+y_{2})\)。直线\(EP\)：\(y=y_{1}/(my_{1}-1)\cdot (x+2)\)；直线\(FQ\)：\(y=y_{2}/(my_{2}-5)\cdot (x-2)\)。联立解得 \(x=2(2my_{1}y_{2}-y_{2}-5y_{1})/(-y_{2}+5y_{1})\)，代入 \(2my_{1}y_{2}=(5/3)(y_{1}+y_{2})\)：\(x=2[(2/3)y_{2}-(10/3)y_{1}]/(-y_{2}+5y_{1})=(4/3)(y_{2}-5y_{1})/(-y_{2}+5y_{1})=-4/3\) 恒成立。故\(M\)恒在定直线 \(x=-4/3\) 上。}
\end{ansblock}

\zux{垂足轨迹·隐圆（定义法）}
% 席注（账面）：⚠照录含后节载体（抛物线焦点，2.7知识，§5"抛物线载体先例照录"）；源详解含图（图注从源）。
\tuhao{03} \tieside{中}已知过抛物线\(y^{2}=4x\)焦点\(F\)的直线交抛物线于\(A\)、\(B\)两点，过原点\(O\)作\(OM\)，使\(OM\perp AB\)，垂足为点\(M\)，求点\(M\)的轨迹方程．
\liubai[16mm]{解答书写区}
\begin{ansblock}[2章-拓-课时10-03]
% ans:2章-拓-课时10-03
\ansitem{03}{\((x-1/2)^{2}+y^{2}=1/4\)}
\ansnote{详解}{\(F(1,0)\)。\(OM\perp AB\) 且垂足\(M\)在过\(F\)的动直线\(AB\)上 \(\Rightarrow  \angle OMF=90^{\circ}\)（\(M=F\) 时直线\(AB\)为\(x=1\)亦合），故\(M\)在以\(OF\)为直径的圆上：圆心\((1/2,0)\)、半径\(1/2\)，方程 \((x-1/2)^{2}+y^{2}=1/4\)。}
\end{ansblock}

\zux{垂足轨迹·齐次化（隐圆）}
% 席注（账面）：源含常规法（设y=kx+m韦达）题后反思版，照录从略；a²=2b²代入口径照源【编注】易错点提示。
\tuhao{04} \tieside{中}\(Q_{1}\)、\(Q_{2}\)为椭圆 \(x^{2}/(2b^{2})+y^{2}/b^{2}=1\) 上的两个动点，且\(OQ_{1}\perp OQ_{2}\)，过原点\(O\)作直线\(Q_{1}Q_{2}\)的垂线\(OD\)，求\(D\)的轨迹方程．
\liubai[16mm]{解答书写区}
\begin{ansblock}[2章-拓-课时10-04]
% ans:2章-拓-课时10-04
\ansitem{04}{\(x^{2}+y^{2}=(2/3)b^{2}\)〔图注：题面示意照录自源〕}
\ansnote{详解}{设弦\(Q_{1}Q_{2}\)：\(mx+ny=1\)（不过原点可如此设）。齐次化：\(x^{2}/(2b^{2})+y^{2}/b^{2}-(mx+ny)^{2}=0\)，展开乘\(2b^{2}\)：\((2-2b^{2}n^{2})y^{2}+(1-2b^{2}m^{2})x^{2}-4mnb^{2}xy=0\)，两边除\(x^{2}\)得关于\(k=y/x\)的方程 \((2-2b^{2}n^{2})k^{2}-4mnb^{2}k+1-2b^{2}m^{2}=0\)，两根即\(OQ_{1}\)、\(OQ_{2}\)斜率，垂直 \(\Rightarrow  k_{1}k_{2}=(1-2b^{2}m^{2})/(2-2b^{2}n^{2})=-1 \Rightarrow  1-2b^{2}m^{2}=-2+2b^{2}n^{2} \Rightarrow  2b^{2}(m^{2}+n^{2})=3\)（\(\ast \)）。设\(D(x_{0},y_{0})\)，\(OD\perp Q_{1}Q_{2}\) 且\(D\)在弦上：直线\(Q_{1}Q_{2}\)即 \(x_{0}x+y_{0}y=x_{0}^{2}+y_{0}^{2}\)（过\(D\)垂直于\(OD\)），对照 \(mx+ny=1\) 得 \(m=x_{0}/(x_{0}^{2}+y_{0}^{2})\)、\(n=y_{0}/(x_{0}^{2}+y_{0}^{2})\)，代入（\(\ast \)）：\(2b^{2}(x_{0}^{2}+y_{0}^{2})/(x_{0}^{2}+y_{0}^{2})^{2}=3 \Rightarrow  x_{0}^{2}+y_{0}^{2}=(2/3)b^{2}\)。\(D\)轨迹为 \(x^{2}+y^{2}=(2/3)b^{2}\)。验算：\(a^{2}=2b^{2}\)、\(c^{2}=b^{2}\)，原点到\("\)斜距\("\)型结论 \(|OD|^{2}=a^{2}b^{2}/(a^{2}+b^{2})=2b^{4}/(3b^{2})=(2/3)b^{2}\) （椭圆切线对偶结构，此处仅作后台校核不作讲解）。}
\end{ansblock}

\zux{向量条件轨迹（中点伸位）}
% 席注（账面）：⚠存疑：坐标消参得x²+(y−1)²=1与几何位似（半径1/2×2=1、圆心(0,1/2)×2=(0,1)）两路同果✓；竖直弦情形P=(0,0)在圆上✓（k→∞极限）。
\tuhao{05} \tieside{中}已知过点\((0,1)\)的直线与圆\(x^{2}+y^{2}=4\)相交于\(A\)、\(B\)两点，若 向量\(OA+\)向量\(OB=\)向量\(OP\)，则点\(P\)的轨迹方程是（　）
\bindopt {\fontsize{10.5pt}{\qpTJgLeadLiti}\selectfont \makebox[\dimexpr0.5\linewidth-0.5\lxhang-0.25em\relax][l]{\(A\)．\(x^{2}+(y-1/2)^{2}=1\)}\makebox[\dimexpr0.5\linewidth-0.5\lxhang-0.25em\relax][l]{\(B\)．\(x^{2}+(y-1)^{2}=1\)}\par}
\bindopt {\fontsize{10.5pt}{\qpTJgLeadLiti}\selectfont \makebox[\dimexpr0.5\linewidth-0.5\lxhang-0.25em\relax][l]{\(C\)．\(x^{2}+(y-1/2)^{2}=2\)}\makebox[\dimexpr0.5\linewidth-0.5\lxhang-0.25em\relax][l]{\(D\)．\(x^{2}+(y-1)^{2}=2\)}\par}
\begin{ansblock}[2章-拓-课时10-05]
% ans:2章-拓-课时10-05
\ansitem{05}{\(B\)}
\ansnote{详解}{设过\((0,1)\)的直线 \(y=kx+1\)（竖直直线\(x=0\)时\(A(0,2)\)、\(B(0,-2)\)，\(OP=OA+OB=(0,0)\)，\(P=(0,0)\)适合下文方程），代入圆：\((1+k^{2})x^{2}+2kx-3=0\)，\(x_{1}+x_{2}=-2k/(1+k^{2})\)，\(y_{1}+y_{2}=k(x_{1}+x_{2})+2=2/(1+k^{2})\)。由向量\(OA+\)向量\(OB=\)向量\(OP\) 得 \(P(x,y)=(x_{1}+x_{2}, y_{1}+y_{2})=(-2k/(1+k^{2}), 2/(1+k^{2}))\)。消\(k\)：\(y-1=(2-(1+k^{2}))/(1+k^{2})=(1-k^{2})/(1+k^{2})\)，于是 \(x^{2}+(y-1)^{2}=[4k^{2}+(1-k^{2})^{2}]/(1+k^{2})^{2}=(1+2k^{2}+k^{4})/(1+k^{2})^{2}=(1+k^{2})^{2}/(1+k^{2})^{2}=1\)，即 \(x^{2}+(y-1)^{2}=1\)。几何法（源【编注】口径）：设\(AB\)中点\(D\)，\(OP=2OD\)（向量），且\(OD\perp AB\)、\(D\)在以\(O(0,0)\)与\((0,1)\)为直径端点的圆 \(x^{2}+(y-1/2)^{2}=1/4\) 上，以\(O\)为中心\(2\)倍伸位即得 \(x^{2}+(y-1)^{2}=1\)。选\(B\)。}
\end{ansblock}

\zux{交轨法（斜率积定值）}
% 席注（账面）：与课内新3（斜率积−1/4）方法同族而载体异（此交轨彼直接法），组合复核由S4按§5三件口径把握。
\tuhao{06} \tieside{中}求两动直线\(y=kx+1\)与\(y=-(2/k)x-1\)的交点\(P\)的轨迹方程．
\liubai[16mm]{解答书写区}
\begin{ansblock}[2章-拓-课时10-06]
% ans:2章-拓-课时10-06
\ansitem{06}{\(x^{2}/(1/2)+y^{2}=1\)（\(x\neq 0\)），即 \(2x^{2}+y^{2}=1\)（\(x\neq 0\)）}
\ansnote{详解}{\(l_{1}\)过定点\(A(0,1)\)、斜率\(k_{1}=k\)；\(l_{2}\)过定点\(B(0,-1)\)、斜率\(k_{2}=-2/k\)，\(k_{1}\cdot k_{2}=-2\)为定值。设\(P(x,y)\)，两斜率存在（\(k\neq 0\)）故 \(x\neq 0\)，\(k_{1}=(y-1)/x\)、\(k_{2}=(y+1)/x\)：\((y-1)(y+1)/x^{2}=-2 \Rightarrow  y^{2}-1=-2x^{2} \Rightarrow  2x^{2}+y^{2}=1\)（\(x\neq 0\)）。}
\end{ansblock}

\zux{交轨法（动线段载体）}
% 席注（账面）：源详解含第二图（轨迹示意，图注从源）；本区双曲线形状语句照录成品（判定语，不涉2.6推导）。　图债注：「如图」指源件配图，印面无图（冻片【注】裁定：去装饰图/条件全在文字/尺寸随注——图债登记汇编报告）
\tuhao{07} \tieside{中}已知点\(P(-2,2)\)、\(Q(0,2)\)以及直线\(l\)：\(y=x\)，设长为\(\sqrt{2}\)的线段\(AB\)在直线\(l\)上移动（如图所示），求直线\(PA\)和\(QB\)的交点\(M\)的轨迹方程．〔图注：照录自源〕
\liubai[16mm]{解答书写区}
\begin{ansblock}[2章-拓-课时10-07]
% ans:2章-拓-课时10-07
\ansitem{07}{\((y+1)^{2}/8-(x+1)^{2}/8=1\)}
\ansnote{详解}{设\(A(a,a)\)、\(B(b,b)\)（\(b>a\)）、\(M(x,y)\)。\(P\)、\(A\)、\(M\)共线：\((x+2)/(y-2)=(a+2)/(a-1) \Rightarrow  a=2(x+y)/(x-y+4)\)（\(y\neq 2\)）；\(Q\)、\(B\)、\(M\)共线：\(x/(y-2)=b/(b-2) \Rightarrow  b=2x/(x-y+2)\)。\(|AB|=\sqrt{2}(b-a)=\sqrt{2} \Rightarrow  b=a+1\)。代入整理：\(2x/(x-y+2)=2(x+y)/(x-y+4)+1 \Rightarrow  x^{2}-y^{2}+2x-2y+8=0\)。\(y=2\)时交点\(M\)与\(P\)或\(Q\)重合，两点坐标均满足该方程，故方程即轨迹方程；配方 \((y+1)^{2}/8-(x+1)^{2}/8=1\)，轨迹为中心\((-1,-1)\)、焦点在直线\(x=-1\)上的双曲线。}
\end{ansblock}

\zux{两线交点轨迹（交轨·隐式）＋定向直线存在性}
% 席注（账面）：⚠照录含后节载体（轨迹为双曲线之判定表述2.6.1；本题只需方程与代数操作，装配可注"曲线名称后置"）。
\tuhao{08} \tieside{难}圆\(O\)：\(x^{2}+y^{2}=4\)与\(x\)轴的负半轴和正半轴分别交于\(A\)、\(B\)两点，\(MN\)是圆与\(x\)轴垂直非直径的弦，直线\(AM\)与直线\(BN\)交于点\(P\)，记动点\(P\)的轨迹为\(E\)．
\((1)\)求轨迹\(E\)的方程；
\((2)\)在平面直角坐标系中，倾斜角确定的直线称为定向直线．是否存在不过点\(A\)的定向直线\(l\)，当直线\(l\)与轨迹\(E\)交于\(C\)、\(D\)时，\(AC\perp AD\)；若存在，求直线\(l\)的一个方向向量；若不存在，说明理由．
\liubai[24mm]{解答书写区}
\begin{ansblock}[2章-拓-课时10-08]
% ans:2章-拓-课时10-08
\ansitem{08}{\((1)x^{2}-y^{2}=4\)（\(x\neq \pm 2\)）；\((2)\)存在，方向向量\((1,0)\)}
\ansnote{详解}{\((1)A(-2,0)\)、\(B(2,0)\)。设\(M(m,n)\)、\(N(m,-n)\)（\(m\neq 0\)，非直径），\(m^{2}+n^{2}=4\)。\(AM\)：\(y=n/(m+2)\cdot (x+2)\)；\(BN\)：\(y=-n/(m-2)\cdot (x-2)\)。两式相乘：\(y^{2}=-n^{2}/(m^{2}-4)\cdot (x^{2}-4)=(-n^{2}/-n^{2})(x^{2}-4)=x^{2}-4\)，即 \(x^{2}-y^{2}=4\)；\(m\neq 0\)排除\(M\)、\(N\)在\(x\)轴上（\(x\neq \pm 2\)）。 \((2)\)①\(l\)倾斜角\(90^{\circ}\)：设\(x=t\)（\(|t|>2\)），\(C\)、\(D=(t,\pm \sqrt{t^{2}-4})\)，\(AC\cdot AD=(t+2)^{2}-(t^{2}-4)=8t+8=0 \rightarrow  t=-1\)与\(|t|>2\)矛盾，无解。②\(l\)不竖直：设\(y=kx+t\)，联立 \((1-k^{2})x^{2}-2ktx-t^{2}-4=0\)（\(1-k^{2}\neq 0\)，\(\Delta >0\)），\(x_{1}+x_{2}=2kt/(1-k^{2})\)、\(x_{1}x_{2}=(-t^{2}-4)/(1-k^{2})\)。\(AC\perp AD\)：\((x_{1}+2)(x_{2}+2)+(kx_{1}+t)(kx_{2}+t)=0 \Rightarrow  (1+k^{2})x_{1}x_{2}+(2+kt)(x_{1}+x_{2})+t^{2}+4=0\)，代入化简得 \(k(t-2k)=0\)：\(k=0\)（恒合，验证\(\Delta =4(t^{2}+4)>0\)且\(1-0\neq 0\) ）或 \(t=2k\)（\(l\)：\(y=k(x+2)\)恒过\(A(-2,0)\)，舍）。综上存在，\(k=0\)，一个方向向量为\((1,0)\)。}
\end{ansblock}

\zux{相关点法·中点双问（多问计1）}
% 席注（账面）：相关点族第3件——与G1、10-中7三件三取二口径照§5（S4复核）〔0913：原中7（-3）置换移入拓15，本族＝正文1（G1）＋拓展2（拓15/本席）〕。
\tuhao{09} \tieside{中}已知圆\(O\)：\(x^{2}+y^{2}=4\)上的一定点\(A(2,0)\)，点\(B(1,1)\)为圆内一点，\(P\)、\(Q\)为圆上的动点．
\((1)\)求线段\(AP\)中点的轨迹方程；\((2)\)若\(\angle PBQ=90^{\circ}\)，求线段\(PQ\)中点的轨迹方程．
\liubai[16mm]{解答书写区}
\begin{ansblock}[2章-拓-课时10-09]
% ans:2章-拓-课时10-09
\ansitem{09}{\((1)(x-1)^{2}+y^{2}=1\)；\((2)(x-1/2)^{2}+(y-1/2)^{2}=3/2\)}
\ansnote{详解}{\((1)\)设\(P(x_{0},y_{0})\)（\(x_{0}^{2}+y_{0}^{2}=4\)）、\(AP\)中点\(M(x,y)\)：\(x=(2+x_{0})/2\)、\(y=y_{0}/2 \Rightarrow  x_{0}=2x-2\)、\(y_{0}=2y\)，代入：\((2x-2)^{2}+(2y)^{2}=4\)，即 \((x-1)^{2}+y^{2}=1\)。 \((2)\)设\(PQ\)中点\(N(x,y)\)。\(\angle PBQ=90^{\circ} \Rightarrow  Rt\triangle PBQ\)斜边中线 \(|BN|=|PN|=|PQ|/2\)；又\(ON\perp PQ\)（垂径定理），\(|OP|^{2}=|ON|^{2}+|PN|^{2} \Rightarrow  |PN|^{2}=4-(x^{2}+y^{2})\)。于是 \(|BN|^{2}=|PN|^{2}\)：\((x-1)^{2}+(y-1)^{2}=4-x^{2}-y^{2}\)，展开整理 \(2x^{2}+2y^{2}-2x-2y-2=0\)，即 \((x-1/2)^{2}+(y-1/2)^{2}=3/2\)。}
\end{ansblock}

\zux{距离差最值·定义转化（隐圆）}
% 席注（账面）：⚠照录后节知识（双曲线定义2.6.1；§5已注"最晚知识2.6.1——建议课时14后回看使用，观察项"），本区不改写。
\tuhao{10} \tieside{难}已知点\(M(1,2)\)，点\(P\)是双曲线\(C\)：\(x^{2}/4-y^{2}/12=1\)左支上的动点，\(N\)是圆\(D\)：\((x+4)^{2}+y^{2}=1\)上的动点，则\(|PM|-|PN|\)的最小值为（　）
\bindopt {\fontsize{10.5pt}{\qpTJgLeadLiti}\selectfont \makebox[\dimexpr0.5\linewidth-0.5\lxhang-0.25em\relax][l]{\(A\)．\(5-\sqrt{10}\)}\makebox[\dimexpr0.5\linewidth-0.5\lxhang-0.25em\relax][l]{\(B\)．\(\sqrt{10}-5\)}\par}
\bindopt {\fontsize{10.5pt}{\qpTJgLeadLiti}\selectfont \makebox[\dimexpr0.5\linewidth-0.5\lxhang-0.25em\relax][l]{\(C\)．\(\sqrt{13}-3\)}\makebox[\dimexpr0.5\linewidth-0.5\lxhang-0.25em\relax][l]{\(D\)．\(3-\sqrt{13}\)}\par}
\begin{ansblock}[2章-拓-课时10-10]
% ans:2章-拓-课时10-10
\ansitem{10}{\(D\)}
\ansnote{详解}{双曲线 \(c=\sqrt{4+12}=4\)，圆心\(D(-4,0)\)即左焦点，右焦点\(F_{2}(4,0)\)；圆半径\(r=1\)。\(|PN|\leqslant |PD|+r\)（\(N\)取\(PD\)延长线与圆交点时等号），故 \(-|PN|\geqslant -|PD|-1\)；又 \(|PM|\geqslant |PF_{2}|-|MF_{2}|=|PF_{2}|-\sqrt{13}\)（\(P\)、\(M\)、\(F_{2}\)共线时等号）。两式相加：\(|PM|-|PN|\geqslant |PF_{2}|-|PD|-\sqrt{13}-1\)。\(P\)在左支：\(|PF_{2}|-|PD|=2a=4\)（双曲线定义），得 \(|PM|-|PN|\geqslant 4-\sqrt{13}-1=3-\sqrt{13}\)。选\(D\)。}
\end{ansblock}

\zux{动弦交点隐圆＋向量最值（多步）}
% 席注（账面）：⚠存疑：源详解"P轨迹圆半径√2"与"|PD|min=|CM|−r₁−r₂"中 r₁ 混用（弦中点圆半径为√3、圆C半径2），本区按|CM|=3√2、√3+√2复原链条，结果4√2−2√3与源选B一致；P=直径端点退化（m致l₁∥l₂无，恒垂直故无退化）✓。
\tuhao{11} \tieside{中}已知直线\(l_{1}\)：\(mx-y-3m+1=0\)与\(l_{2}\)：\(x+my-3m-1=0\)相交于点\(P\)，线段\(AB\)是圆\(C\)：\((x+1)^{2}+(y+1)^{2}=4\)的一条动弦，且\(|AB|=2\)，则\(|\)向量\(PA+\)向量\(PB|\)的最小值是（　）
\bindopt {\fontsize{10.5pt}{\qpTJgLeadLiti}\selectfont \makebox[\dimexpr0.5\linewidth-0.5\lxhang-0.25em\relax][l]{\(A\)．\(2\sqrt{2}-\sqrt{3}\)}\makebox[\dimexpr0.5\linewidth-0.5\lxhang-0.25em\relax][l]{\(B\)．\(4\sqrt{2}-2\sqrt{3}\)}\par}
\bindopt {\fontsize{10.5pt}{\qpTJgLeadLiti}\selectfont \makebox[\dimexpr0.5\linewidth-0.5\lxhang-0.25em\relax][l]{\(C\)．\(2\sqrt{2}-1\)}\makebox[\dimexpr0.5\linewidth-0.5\lxhang-0.25em\relax][l]{\(D\)．\(4\sqrt{2}-2\)}\par}
\begin{ansblock}[2章-拓-课时10-11]
% ans:2章-拓-课时10-11
\ansitem{11}{\(B\)}
\ansnote{详解}{\(l_{1}\perp l_{2}\)（\(m\cdot 1+(-1)\cdot m=0\)），\(l_{1}\)恒过定点\(S(3,1)\)、\(l_{2}\)恒过定点\(T(1,3)\)（\(m=0\)与一般值均合），\(\angle SPT=90^{\circ} \Rightarrow  P\)在以\(ST\)为直径的圆上：圆心\(M(2,2)\)、半径 \(r_{2}=|ST|/2=\sqrt{2}\)，即 \((x-2)^{2}+(y-2)^{2}=2\)。设\(AB\)中点\(D\)：\(|CD|=\sqrt{2^{2}-1^{2}}=\sqrt{3} \Rightarrow  D\)在以\(C(-1,-1)\)为心、\(\sqrt{3}\)为半径的圆上。\(PA+PB=2PD\)，\(|PD|\geqslant |MC|-r_{2}-\sqrt{3}=\sqrt{18}-\sqrt{2}-\sqrt{3}=3\sqrt{2}-\sqrt{2}-\sqrt{3}=2\sqrt{2}-\sqrt{3}\)（\(P\)、\(D\)在连心线上时等号可达），故 \(|PA+PB|_{2=}2|PD|\geqslant 4\sqrt{2}-2\sqrt{3}\)。选\(B\)。}
\end{ansblock}

\zux{垂直双弦中点轨迹＋距离最值（隐圆·向量）}
% 席注（账面）：源详解含图（图注从源）；"2NQ²+2QM²=4"步源dump作"NQ²+QM²=2"等价，照录并注。
\tuhao{12} \tieside{难}平面直角坐标系\(xOy\)中，\(M\)：\((x-1)^{2}+(y-1)^{2}=4\)，过点\(P(2,2)\)作两条直线，被圆\(M\)截得弦\(AB\)、\(CD\)，满足\(AB\perp CD\)．设线段\(AC\)的中点为\(N\)，则\(|ON|\)的最小值为\kongda{\((3\sqrt{2}-\sqrt{6})/2\)}．
\begin{ansblock}[2章-拓-课时10-12]
% ans:2章-拓-课时10-12
\ansitem{12}{\((3\sqrt{2}-\sqrt{6})/2\)}
\ansnote{详解}{\(M(1,1)\)，\(P(2,2)\)，\(|MP|=\sqrt{2}\)。设\(MP\)中点\(Q(3/2,3/2)\)。\(AB\perp CD\) 于 \(P \Rightarrow  Rt\triangle APC\)中 \(N\)为斜边\(AC\)中点：\(|PN|=|AN|=|CN|\)。垂径：\(AN^{2}+MN^{2}=AM^{2}=4\)，故 \(PN^{2}+MN^{2}=4\)。向量：\(PN=NQ+QP\)、\(MN=NQ+QM\)，\(QP=-QM\) 且 \(|QM|=|QP|=\sqrt{2}/2\)，相加：\(PN^{2}+MN^{2}=2NQ^{2}+2QM^{2}=2NQ^{2}+1=4 \Rightarrow  NQ^{2}=3/2\)，\(N\)在以\(Q\)为心、\(\sqrt{3/2}=\sqrt{6}/2\)为半径的圆上。\(|ON|\geqslant |OQ|-\sqrt{6}/2=(3\sqrt{2})/2-\sqrt{6}/2=(3\sqrt{2}-\sqrt{6})/2\)。验算存在性：\(N=OQ\)连线与圆交点确对应某对垂直弦（\(|PN|^{2}+|MN|^{2}=4\) 与 \(PN=MN\) 型中点可反解），等号可达。答 \((3\sqrt{2}-\sqrt{6})/2\)。}
\end{ansblock}

\zux{定比轨迹·含参分类（直接法＋曲线形状判定）}
% 席注（账面）：与简5（习题2-4 B①）同族第2件（定值√2具体vs参数k分类，判别成立，§7已注）。
\tuhao{13} \tieside{中}已知某曲线上的点到定点\(O(0,0)\)与到定点\(A(a,0)\)（\(a\neq 0\)）的距离的比值为\(k\)，求此曲线的方程，并判定曲线的形状．
\liubai[16mm]{解答书写区}
\begin{ansblock}[2章-拓-课时10-13]
% ans:2章-拓-课时10-13
\ansitem{13}{\(k=1\)：\(x=a/2\)（线段\(OA\)的垂直平分线）；\(k\neq 1\)（\(k>0\)）：\((k^{2}-1)x^{2}+(k^{2}-1)y^{2}-2k^{2}ax+k^{2}a^{2}=0\)，即圆心\((k^{2}a/(k^{2}-1),0)\)、半径\(|ka/(k^{2}-1)|\)的圆（阿波罗尼斯圆）}
\ansnote{详解}{设\(M(x,y)\)为曲线上任意一点，依题意 \(|MO|/|MA|=k\)（\(k>0\)）：\(\sqrt{x^{2}+y^{2}}/\sqrt{}((x-a)^{2}+y^{2})=k\)，两边平方并展开：\(x^{2}+y^{2}=k^{2}[(x-a)^{2}+y^{2}]\)，整理得 \((k^{2}-1)x^{2}+(k^{2}-1)y^{2}-2k^{2}ax+k^{2}a^{2}=0\)（化简步步同解，两性保住了）。分类： ①\(k=1\)：方程化为 \(-2ax+a^{2}=0\)，即 \(x=a/2\)——过\(OA\)中点\((a/2,0)\)且垂直于\(x\)轴的直线，即线段\(OA\)的垂直平分线。 ②\(k\neq 1\)：两边除以 \(k^{2}-1\)：\(x^{2}+y^{2}-[2k^{2}a/(k^{2}-1)]x+k^{2}a^{2}/(k^{2}-1)=0\)。配方：\((x-k^{2}a/(k^{2}-1))^{2}+y^{2}=k^{4}a^{2}/(k^{2}-1)^{2}-k^{2}a^{2}/(k^{2}-1)=k^{2}a^{2}[k^{2}-(k^{2}-1)]/(k^{2}-1)^{2}=k^{2}a^{2}/(k^{2}-1)^{2}>0\) 恒成立，故恒表示圆，圆心\((k^{2}a/(k^{2}-1),0)\)、半径\(|ka/(k^{2}-1)|\)（无需再验 \(D^{2}+E^{2}-4F>0\)，配方已直接给出正的半径平方）。 验算（\(k=1/2\)、\(a=2\) 特值）：圆心\((-2/3,0)\)、半径\(4/3\)，圆与\(x\)轴交点 \((2/3,0)\)：\(|MO|/|MA|=(2/3)/(4/3)=1/2\)；\((-2,0)\)：\(2/4=1/2\)。}
\end{ansblock}

\zux{圆系方程＋完备性辨析（退化陷阱位）}
% 席注（账面）：本题即§3"完备性陷阱位"——设C₁+λC₂漏掉λ"取无穷"端（C₂自身），是"方程与曲线对应不完备"的活例；验算交点：联立C₁、C₂得公共弦y=−x，交点(√2/2,−√2/2)、(−√2/2,√2/2)，均在C₂上✓。
\tuhao{14} \tieside{中}求通过圆\(x^{2}+y^{2}=1\)与\(x^{2}+y^{2}-4x-4y-1=0\)的交点，并且过点\((2,-1)\)的圆的方程．
\liubai[16mm]{解答书写区}
\begin{ansblock}[2章-拓-课时10-14]
% ans:2章-拓-课时10-14
\ansitem{14}{\(x^{2}+y^{2}-4x-4y-1=0\)（即圆\(C_{2}\)本身，\((x-2)^{2}+(y-2)^{2}=9\)）}
\ansnote{详解}{记 \(C_{1}\)：\(x^{2}+y^{2}-1=0\)，\(C_{2}\)：\(x^{2}+y^{2}-4x-4y-1=0\)。①两圆确有两个交点：\(C_{1}\)圆心\((0,0)\)半径\(1\)，\(C_{2}\)圆心\((2,2)\)半径\(3\)，圆心距\(2\sqrt{2}\)，满足 \(3-1<2\sqrt{2}<3+1\) 。②过两交点的一切圆可表为 \(C_{1}+\lambda C_{2}=0\)（\(\lambda \neq -1\)；\(\lambda =-1\) 时退化为公共弦所在直线 \(x+y=0\)）——注意该系不含圆\(C_{2}\)本身。③代点\((2,-1)\)：\(C_{1}\)值\(=4+1-1=4\)，\(C_{2}\)值\(=4+1-8+4-1=0\)，系内圆过\((2,-)\)须 \(4+\lambda \cdot 0=0\)，无解！④而点\((2,-1)\)恰在\(C_{2}\)上（上式\(=0\)），\(C_{2}\)本身即\("\)过两交点且过\((2,-1)"\)的圆；两相交点加线外一点定一圆，故唯一，所求圆即 \(C_{2}\)：\(x^{2}+y^{2}-4x-4y-1=0\)。}
\end{ansblock}

\zux{轨迹方程·相关点法（重心）}
% 席注（账面）：⚠照录含后节载体（双曲线方程与焦点，2.6.1知识）——成品自带"载体先例照录"注记（§3），本区不改；与G1、拓展-13同族三件三取二口径照§5。
\tuhao{15} \tieside{中}双曲线 \(x^{2}/9-y^{2}=1\) 有动点\(P\)，\(F_{1}\)、\(F_{2}\)是曲线的两个焦点，求\(\triangle PF_{1}F_{2}\)的重心\(M\)的轨迹方程．
\liubai[16mm]{解答书写区}
\begin{ansblock}[2章-拓-课时10-15]
% ans:2章-拓-课时10-15
\ansitem{15}{\(x^{2}-9y^{2}=1\)（\(y\neq 0\)）}
\ansnote{详解}{设\(P(x_{1},y_{1})\)、\(M(x,y)\)。双曲线中 \(a=3\)、\(b=1\)、\(c=\sqrt{10}\)，\(F_{1}(-\sqrt{10},0)\)、\(F_{2}(\sqrt{10},0)\)。\(P\)为三角形顶点，\(y_{1}\neq 0\)。重心坐标公式：\(x=(x_{1}-\sqrt{10}+\sqrt{10})/3=x_{1}/3\)，\(y=y_{1}/3\)，即 \(x_{1}=3x\)、\(y_{1}=3y\)。\(P\)在双曲线上：\((3x)^{2}/9-(3y)^{2}=1\)，即 \(x^{2}-9y^{2}=1\)；由 \(y_{1}\neq 0\) 得 \(y\neq 0\)。故重心\(M\)的轨迹方程为 \(x^{2}-9y^{2}=1\)（\(y\neq 0\)）。}
\end{ansblock}

\jietitle{2.5.1　椭圆的标准方程（课时11·拓区）}
% 【11-拓2｜撤位登记（案2令2：与11-简7 逐字同题撤席让位，唯一活位＝11-简7）】本席不设题——跳号留痕（拓3 起），注记随席（台账抄件）；防双收 B21。
\zux{利用定义求方程（根式和式）}
% 席注（账面）：〔0913：原「与拓2同题型拓展两件注」随拓2撤位销注——案2裁撤席让位后，本席为定义求方程组合正文+G5外的唯一拓展件（§5注随动）。〕
\tuhao{01} \tieside{中}已知椭圆\(C\)上任意一点\(P(x,y)\)都满足关系式 \(\sqrt{}((x-1)^{2}+y^{2})+\sqrt{}((x+1)^{2}+y^{2})=4\)，则椭圆\(C\)的标准方程为（　）
\bindopt {\fontsize{10.5pt}{\qpTJgLeadLiti}\selectfont \makebox[\dimexpr0.5\linewidth-0.5\lxhang-0.25em\relax][l]{\(A\)．\(x^{2}/3+y^{2}/4=1\)}\makebox[\dimexpr0.5\linewidth-0.5\lxhang-0.25em\relax][l]{\(B\)．\(x^{2}/4+y^{2}/3=1\)}\par}
\bindopt {\fontsize{10.5pt}{\qpTJgLeadLiti}\selectfont \makebox[\dimexpr0.5\linewidth-0.5\lxhang-0.25em\relax][l]{\(C\)．\(x^{2}/16+y^{2}/15=1\)}\makebox[\dimexpr0.5\linewidth-0.5\lxhang-0.25em\relax][l]{\(D\)．\(x^{2}/4+y^{2}=1\)}\par}
\begin{ansblock}[2章-拓-课时11-01]
% ans:2章-拓-课时11-01
\ansitem{01}{\(B\)}
\ansnote{详解}{关系式即\("\)到两定点\((1,0)\)、\((-1,0)\)距离和为常数\(4"\)，\(4>|F_{1}F_{2}|=2\)，\(P\)轨迹为椭圆：焦点在\(x\)轴、\(c=1\)、\(2a=4\Rightarrow a=2\)、\(b^{2}=a^{2}-c^{2}=3\)。标准方程 \(x^{2}/4+y^{2}/3=1\)。选\(B\)。}
\end{ansblock}

\zux{动圆相切条件下的椭圆轨迹（挖点）}
% 席注（账面）：圆切轨迹族（跨课时↔课时13件8-#5）§7注记照。
\tuhao{03} \tieside{中}一动圆\(M_{1}\)与圆\(M\)：\((x+1)^{2}+y^{2}=1\)外切，与圆\(M_{2}\)：\((x-1)^{2}+y^{2}=9\)内切，则动圆圆心\(M_{1}\)的轨迹方程为（　）
\lxopt{\(A\)．\(x^{2}/4+y^{2}/3=1\)}
\lxopt{\(B\)．\(x^{2}/4+y^{2}/3=1\)（\(x\neq \pm 2\)）}
\lxopt{\(C\)．\(x^{2}/16+y^{2}/15=1\)}
\lxopt{\(D\)．\(x^{2}/4+y^{2}/3=1\)（\(x\neq -2\)）}
\begin{ansblock}[2章-拓-课时11-03]
% ans:2章-拓-课时11-03
\ansitem{03}{\(D\)}
\ansnote{详解}{圆\(M\)圆心\(M(-1,0)\)半径\(1\)；圆\(M_{2}\)圆心\((1,0)\)半径\(3\)。\(|MM_{2}|=2=3-1\)，两定圆内切于点\(A(-2,0)\)（左端点）。设动圆半径\(r_{1}\)：外切 \(\Rightarrow  |M_{1}M|=r_{1}+1\)；内切（动圆含于大圆内） \(\Rightarrow  |M_{1}M_{2}|=3-r_{1}\)。相加：\(|M_{1}M|+|M_{1}M_{2}|=4>|MM_{2}|=2\)，\(M_{1}\)轨迹以\(M\)、\(M_{2}\)为焦点的椭圆：\(a=2\)、\(c=1\)、\(b^{2}=3\)，\(x^{2}/4+y^{2}/3=1\)。退化剔除：\(r_{1}=0\) 时 \(M_{1}\) 与切点\(A(-2,0)\)重合（\("\)动圆\("\)退化为点，圆\(M_{1}\)不存在），故去 \(x=-2\)（右端点\(x=2\)处 \(r_{1}=\)… 检验：\(M_{1}(2,0)\)：\(|M_{1}M|=3=r_{1}+1\rightarrow r_{1}=2\)；\(|M_{1}M_{2}|=1=3-2\)  合法保留）。选\(D\)。}
\end{ansblock}

\zux{到焦点距离计算（两点距离＋定义）}
% 席注（账面）：多问计1；焦半径公式未用（2.5.2内容），全程定义式。
\tuhao{04} \tieside{中}已知椭圆\(x^{2}/25+y^{2}/9=1\)上一点\(P\)的横坐标是\(2\)，求：\((1)\)点\(P\)到椭圆左焦点的距离\(|PF_{1}|\)；\((2)\)点\(P\)到椭圆右焦点的距离\(|PF_{2}|\)．
\liubai[16mm]{解答书写区}
\begin{ansblock}[2章-拓-课时11-04]
% ans:2章-拓-课时11-04
\ansitem{04}{\((1)33/5\)；\((2)17/5\)}
\ansnote{详解}{\(a=5\)、\(b=3\)、\(c=4\)，\(F_{1}(-4,0)\)、\(F_{2}(4,0)\)。\(P\)横坐标\(2\)：\(4/25+y^{2}/9=1 \Rightarrow  y^{2}=9\cdot 21/25=189/25\)。\((1)|PF_{1}|=\sqrt{}((2+4)^{2}+y^{2})=\sqrt{36+189/25}=\sqrt{1089/25}=33/5\)。\((2)\)定义：\(|PF_{2}|=2a-|PF_{1}|=10-33/5=17/5\)。（直算复核：\(|PF_{2}|=\sqrt{}((2-4)^{2}+189/25)=\sqrt{4+189/25}=\sqrt{289/25}=17/5\)  两路一致）}
\end{ansblock}

\zux{椭圆上点到定点距离最值（二次函数）}
% 席注（账面）："到焦点距离两问"标签照§5枚举为-5、-6两题（本席=-6｜PA最值）。
\tuhao{05} \tieside{中}已知\(P\)是椭圆\(x^{2}/4+y^{2}/36=1\)上一点，\(A(0,5)\)，求\(|PA|\)的最小值与最大值．
\liubai[16mm]{解答书写区}
\begin{ansblock}[2章-拓-课时11-05]
% ans:2章-拓-课时11-05
\ansitem{05}{最小值\(\sqrt{14}/4\)，最大值\(11\)}
\ansnote{详解}{焦点在\(y\)轴，\(a=6\)、\(b=2\)。设\(P(x_{0},y_{0})\)：\(x_{0}^{2}/4+y_{0}^{2}/36=1 \Rightarrow  x_{0}^{2}=4(1-y_{0}^{2}/36)=4-y_{0}^{2}/9\)，\(y_{0}\in [-6,6]\)。\(|PA|^{2}=x_{0}^{2}+(y_{0}-5)^{2}=4-y_{0}^{2}/9+y_{0}^{2}-10y_{0}+25=(8/9)y_{0}^{2}-10y_{0}+29=(8/9)(y_{0}-45/8)^{2}+7/8\)。顶点\(y_{0}=45/8\in [-6,6]\)，最小值\(|PA|^{2}=7/8 \Rightarrow  |PA|min=\sqrt{7/8}=\sqrt{14}/4\)；端点比较：\(y_{0}=-6\) 时 \((8/9)\cdot 36+60+29=121 \Rightarrow  |PA|=11\)；\(y_{0}=6\) 时 \(32-60+29=1 \Rightarrow  |PA|=1\)。最大\(11\)、最小\(\sqrt{14}/4\)。}
\end{ansblock}

\zux{椭圆动点与圆动点距离最值（隐圆转化）}
% 席注（账面）：与"圆上动点|MN|最大"§5枚举-7对位本席（-6/-7两题次序照§5行序）。
\tuhao{06} \tieside{中}已知点\(M\)在椭圆\(x^{2}/18+y^{2}/9=1\)上运动，点\(N\)在圆\(x^{2}+(y-1)^{2}=1\)上运动，则\(|MN|\)的最大值为（　）
\bindopt {\fontsize{10.5pt}{\qpTJgLeadLiti}\selectfont \makebox[\dimexpr0.5\linewidth-0.5\lxhang-0.25em\relax][l]{\(A\)．\(1+\sqrt{19}\)}\makebox[\dimexpr0.5\linewidth-0.5\lxhang-0.25em\relax][l]{\(B\)．\(1+2\sqrt{5}\)}\par}
\bindopt {\fontsize{10.5pt}{\qpTJgLeadLiti}\selectfont \makebox[\dimexpr0.5\linewidth-0.5\lxhang-0.25em\relax][l]{\(C\)．\(5\)}\makebox[\dimexpr0.5\linewidth-0.5\lxhang-0.25em\relax][l]{\(D\)．\(6\)}\par}
\begin{ansblock}[2章-拓-课时11-06]
% ans:2章-拓-课时11-06
\ansitem{06}{\(B\)}
\ansnote{详解}{圆心\(C(0,1)\)、半径\(1\)：\(|MN|\leqslant |MC|+1\)（\(N\)在\(MC\)延长线与圆交点时等号）。设\(M(x_{0},y_{0})\)，\(x_{0}^{2}=18-2y_{0}^{2}\)（\(y_{0}\in [-3,3]\)）。\(|MC|^{2}=x_{0}^{2}+(y_{0}-1)^{2}=18-2y_{0}^{2}+y_{0}^{2}-2y_{0}+1=-(y_{0}+1)^{2}+20\leqslant 20\)，\(y_{0}=-1\in [-3,3]\) 可取，\(|MC|max=2\sqrt{5}\)。\(|MN|max=2\sqrt{5}+1\)。选\(B\)。}
\end{ansblock}

\zux{定义与焦距离最值·多选辨析}
% 席注（账面）："到焦点距离最值"界内推导本区补一行口径（端点直算型）。
\tuhao{07}\duoxuan\hspace{0.6em} \tieside{中}设\(P\)是椭圆\(x^{2}/5+y^{2}/3=1\)上的动点，则（　）
\lxopt{\(A\)．点\(P\)到该椭圆的两个焦点的距离之和为\(2\sqrt{5}\)}
\lxopt{\(B\)．点\(P\)到该椭圆的两个焦点的距离之和为\(2\sqrt{2}\)}
\lxopt{\(C\)．点\(P\)到左焦点距离的最大值为\(\sqrt{5}+\sqrt{2}\)}
\lxopt{\(D\)．点\(P\)到左焦点距离的最大值为\(\sqrt{5}+2\sqrt{2}\)}
\begin{ansblock}[2章-拓-课时11-07]
% ans:2章-拓-课时11-07
\ansitem{07}{\(AC\)}
\ansnote{详解}{\(a^{2}=5\)、\(b^{2}=3\)、\(c^{2}=2\)：\(a=\sqrt{5}\)、\(c=\sqrt{2}\)。\(A\)对\(B\)错：定义距离和\(=2a=2\sqrt{5}\)。\(C\)对\(D\)错：\(|PF_{1}|=2a-|PF_{2}|\)，要最大须\(|PF_{2}|\)最小\(=a-c\)（右顶点），故\(max=2a-(a-c)=a+c=\sqrt{5}+\sqrt{2}\)。（\(|PF_{2}|\in [a-c,a+c]\)由拓\(4\)型距离直算可得，未引\(2.5.2\)性质表。）}
\end{ansblock}

\zux{到焦点和定点距离和最值（左焦点转化定和）}
% 席注（账面）：源详解含图（图注从源）；源等号行"AF₂的延长线"表述方向易误（应为过F₂一侧），本区照正确构式校订；A内点判定为本区补行。
\tuhao{08} \tieside{中}已知\(F_{1}\)是椭圆\(9x^{2}+25y^{2}=225\)的左焦点，\(P\)是此椭圆上的动点，\(A(1,2)\)是一定点，则\(|PA|+|PF_{1}|\)的最大值为\kongda{\(10+\sqrt{13}\)}．
\begin{ansblock}[2章-拓-课时11-08]
% ans:2章-拓-课时11-08
\ansitem{08}{\(10+\sqrt{13}\)}
\ansnote{详解}{标准形\(x^{2}/25+y^{2}/9=1\)：\(a=5\)、\(c=4\)，\(F_{1}(-4,0)\)、右焦点\(F_{2}(4,0)\)。定义转化：\(|PF_{1}|=10-|PF_{2}|\)，\(|PA|+|PF_{1}|=10+(|PA|-|PF_{2}|)\leqslant 10+|AF_{2}|=10+\sqrt{}((1-4)^{2}+2^{2})=10+\sqrt{13}\)（三角不等式，\(|PA|-|PF_{2}|\leqslant |AF_{2}|\) 当\(F_{2}\)在线段\(PA\)上取等）。等号可取性：\(A(1,2)\)在椭圆内（\(1/25+4/9=109/225<1\)），自\(A\)出发过\(F_{2}\)的射线必与椭圆交于一点\(P\)（\(F_{2}\)亦椭圆内点，交点在\(F_{2}\)外侧），此时 \(|PA|-|PF_{2}|=|AF_{2}|=\sqrt{13}\)。最大值\(10+\sqrt{13}\)。}
\end{ansblock}

\zux{垂直条件求P坐标（隐圆联立）}
% 席注（账面）：多解全列照源。
\tuhao{09} \tieside{中}已知椭圆\(x^{2}/45+y^{2}/20=1\)上一点\(P\)与两个焦点的连线互相垂直，求点\(P\)的坐标．
\liubai[16mm]{解答书写区}
\begin{ansblock}[2章-拓-课时11-09]
% ans:2章-拓-课时11-09
\ansitem{09}{\((3,4)\)或\((3,-4)\)或\((-3,4)\)或\((-3,-4)\)}
\ansnote{详解}{\(c^{2}=45-20=25\)：\(F_{1}(-5,0)\)、\(F_{2}(5,0)\)。设\(P(x,y)\)：垂直 \(\Rightarrow \) 向量\(PF_{1}\cdot \)向量\(PF_{2}=(-5-x,-y)\cdot (5-x,-y)=(x^{2}-25)+y^{2}=0\)，即\(P\)在圆\(x^{2}+y^{2}=25\)上；联立 \(x^{2}/45+y^{2}/20=1\)：代\(x^{2}=25-y^{2}\)：\((25-y^{2})/45+y^{2}/20=1\)，\(\times 180\)：\(100-4y^{2}+9y^{2}=180 \Rightarrow  y^{2}=16\)、\(x^{2}=9\)。\(P=(\pm 3,\pm 4)\)四解。验算任一组：\(9/45+16/20=1/5+4/5=1\)；\((-5-3)(5-3)+(-4)(-4)?\) 取\((3,4)\)：\((x^{2}-25)+y^{2}=9-25+16=0\)。}
\end{ansblock}

\zux{依条件求方程（共焦＋轴不定分类·多问计1）}
% 席注（账面）：源详解双法（分类法一/设系法二）照录取要；与简1同题型正文1＋拓展1（§5注）。
\tuhao{10} \tieside{中}求满足下列条件的椭圆的标准方程．
\((1)\)过点\(Q(2\sqrt{2},1)\)，且与椭圆\(x^{2}/9+y^{2}/4=1\)有公共的焦点；
\((2)\)中心在原点，焦点在坐标轴上，且经过两点\(P(1/3,1/3)\)、\(Q(0,-1/2)\)．
\liubai[24mm]{解答书写区}
\begin{ansblock}[2章-拓-课时11-10]
% ans:2章-拓-课时11-10
\ansitem{10}{\((1)x^{2}/10+y^{2}/5=1\)；\((2)y^{2}/(1/4)+x^{2}/(1/5)=1\)（即\(5x^{2}+4y^{2}=1\)）}
\ansnote{详解}{\((1)\)共焦：\(c^{2}=9-4=5\)，设 \(x^{2}/a^{2}+y^{2}/(a^{2}-5)=1\)，代\(Q\)：\(8/a^{2}+1/(a^{2}-5)=1\)，令\(u=a^{2}\)：\(8(u-5)+u=u(u-5) \Rightarrow  u^{2}-14u+40=0 \Rightarrow  u=10\)或\(4\)（\(u=4<5\)舍，\(b^{2}<0\)）：\(a^{2}=10\)、\(b^{2}=5\)，\(x^{2}/10+y^{2}/5=1\)。验算\(Q\)：\(8/10+1/5=1\) 。（源法二共焦系\(x^{2}/(9+\lambda )+y^{2}/(4+\lambda )=1\)代点\(\lambda =1\)同果。） \((2)\)焦点轴不定，设 \(mx^{2}+ny^{2}=1\)（\(m>0\)、\(n>0\)、\(m\neq n\)）。代\(P(1/3,1/3)\)：\((m+n)/9=1\)；代\(Q(0,-1/2)\)：\(n/4=1 \Rightarrow  n=4\)、\(m=5\)：\(5x^{2}+4y^{2}=1\)，标准形 \(y^{2}/(1/4)+x^{2}/(1/5)=1\)（焦点在\(y\)轴）。验算：\(P\)：\(5/9+4/9=1\)。}
\end{ansblock}

\zux{祖暅原理·椭球体积（文化）}
% 席注（账面）：与拓12同题型拓展两件注（§5）；"长轴为4"取a=2（该题以半轴记，口径照源）。　图债注：「如图」指源件配图，印面无图（冻片【注】裁定：去装饰图/条件全在文字/尺寸随注——图债登记汇编报告）
\tuhao{11} \tieside{中}祖{\fangsong 暅}（公元\(5-6\)世纪，祖冲之之子，是我国齐梁时代的数学家）提出了一条原理：\("\)幂势既同，则积不容异．\("\)意思是：两个等高的几何体若在所有等高处的水平截面的面积相等，则这两个几何体的体积相等．该原理在西方直到十七世纪才由意大利数学家卡瓦列利发现，比祖{\fangsong 暅}晚一千一百多年．椭球体是椭圆绕其轴旋转所成的旋转体．如图，将底面直径皆为\(2b\)、高皆为\(a\)的椭半球体及已被挖去了圆锥体的圆柱体放置于同一平面\(\beta \)上，用平行于平面\(\beta \)的平面于距平面\(\beta \)任意高\(d\)处截，得到\(S\_\)圆及\(S\_\)环两截面，可以证明\(S\_\)圆\(=S\_\)环总成立．据此，短轴长为\(2cm\)、长轴为\(4cm\)的椭球体的体积是（　）\(cm^{3}\)．〔图注：半球体与挖锥圆柱并置示意，照录自源〕
\bindopt {\fontsize{10.5pt}{\qpTJgLeadLiti}\selectfont \makebox[\dimexpr0.5\linewidth-0.5\lxhang-0.25em\relax][l]{\(A\)．\(2\pi /3\)}\makebox[\dimexpr0.5\linewidth-0.5\lxhang-0.25em\relax][l]{\(B\)．\(4\pi /3\)}\par}
\bindopt {\fontsize{10.5pt}{\qpTJgLeadLiti}\selectfont \makebox[\dimexpr0.5\linewidth-0.5\lxhang-0.25em\relax][l]{\(C\)．\(8\pi /3\)}\makebox[\dimexpr0.5\linewidth-0.5\lxhang-0.25em\relax][l]{\(D\)．\(16\pi /3\)}\par}
\begin{ansblock}[2章-拓-课时11-11]
% ans:2章-拓-课时11-11
\ansitem{11}{\(C\)}
\ansnote{详解}{由\(S\_\)圆\(=S\_\)环恒成立，半椭球体积\(=\)柱\(-\)锥\(=\pi b^{2}a-(1/3)\pi b^{2}a=(2/3)\pi b^{2}a\)。短轴长\(2 \Rightarrow  b=1\)；长轴\(4 \Rightarrow \) 半长轴\(a=2\)。\(V\)椭球\(=2\cdot (2/3)\pi \cdot 1^{2}\cdot 2=8\pi /3\)。选\(C\)。}
\end{ansblock}

\zux{祖暅原理·旋转体体积（类比）}
% 席注（账面）：源详解首步"截面与顶点距离h(0≤h≤4)"界误（旋转半轴为3），已照理订正为0≤h≤3，答案16π不变；§5注同题型两件（-24/-25）照。　图债注：「如图」指源件配图，印面无图（冻片【注】裁定：去装饰图/条件全在文字/尺寸随注——图债登记汇编报告）
\tuhao{12} \tieside{中}运用祖{\fangsong 暅}原理计算球的体积时，夹在两个平行平面之间的两个几何体，被平行于这两个平面的任意一个平面所截，若截面面积都相等，则这两个几何体的体积相等．构造一个底面半径和高都与球的半径相等的圆柱，与半球（如图①）放置在同一平面上，然后在圆柱内挖去一个以圆柱下底面圆心为顶点、圆柱上底面为底面的圆锥后得到一新几何体（如图②），用任何一个平行于底面的平面去截它们时，可证得所截得的两个截面面积相等，由此可证明新几何体与半球体积相等．现将椭圆\(x^{2}/4+y^{2}/9=1\)绕\(y\)轴旋转一周后得一橄榄状的几何体（如图③），类比上述方法，运用祖{\fangsong 暅}原理可求得其体积等于（　）．〔图注：图①②③三联示意，照录自源〕
\bindopt {\fontsize{10.5pt}{\qpTJgLeadLiti}\selectfont \makebox[\dimexpr0.25\linewidth-0.25\lxhang-0.25em\relax][l]{\(A\)．\(8\pi \)}\makebox[\dimexpr0.25\linewidth-0.25\lxhang-0.25em\relax][l]{\(B\)．\(16\pi \)}\makebox[\dimexpr0.25\linewidth-0.25\lxhang-0.25em\relax][l]{\(C\)．\(24\pi \)}\makebox[\dimexpr0.25\linewidth-0.25\lxhang-0.25em\relax][l]{\(D\)．\(32\pi \)}\par}
\begin{ansblock}[2章-拓-课时11-12]
% ans:2章-拓-课时11-12
\ansitem{12}{\(B\)}
\ansnote{详解}{绕\(y\)轴旋转：截面为半径\(|x|\)的圆盘。构造底面半径\(2\)、高\(3\)的圆柱，挖去以下底圆心为顶点、以上底为底的圆锥。高\(h\)处（\(0\leqslant h\leqslant 3\)）：圆锥截面半径\(r=(2/3)h\)，环面积\(=4\pi -(4h^{2}\pi )/9\)；橄榄体同高处截面\(=\pi x^{2}\)，由 \(x^{2}/4+h^{2}/9=1\) 得 \(x^{2}=4(1-h^{2}/9)\)，截面\(=4\pi -(4h^{2}\pi )/9\)，两处恒等。半橄榄\(=\)柱\(-\)锥\(=\pi \cdot 2^{2}\cdot 3-(1/3)\pi \cdot 2^{2}\cdot 3=12\pi -4\pi =8\pi \)，\(V=2\times 8\pi =16\pi \)。选\(B\)。}
\end{ansblock}

\zux{椭圆定义实际应用（绳-桅杆模型）}
% 席注（账面）：多问不涉，计1题。　图债注：「如图」指源件配图，印面无图（冻片【注】裁定：去装饰图/条件全在文字/尺寸随注——图债登记汇编报告）
\tuhao{13} \tieside{中}船上两根高\(7.5m\)的桅杆相距\(15m\)，一条\(30m\)长的绳子两端系在桅杆的顶上，并按如图所示的方式绷紧．假设绳子位于两根桅杆所在的平面内，求绳子与甲板接触点\(P\)到桅杆\(AB\)的距离（精确到\(0.01m\)）．〔图注：双桅杆\(-\)甲板\(-\)绷绳示意，照录自源〕
\liubai[16mm]{解答书写区}
\begin{ansblock}[2章-拓-课时11-13]
% ans:2章-拓-课时11-13
\ansitem{13}{\(4.75m\)}
\ansnote{详解}{接触点\(P\)处绳两直段和\(=30\)：\(|PA|+|PC|=30\)（\(A\)、\(C\)为两桅杆顶），\(P\)在以\(A\)、\(C\)为焦点的椭圆上：\(2a=30\)、\(a=15\)；两焦点距离\(=\)两桅杆顶间距\(=15\)（两顶等高\(7.5m\)、水平相距\(15m\)），\(2c=15\)、\(c=7.5\)；\(b^{2}=a^{2}-c^{2}=225-56.25=168.75=675/4\)。以\(AC\)中点为原点建系：\(x^{2}/225+y^{2}/(675/4)=1\)。甲板在杆顶下方\(7.5m\)：\(y\_P=-7.5\)，代得 \(x\_P^{2}=225(1-56.25/168.75)=225\cdot (2/3)=150\)，\(P\)偏\(AB\)侧取 \(x\_P=-\sqrt{150}\)。\(P\)到桅杆\(AB\)（所在竖直线\(x=-7.5\)）的距离\(=|x\_P|-7.5=\sqrt{150}-7.5\approx 12.25-7.5=4.75(m)\)。}
\end{ansblock}

\zux{垂直动直线交点轨迹（隐圆）＋焦点距离积范围}
% 席注（账面）：§5枚举行"交轨圆Q＋|QF₁||QF₂|范围"对位；交轨法属2.4知识、圆的参数表示属2.3，界内合规。
\tuhao{14} \tieside{中}已知直线\(l_{1}\)：\(mx-y+m=0\)与直线\(l_{2}\)：\(x+my-1=0\)的交点为\(Q\)，椭圆\(x^{2}/5+y^{2}/2=1\)的焦点为\(F_{1}\)、\(F_{2}\)，则\(|QF_{1}|\cdot |QF_{2}|\)的取值范围是（　）
\bindopt {\fontsize{10.5pt}{\qpTJgLeadLiti}\selectfont \makebox[\dimexpr0.5\linewidth-0.5\lxhang-0.25em\relax][l]{\(A\)．\([2,+\infty )\)}\makebox[\dimexpr0.5\linewidth-0.5\lxhang-0.25em\relax][l]{\(B\)．\([2\sqrt{3},+\infty )\)}\par}
\bindopt {\fontsize{10.5pt}{\qpTJgLeadLiti}\selectfont \makebox[\dimexpr0.5\linewidth-0.5\lxhang-0.25em\relax][l]{\(C\)．\([2,4]\)}\makebox[\dimexpr0.5\linewidth-0.5\lxhang-0.25em\relax][l]{\(D\)．\([2\sqrt{3},4]\)}\par}
\begin{ansblock}[2章-拓-课时11-14]
% ans:2章-拓-课时11-14
\ansitem{14}{\(C\)}
\ansnote{详解}{\(l_{1}\)恒过\((-1,0)\)、\(l_{2}\)恒过\((1,0)\)（\(m=0\)与\(m\neq 0\)均合），且\(l_{1}\perp l_{2}\)（系数\(m\cdot 1+(-1)\cdot m=0\)），交点\(Q\)满足\(\angle (-1,0)Q(1,0)=90^{\circ}\)或重合端点，即\(Q\)在圆\(x^{2}+y^{2}=1\)上（含\((\pm 1,0)\)退化界点——\(m\)致两线共线时无交，逐\(m\)检验交点总存在，轨迹即整圆）。椭圆\(c^{2}=3\)：\(F_{1}(-\sqrt{3},0)\)、\(F_{2}(\sqrt{3},0)\)。设\(Q(cos\theta ,sin\theta )\)：\(|QF_{1}|=\sqrt{}((cos\theta +\sqrt{3})^{2}+sin^{2}\theta )=\sqrt{4+2cos\theta }\)，\(|QF_{2}|=\sqrt{4-2cos\theta }\)，乘积\(=\sqrt{16-12cos^{2}\theta }\in [\sqrt{4},\sqrt{16}]=[2,4]\)（\(cos^{2}\theta \in [0,1]\)）。选\(C\)。}
\end{ansblock}

\zux{焦点三角形面积通式（定义＋余弦定理消元）}
% 席注（账面）：通式结论为课时12常引件，本课时只到"推导"为止（照§3组合注）；α=π/2特例与简6联动由S4把握。
\tuhao{15} \tieside{中}已知点\(P\)为椭圆\(x^{2}/a^{2}+y^{2}/b^{2}=1\)（\(a>b>0\)）上任一点，\(F_{1}\)、\(F_{2}\)为两焦点，\(\angle F_{1}PF_{2}=\alpha \)，求\(\triangle F_{1}PF_{2}\)的面积．
\liubai[16mm]{解答书写区}
\begin{ansblock}[2章-拓-课时11-15]
% ans:2章-拓-课时11-15
\ansitem{15}{\(b^{2}\cdot tan(\alpha /2)\)}
\ansnote{详解}{\(|PF_{1}|+|PF_{2}|=2a\)、\(|F_{1}F_{2}|=2c\)。余弦定理：\(cos\alpha =(|PF_{1}|^{2}+|PF_{2}|^{2}-4c^{2})/(2|PF_{1}||PF_{2}|)=[(2a)^{2}-2|PF_{1}||PF_{2}|-4c^{2}]/(2|PF_{1}||PF_{2}|)=(4b^{2}-2|PF_{1}||PF_{2}|)/(2|PF_{1}||PF_{2}|)=(2b^{2})/(|PF_{1}||PF_{2}|)-1\)，解出 \(|PF_{1}||PF_{2}|=2b^{2}/(1+cos\alpha )\)。\(S=\)½\(|PF_{1}||PF_{2}|sin\alpha =b^{2}\cdot sin\alpha /(1+cos\alpha )=b^{2}\cdot [2sin(\alpha /2)cos(\alpha /2)]/[2cos^{2}(\alpha /2)]=b^{2}\cdot tan(\alpha /2)\)。}
\end{ansblock}

\zux{焦点三角形内心模型（面积法）}
% 席注（账面）：源详解以离心率e=c/a列方程(1+e)²+4e²=4解e=3/5再取倒数——**离心率为2.5.2概念（前引）**，本区校订为纯a、b、c齐次运算路线，同果5/3；"内心"知识属必修平面几何/向量，可讲。
\tuhao{16} \tieside{中}已知椭圆\(C\)：\(x^{2}/a^{2}+y^{2}/b^{2}=1\)（\(a>b>0\)）的左、右焦点分别为\(F_{1}\)、\(F_{2}\)，\(M\)为\(C\)上一点，且\(\triangle MF_{1}F_{2}\)的内心为\(I(x_{0},2)\)，若\(\triangle MF_{1}F_{2}\)的面积为\(4b\)，则\((|MF_{1}|+|MF_{2}|)/|F_{1}F_{2}|=\)（　）
\bindopt {\fontsize{10.5pt}{\qpTJgLeadLiti}\selectfont \makebox[\dimexpr0.5\linewidth-0.5\lxhang-0.25em\relax][l]{\(A\)．\(3/2\)}\makebox[\dimexpr0.5\linewidth-0.5\lxhang-0.25em\relax][l]{\(B\)．\(5/3\)}\par}
\bindopt {\fontsize{10.5pt}{\qpTJgLeadLiti}\selectfont \makebox[\dimexpr0.5\linewidth-0.5\lxhang-0.25em\relax][l]{\(C\)．\(\sqrt{13}/2\)}\makebox[\dimexpr0.5\linewidth-0.5\lxhang-0.25em\relax][l]{\(D\)．\(4/3\)}\par}
\begin{ansblock}[2章-拓-课时11-16]
% ans:2章-拓-课时11-16
\ansitem{16}{\(B\)}
\ansnote{详解}{内心到三边距离\(=\)内切圆半径\(r\)；\(I(x_{0},2)\)到\(x\)轴（边\(F_{1}F_{2}\)所在直线）距离为\(2\)，故 \(r=2\)。面积\(=\)½\(r\times \)周长：\(4b=\)½\(\cdot 2\cdot (|MF_{1}|+|MF_{2}|+|F_{1}F_{2}|)=2a+2c \Rightarrow  a+c=2b\)。平方：\((a+c)^{2}=4b^{2}=4(a^{2}-c^{2})=4(a-c)(a+c) \Rightarrow  a+c=4(a-c) \Rightarrow  5c=3a\)，所求\(=(2a)/(2c)=a/c=5/3\)。选\(B\)。验算：\(a/c=5/3\)时取\(a=5k\)、\(c=3k\)、\(b=4k\)：\(a+c=8k=2b=8k\)  面积式闭合。}
\end{ansblock}

\zux{方程表示椭圆的条件（含参三段分类）}
% 席注（账面）："焦点位置看大分母"为本课时2.5.1口径（§0），无前引。
\tuhao{17} \tieside{中}已知\(x^{2}/(m-3)+y^{2}/(11-m)=1\)，当\(m\)为何值时，\((1)\)方程表示椭圆；\((2)\)方程表示焦点在\(x\)轴上的椭圆；\((3)\)方程表示焦点在\(y\)轴上的椭圆．
\liubai[16mm]{解答书写区}
\begin{ansblock}[2章-拓-课时11-17]
% ans:2章-拓-课时11-17
\ansitem{17}{\((1)3<m<7\)或\(7<m<11\)；\((2)7<m<11\)；\((3)3<m<7\)}
\ansnote{详解}{\((1)\)椭圆标准方程要求两分母同正且不相等（相等即圆）：\(m-3>0\)、\(11-m>0\)、\(m-3\neq 11-m \Rightarrow  3<m<11\) 且 \(m\neq 7\)。\((2)\)焦点在\(x\)轴\(\Leftrightarrow x^{2}\)项分母大：\(m-3>11-m>0 \Rightarrow  m>7\) 且 \(m<11\)。\((3)\)焦点在\(y\)轴\(\Leftrightarrow y^{2}\)项分母大：\(11-m>m-3>0 \Rightarrow  3<m<7\)。}
\end{ansblock}

\jietitle{2.5.2　椭圆的几何性质（课时12·拓区）}
% 【12-拓27｜撤席登记（0914 主裁：保 12-拓6 撤 12-拓27，同题双收）】本席不设题——跳号留痕（拓28 起），注记随席（台账抄件）；防双收 B23。
\tuhao{01} \tieside{简}求下列椭圆的焦点坐标：\((1) x^{2}/9+y^{2}/4=1\)；\((2) x^{2}/3+y^{2}/9=1\)；\((3) 16x^{2}+7y^{2}=112\)；\((4) x^{2}+2y^{2}=1\)．
\liubai[16mm]{解答书写区}
\begin{ansblock}[2章-拓-课时12-01]
% ans:2章-拓-课时12-01
\ansitem{01}{\((1)(\sqrt{5},0)\)、\((-\sqrt{5},0)\)；\((2)(0,\sqrt{6})\)、\((0,-\sqrt{6})\)；\((3)(0,-3)\)、\((0,3)\)；\((4)(\sqrt{2}/2,0)\)、\((-\sqrt{2}/2,0)\)}
\ansnote{详解}{化标准式后比较分母定焦点所在轴，按 \(c^{2}=a^{2}-b^{2}\) 求 \(c\)。 \((1)a^{2}=9\)、\(b^{2}=4\)，\(c^{2}=5\)，焦点在 \(x\) 轴：\((\pm \sqrt{5},0)\)； \((2)a^{2}=9\)、\(b^{2}=3\)，\(c^{2}=6\)，焦点在 \(y\) 轴：\((0,\pm \sqrt{6})\)； \((3)\)化为 \(x^{2}/7+y^{2}/16=1\)，\(a^{2}=16\)、\(b^{2}=7\)，\(c^{2}=9\)，焦点在 \(y\) 轴：\((0,\pm 3)\)； \((4)\)化为 \(x^{2}+y^{2}/(1/2)=1\)，\(a^{2}=1\)、\(b^{2}=1/2\)，\(c^{2}=1/2\)，\(c=\sqrt{2}/2\)，焦点在 \(x\) 轴：\((\pm \sqrt{2}/2,0)\)．}
\end{ansblock}

\tuhao{02} \tieside{中}已知椭圆 \(x^{2}/a^{2}+y^{2}/b^{2}=1 (a>b>0)\) 的左右两个焦点分别为 \(F_{1}\)，\(F_{2}\)，以 \(F_{1}F_{2}\) 为斜边的等腰直角三角形 \(PF_{1}F_{2}\) 与椭圆有两个不同的交点 \(M\)，\(N\)，且 \(MN=(1/3)F_{1}F_{2}\)，则该椭圆的离心率为\kongda{\(\sqrt{5}-\sqrt{2}\)}．
\begin{ansblock}[2章-拓-课时12-02]
% ans:2章-拓-课时12-02
\ansitem{02}{\(\sqrt{5}-\sqrt{2}\)}
\ansnote{详解}{以 \(O\) 为原点、\(F_{1}F_{2}\) 为 \(x\) 轴，则 \(F_{1}(-c,0)\)、\(F_{2}(c,0)\)，等腰直角三角形顶点 \(P(0,c)\)（\(|OP|=\)½\(|F_{1}F_{2}|=c\)）． 斜边 \(F_{1}F_{2}\) 外的两腰 \(PF_{1}\)：\(y=x+c\)（\(x\in [-c,0]\)）、\(PF_{2}\)：\(x+y=c\)（\(x\in [0,c]\)）．由对称性，\(M\)、\(N\) 分别在两腰上且关于 \(y\) 轴对称，\(MN\) 水平，\(|MN|=(1/3)|F_{1}F_{2}|=2c/3\)，故 \(N\) 的横坐标为 \(c/3\)、\(M\) 的横坐标为 \(-c/3\)， 代入相应腰的方程得 \(N(c/3, 2c/3)\)、\(M(-c/3, 2c/3)\)． 由椭圆定义：\(2a=|NF_{1}|+|NF_{2}|=\sqrt{}((c/3+c)^{2}+(2c/3)^{2})+\sqrt{}((c/3-c)^{2}+(2c/3)^{2})=\sqrt{16c^{2}/9+4c^{2}/9}+\sqrt{4c^{2}/9+4c^{2}/9}=(2\sqrt{5}/3)c+(2\sqrt{2}/3)c\)， 故 \(e=c/a=3/(\sqrt{5}+\sqrt{2})=3(\sqrt{5}-\sqrt{2})/(5-2)=\sqrt{5}-\sqrt{2}\)． （数值核验：\(\sqrt{5}-\sqrt{2}\approx 2.236-1.414=0.822\in (0,1)\) ）}
\end{ansblock}

\tuhao{03} \tieside{中}设圆锥曲线 \(\Gamma \) 的两个焦点分别为 \(F_{1}\)，\(F_{2}\)．若曲线 \(\Gamma \) 上存在点 \(P\) 满足 \(|PF_{1}|:|F_{1}F_{2}|:|PF_{2}|=4:3:2\)，则曲线 \(\Gamma \) 的离心率等于\kongda{\(1/2\) 或 \(3/2\)}．
\begin{ansblock}[2章-拓-课时12-03]
% ans:2章-拓-课时12-03
\ansitem{03}{\(1/2\) 或 \(3/2\)}
\ansnote{详解}{设 \(|PF_{1}|=4t\)、\(|F_{1}F_{2}|=2c=3t\)、\(|PF_{2}|=2t\)（\(t>0\)），则 \(c=3t/2\)。 椭圆支：\(2a=|PF_{1}|+|PF_{2}|=6t \Rightarrow  a=3t\)，\(e=c/a=(3t/2)/(3t)=1/2\) （且 \(2a>2c\)：\(6t>3t\) ，可成椭圆）； 双曲线支：\(2a=||PF_{1}|-|PF_{2}||=2t \Rightarrow  a=t\)，\(e=c/a=(3t/2)/t=3/2\) （\(2a<2c\)：\(2t<3t\) ，可成双曲线）．}
\end{ansblock}

\tuhao{04} \tieside{中}设椭圆 \(C\)：\(x^{2}/a^{2}+y^{2}/b^{2}=1\)（\(a>b>0\)），长轴的两个端点分别为 \(A_{1}\)，\(A_{2}\)，短轴的两个端点分别为 \(B_{1}\)，\(B_{2}\)．
\((1)\)证明：四边形 \(A_{1}B_{1}A_{2}B_{2}\) 为菱形；\((2)\)若四边形 \(A_{1}B_{1}A_{2}B_{2}\) 的面积为\(120\)，边长为\(13\)，求椭圆 \(C\) 的方程．
\liubai[16mm]{解答书写区}
\begin{ansblock}[2章-拓-课时12-04]
% ans:2章-拓-课时12-04
\ansitem{04}{\((1)\)证明见详解；\((2)x^{2}/144+y^{2}/25=1\)}
\ansnote{详解}{\((1)A_{1}(-a,0)\)、\(A_{2}(a,0)\)、\(B_{1}(0,-b)\)、\(B_{2}(0,b)\)，两对角线 \(A_{1}A_{2}\) 与 \(B_{1}B_{2}\) 同为对方的中垂线（互相垂直平分），故 \(A_{1}B_{1}A_{2}B_{2}\) 是平行四边形且对角线垂直，即为菱形． \((2)\)菱形面积\(=\)½\(\cdot |A_{1}A_{2}|\cdot |B_{1}B_{2}|=\)½\(\cdot 2a\cdot 2b=2ab=120 \Rightarrow  ab=60\)；相邻顶点 \(A_{2}(a,0)\)、\(B_{2}(0,b)\) 的距离即边长 \(\sqrt{a^{2}+b^{2}}=13 \Rightarrow  a^{2}+b^{2}=169\)． 联立 \(a^{2}+b^{2}=169\) 与 \(ab=60\)：\(a^{2}\)、\(b^{2}\) 是 \(u^{2}-169u+3600=0\) 的两根，\(\Delta =169^{2}-4\cdot 3600=28561-14400=14161=119^{2}\)， \(u=(169\pm 119)/2=144\) 或 \(25\)，由 \(a>b\) 得 \(a^{2}=144\)、\(b^{2}=25\)，椭圆 \(x^{2}/144+y^{2}/25=1\)．}
\end{ansblock}

\tuhao{05} \tieside{中}已知 \(F_{1}\)，\(F_{2}\) 为椭圆 \(C\)：\(x^{2}/16+y^{2}/4=1\) 的两个焦点，\(P\)，\(Q\) 为 \(C\) 上关于坐标原点对称的两点，且 \(|PQ|=|F_{1}F_{2}|\)，则四边形 \(PF_{1}QF_{2}\) 的面积为\kongda{\(8\)}．
\begin{ansblock}[2章-拓-课时12-05]
% ans:2章-拓-课时12-05
\ansitem{05}{\(8\)}
\ansnote{详解}{\(P\)、\(Q\) 关于 \(O\) 对称，\(F_{1}\)、\(F_{2}\) 关于 \(O\) 对称，故 \(PF_{1}QF_{2}\) 的对角线 \(PQ\) 与 \(F_{1}F_{2}\) 互相平分（是平行四边形）；又 \(|PQ|=|F_{1}F_{2}|\)，对角线相等的平行四边形为矩形，于是 \(\angle F_{1}PF_{2}=90^{\circ}\)． \(a^{2}=16\)、\(b^{2}=4 \Rightarrow  c^{2}=12\)，\(|F_{1}F_{2}|^{2}=4c^{2}=48\)．设 \(|PF_{1}|=m\)、\(|PF_{2}|=n\)，则 \(m+n=2a=8\)，\(m^{2}+n^{2}=48\)， \(\Rightarrow  2mn=(m+n)^{2}-(m^{2}+n^{2})=64-48=16 \Rightarrow  mn=8\)． 矩形面积\(=mn=8\)．}
\end{ansblock}

\tuhao{06} \tieside{中}若 \(P\) 为椭圆 \(x^{2}/25+y^{2}/(25/2)=1\) 上的一点，\(F_{1}\)，\(F_{2}\) 分别是椭圆的左、右焦点，则 \(\angle F_{1}PF_{2}\) 的最大值为（　　）
\bindopt {\fontsize{10.5pt}{\qpTJgLeadLiti}\selectfont \makebox[\dimexpr0.25\linewidth-0.25\lxhang-0.25em\relax][l]{\(A\)．\(30^{\circ}\)}\makebox[\dimexpr0.25\linewidth-0.25\lxhang-0.25em\relax][l]{\(B\)．\(45^{\circ}\)}\makebox[\dimexpr0.25\linewidth-0.25\lxhang-0.25em\relax][l]{\(C\)．\(60^{\circ}\)}\makebox[\dimexpr0.25\linewidth-0.25\lxhang-0.25em\relax][l]{\(D\)．\(90^{\circ}\)}\par}
\begin{ansblock}[2章-拓-课时12-06]
% ans:2章-拓-课时12-06
\ansitem{06}{\(D\)}
\ansnote{详解}{\(a^{2}=25\)、\(b^{2}=25/2\)，\(c^{2}=a^{2}-b^{2}=25/2\)，\(e^{2}=c^{2}/a^{2}=1/2\)，故 \(b=c=5\sqrt{2}/2\)． 设 \(P(x,y)\)，\(x^{2}\in [0,25]\)．由焦半径 \(|PF_{1}|=a+ex\)、\(|PF_{2}|=a-ex\) 得 \(|PF_{1}|^{2}+|PF_{2}|^{2}=(a+ex)^{2}+(a-ex)^{2}=2a^{2}+2e^{2}x^{2}=50+x^{2}\)，\(|PF_{1}|\cdot |PF_{2}|=a^{2}-e^{2}x^{2}=25-x^{2}/2\)，且 \(|F_{1}F_{2}|^{2}=4c^{2}=50\)， 故 \(cos\angle F_{1}PF_{2}=(|PF_{1}|^{2}+|PF_{2}|^{2}-|F_{1}F_{2}|^{2})/(2|PF_{1}|\cdot |PF_{2}|)=(50+x^{2}-50)/[2(25-x^{2}/2)]=x^{2}/(50-x^{2})\geqslant 0\)， 分母 \(50-x^{2}\geqslant 25>0\) 恒正，等号当且仅当 \(x=0\)，即 \(P\) 为短轴端点 \((0,\pm 5\sqrt{2}/2)\)． \(cos\angle F_{1}PF_{2}\geqslant 0\) 说明 \(\angle F_{1}PF_{2}\leqslant 90^{\circ}\)，且 \(cos\) 在 \([0,\pi ]\) 上随角单调递减，故 \(cos\) 取最小值 \(0\) 时顶角最大： \(x=0\) 时 \(|PF_{1}|=|PF_{2}|=\sqrt{b^{2}+c^{2}}=\sqrt{25/2+25/2}=5=a\)，\(|F_{1}F_{2}|=2c=5\sqrt{2}\)，\(5^{2}+5^{2}=50=(5\sqrt{2})^{2} \Rightarrow  \angle F_{1}PF_{2}=90^{\circ}\)； \(x^{2}=25\)（长轴端点，三点共线）时 \(cos=25/(50-25)=1\)，顶角 \(0^{\circ}\)，与图形一致． 所以 \(\angle F_{1}PF_{2}\) 的最大值为 \(90^{\circ}\)．故选：\(D\)．}
\end{ansblock}

\tuhao{07} \tieside{中}已知椭圆 \(x^{2}/4+y^{2}=1\) 的两个焦点为 \(F_{1}\)，\(F_{2}\)，\(P(x,y)\) 为椭圆上任意一点，求使 \(\angle F_{1}PF_{2}\geqslant 90^{\circ}\) 的 \(x\) 的取值范围．
\liubai[16mm]{解答书写区}
\begin{ansblock}[2章-拓-课时12-07]
% ans:2章-拓-课时12-07
\ansitem{07}{\([-2\sqrt{6}/3, 2\sqrt{6}/3]\)}
\ansnote{详解}{\(c^{2}=4-1=3\)，\(F_{1}(-\sqrt{3},0)\)、\(F_{2}(\sqrt{3},0)\)．\(P\) 不在 \(x\) 轴上的线段 \(F_{1}F_{2}\) 内（椭圆上点 \(|y|\) 可取\(0\)但此时 \(x=\pm 2\)，在 \(F_{1}F_{2}\) 之外），故 \(\angle F_{1}PF_{2}\geqslant 90^{\circ} \Leftrightarrow \) 向量 \(PF_{1}\cdot PF_{2}\leqslant 0 \Leftrightarrow  (-\sqrt{3}-x)(\sqrt{3}-x)+y^{2}\leqslant 0 \Leftrightarrow  x^{2}+y^{2}-3\leqslant 0\)，即 \(P\) 在以原点为圆心、\(\sqrt{3}\) 为半径的圆内（含界）． 以 \(y^{2}=1-x^{2}/4\) 代入：\(x^{2}+1-x^{2}/4-3\leqslant 0 \Rightarrow  (3/4)x^{2}\leqslant 2 \Rightarrow  x^{2}\leqslant 8/3 \Rightarrow  |x|\leqslant 2\sqrt{6}/3\)， 又 \(2\sqrt{6}/3\approx 1.633<2\)  在椭圆 \(x\) 范围内，故 \(x\in [-2\sqrt{6}/3, 2\sqrt{6}/3]\)．}
\end{ansblock}

% 图债注：「如图」指源件配图，印面无图（冻片【注】裁定：去装饰图/条件全在文字/尺寸随注——图债登记汇编报告）
\tuhao{08} \tieside{中}如图所示，\(A_{1}\)，\(A_{2}\) 是椭圆 \(C\)：\(x^{2}/18+y^{2}/9=1\) 的短轴端点，点 \(M\) 在椭圆上运动，且点 \(M\) 不与 \(A_{1}\)，\(A_{2}\) 重合，点 \(N\) 满足 \(NA_{1}\perp MA_{1}\)，\(NA_{2}\perp MA_{2}\)，则 \(S\triangle MA_{1}A_{2}/S\triangle NA_{1}A_{2}=\)（　　）
\lxopt{\(A\)．\(2\)}
\lxopt{\(B\)．\(3\)}
\lxopt{\(C\)．\(4\)}
\lxopt{\(D\)．\(5/2\) 【图注】短轴在 \(y\) 轴上（\(a^{2}=18\) 在 \(x\) 项），\(A_{1}(0,3)\)、\(A_{2}(0,-3)\)；\(M\) 在椭圆上，\(N\) 为两垂线交点。}
\begin{ansblock}[2章-拓-课时12-08]
% ans:2章-拓-课时12-08
\ansitem{08}{\(A\)}
\ansnote{详解}{\(a^{2}=18\)、\(b^{2}=9\)，短轴端点 \(A_{1}(0,3)\)、\(A_{2}(0,-3)\)．设 \(M(x_{0},y_{0})\)（\(x_{0}\neq 0\)），\(N(x_{1},y_{1})\)． \(k_{MA_{1}}=(y_{0}-3)/x_{0}\)，\(NA_{1}\perp MA_{1} \Rightarrow  k_{NA_{1}}=-x_{0}/(y_{0}-3)\)，\(NA_{1}\)：\(y=-x_{0}/(y_{0}-3)\cdot x+3\)； 同理 \(NA_{2}\)：\(y=-x_{0}/(y_{0}+3)\cdot x-3\)． 两式相减消 \(y\)：\(0=-x_{0}x[1/(y_{0}-3)-1/(y_{0}+3)]+6 \Rightarrow  x_{0}x\cdot 6/(y_{0}^{2}-9)=6 \Rightarrow  x_{1}=(y_{0}^{2}-9)/x_{0}\)． 由 \(M\) 在椭圆上：\(x_{0}^{2}/18+y_{0}^{2}/9=1 \Rightarrow  y_{0}^{2}-9=-x_{0}^{2}/2\)，故 \(x_{1}=-x_{0}/2\)． 两三角形共底 \(A_{1}A_{2}\)（在 \(y\) 轴上），面积比等于高的比\(=|x_{0}|/|x_{1}|=2\)．故选：\(A\)．}
\end{ansblock}

\tuhao{09} \tieside{中}已知动点 \(P(x,y)\) 在椭圆 \(x^{2}/25+y^{2}/16=1\) 上，若 \(A\) 点坐标为 \((3,0)\)，\(|AM|=1\)，且向量\(PM\cdot \)向量\(AM=0\)，则 \(|PM|\) 的最小值是（　　）
\bindopt {\fontsize{10.5pt}{\qpTJgLeadLiti}\selectfont \makebox[\dimexpr0.25\linewidth-0.25\lxhang-0.25em\relax][l]{\(A\)．\(\sqrt{2}\)}\makebox[\dimexpr0.25\linewidth-0.25\lxhang-0.25em\relax][l]{\(B\)．\(\sqrt{3}\)}\makebox[\dimexpr0.25\linewidth-0.25\lxhang-0.25em\relax][l]{\(C\)．\(2\)}\makebox[\dimexpr0.25\linewidth-0.25\lxhang-0.25em\relax][l]{\(D\)．\(3\)}\par}
\begin{ansblock}[2章-拓-课时12-09]
% ans:2章-拓-课时12-09
\ansitem{09}{\(B\)}
\ansnote{详解}{向量\(PM\cdot \)向量\(AM=0 \Rightarrow  PM\perp AM\)，故 \(\triangle PMA\) 是以 \(M\) 为直角顶点的直角三角形， \(|PM|^{2}=|PA|^{2}-|AM|^{2}=|PA|^{2}-1\)，即 \(|PM|=\sqrt{|PA|^{2}-1}\)（\(|PA|>1\) 时 \(M\) 即以 \(A\) 为圆心、\(1\) 为半径的圆上自 \(P\) 所引切线的切点，必存在）， 所以 \(|PM|\) 随 \(|PA|\) 单调增大，只需求 \(|PA|\) 的最小值． 椭圆 \(a^{2}=25\)、\(b^{2}=16 \Rightarrow  c^{2}=a^{2}-b^{2}=9\)、\(c=3\)，右焦点为 \((3,0)\)，恰与点 \(A\) 重合，故 \(|PA|\) 即焦半径 \(|PF_{2}|\)， 其最小值为 \(a-c=5-3=2\)（\(P\) 为右顶点 \((5,0)\)）． 于是 \(|PM|\geqslant \sqrt{2^{2}-1}=\sqrt{3}\)，当 \(P\) 为右顶点时取等（此时 \(|PA|=2>1\)，切点 \(M\) 存在，\(|PM|=\sqrt{3}\)）．故选：\(B\)．}
\end{ansblock}

\tuhao{10} \tieside{中}若曲线 \(x=\sqrt{a-y^{2}}\) 与椭圆 \(x^{2}/9+y^{2}/7=1\) 有两个不同的交点，则 \(a\) 的取值范围是\kongda{\([7,9)\)}．
\begin{ansblock}[2章-拓-课时12-10]
% ans:2章-拓-课时12-10
\ansitem{10}{\([7,9)\)}
\ansnote{详解}{\(x=\sqrt{a-y^{2}} \Leftrightarrow  x^{2}+y^{2}=a\) 且 \(x\geqslant 0\)（须 \(a>0\)），表示以原点为圆心、\(\sqrt{a}\) 为半径的右半圆（含 \(x\) 轴上的两端点）． 椭圆 \(a_{1}^{2}=9\)、\(b_{1}^{2}=7\)，即 \(a_{1}=3\)、\(b_{1}=\sqrt{7}\)．椭圆上任一点 \(P(x,y)\) 到原点的距离平方： \(|OP|^{2}=x^{2}+y^{2}\)，由 \(x^{2}=9(1-y^{2}/7)\) 得 \(|OP|^{2}=9-(2/7)y^{2}\)，\(y^{2}\in [0,7]\)，故 \(|OP|^{2}\in [7,9]\)； 其中 \(|OP|^{2}=7\) 仅在上下顶点 \((0,\pm \sqrt{7})\)，\(|OP|^{2}=9\) 仅在左右顶点 \((\pm 3,0)\)，介于两者之间的每个值对应 \(x=\pm \)两个、\(y=\pm \)两个共四个点． 按 \(a=|OP|^{2}\) 逐段数交点（半圆只取 \(x\geqslant 0\)）： \(a<7\)：\(|OP|^{2}\) 在椭圆上最小为 \(7\)，半圆落在椭圆内部，无交点； \(a=7\)：交点 \((0,\sqrt{7})\)、\((0,-\sqrt{7})\)，恰 \(2\) 个 ； \(7<a<9\)：由 \(9-(2/7)y^{2}=a\) 得 \(y^{2}=7(9-a)/2\in (0,7)\)，两个 \(y\) 值各配 \(x=\sqrt{}(9(a-7)/2)>0\) 一个交点，共 \(2\) 个 ； \(a=9\)：仅 \(y=0\)、\(x=3\)，\(1\) 个交点（左顶点 \(x=-3\) 不在半圆上）； \(a>9\)：无交点． 综上 \(a\in [7,9)\)．}
\end{ansblock}

\tuhao{11} \tieside{中}若点 \(O\) 和点 \(F\) 分别为椭圆 \(x^{2}/4+y^{2}/3=1\) 的中心和左焦点，点 \(P\) 为椭圆上的任意一点，则向量\(OP\cdot \)向量\(FP\) 的最大值为\kongda{\(6\)}．
\begin{ansblock}[2章-拓-课时12-11]
% ans:2章-拓-课时12-11
\ansitem{11}{\(6\)}
\ansnote{详解}{\(a^{2}=4\)、\(b^{2}=3 \Rightarrow  c^{2}=1\)，故 \(F(-1,0)\)、\(O(0,0)\)．设 \(P(x,y)\)，\(x\in [-2,2]\)， 向量\(OP\cdot \)向量\(FP=(x,y)\cdot (x+1,y)=x^{2}+x+y^{2}\)． 以 \(y^{2}=3(1-x^{2}/4)=3-(3/4)x^{2}\) 代入：原式\(=x^{2}+x+3-(3/4)x^{2}=(1/4)x^{2}+x+3=(1/4)(x+2)^{2}+2\)． \(x\in [-2,2]\) 上该二次函数自对称轴 \(x=-2\) 起单调递增，故 \(x=2\)（\(P\) 为右顶点）时取最大值 \((1/4)\cdot 16+2=6\)．}
\end{ansblock}

\tuhao{12} \tieside{中}已知椭圆 \(C\)：\(x^{2}/4+y^{2}/3=1\)，\(A_{1}\)，\(A_{2}\) 为长轴的两个端点，点 \(P\) 是椭圆上的一点，且满足直线 \(PA_{1}\) 的斜率的取值范围是 \([1,2]\)，则直线 \(PA_{2}\) 的斜率的取值范围是\kongda{\([-3/4, -3/8]\)}．
\begin{ansblock}[2章-拓-课时12-12]
% ans:2章-拓-课时12-12
\ansitem{12}{\([-3/4, -3/8]\)}
\ansnote{详解}{\(a=2\)，\(A_{1}(-2,0)\)、\(A_{2}(2,0)\)．设 \(P(m,n)\)（\(m\neq \pm 2\)，两斜率均存在且非零），则 \(n^{2}=3(1-m^{2}/4)=3(4-m^{2})/4\)， \(k_{PA_{1}}\cdot k_{PA_{2}}=[n/(m+2)]\cdot [n/(m-2)]=n^{2}/(m^{2}-4)=[3(4-m^{2})/4]/(m^{2}-4)=-3/4\)， 即第三定义的斜率积定值 \(-b^{2}/a^{2}=-3/4\)． 由 \(1\leqslant k_{PA_{1}}\leqslant 2\) 且 \(k_{PA_{2}}=-3/(4k_{PA_{1}})\)；\(k_{PA_{1}}>0\) 时 \(k\mapsto -3/(4k)\) 单调递增， \(k_{PA_{1}}=1 \Rightarrow  k_{PA_{2}}=-3/4\)；\(k_{PA_{1}}=2 \Rightarrow  k_{PA_{2}}=-3/8\)， 故 \(k_{PA_{2}}\in [-3/4, -3/8]\)．}
\end{ansblock}

\tuhao{13} \tieside{中}点 \(P\) 在圆 \(C\)：\(x^{2}+(y-2)^{2}=1/4\) 上移动，点 \(Q\) 在椭圆 \(x^{2}+4y^{2}=4\) 上移动，则 \(|PQ|\) 的最大值为\(\kongbai{}\)，相应的点 \(Q\) 的坐标为\(\kongbai{}\)．
\begin{ansblock}[2章-拓-课时12-13]
% ans:2章-拓-课时12-13
\ansitem{13}{\(1/2+2\sqrt{21}/3\)；\((\pm 2\sqrt{5}/3, -2/3)\)}
\ansnote{详解}{椭圆化为 \(x^{2}/4+y^{2}=1\)，设 \(Q(2cos\alpha , sin\alpha )\)；圆心 \(C(0,2)\)、半径 \(1/2\)． 三角不等式 \(|PQ|\leqslant |PC|+|CQ|=1/2+|CQ|\)，等号当 \(P\) 为半径 \(CQ\) 延长线与圆的交点时成立，故只需求 \(|CQ|\) 的最大值． \(|CQ|^{2}=4cos^{2}\alpha +(sin\alpha -2)^{2}=4(1-sin^{2}\alpha )+sin^{2}\alpha -4sin\alpha +4=8-4sin\alpha -3sin^{2}\alpha =-3(sin\alpha +2/3)^{2}+28/3\leqslant 28/3\)， 等号当 \(sin\alpha =-2/3\)（可取）时成立，此时 \(cos\alpha =\pm \sqrt{1-4/9}=\pm \sqrt{5}/3\)， \(|CQ|max=\sqrt{28/3}=2\sqrt{7}/\sqrt{3}=2\sqrt{21}/3\)，相应 \(Q(\pm 2\sqrt{5}/3, -2/3)\)． 故 \(|PQ|max=1/2+2\sqrt{21}/3\)，相应点 \(Q\) 为 \((\pm 2\sqrt{5}/3, -2/3)\)．}
\end{ansblock}

% 图债注：「如图」指源件配图，印面无图（冻片【注】裁定：去装饰图/条件全在文字/尺寸随注——图债登记汇编报告）
\tuhao{14} \tieside{中}人造地球卫星绕地球运行遵循开普勒行星运动定律：如图，卫星在以地球的中心为焦点的椭圆轨道上绕地球运行时，其运行速度是变化的，速度的变化服从面积守恒规律，即卫星的向径（卫星与地心的连线）在相同的时间内扫过的面积相等．设该椭圆的长轴长、焦距分别为 \(2a\)，\(2c\)．某同学根据所学知识，得到下列结论：
①卫星向径的取值范围是 \([a-c, a+c]\)；
②卫星向径的最小值与最大值的比值越大，椭圆轨道越扁；
③卫星在左半椭圆弧的运行时间大于其在右半椭圆弧的运行时间；
④卫星运行速度在近地点时最小，在远地点时最大．
其中正确的结论是（　　）
\lxopt{\(A\)．①②}
\lxopt{\(B\)．①③}
\lxopt{\(C\)．②④}
\lxopt{\(D\)．①③④ 【图注】源图未随文本提取；据③为真相应读作地球（焦点）位于右焦点、近地点在右侧顶点（若焦点画在左侧，③须反向表述）。}
\begin{ansblock}[2章-拓-课时12-14]
% ans:2章-拓-课时12-14
\ansitem{14}{\(B\)}
\ansnote{详解}{①向径即焦半径，\(|PF|=a-ex\)（设焦点为右焦点、\(x\in [-a,a]\)），故取值范围为 \([a-c, a+c]\)，①正确． ②比值为 \((a-c)/(a+c)=(1-e)/(1+e)\)，对 \(e\) 求导得 \([(1-e)/(1+e)]'=-2/(1+e)^{2}<0\)，故比值越大 \(\Leftrightarrow  e\) 越小 \(\Leftrightarrow \) 轨道越圆，②错误． ③以焦点 \(F(c,0)\) 为顶点看两半弧扫过的面积：左半弧（\(x\leqslant 0\) 段）扫过面积\(=\)半个椭圆面积\(+\triangle F B_{1}B_{2}\)（\(B_{1}\)、\(B_{2}\) 为短轴端点）\(=\pi ab/2+bc\)；右半弧扫过面积\(=\pi ab/2-bc\)．由面积守恒，扫过面积与用时成正比，左半弧用时长，③正确． ④近地点向径短、单位时间扫过同面积须弧长长，故速度最大；远地点速度最小，④所说相反，错误． 综上正确的是①③．故选：\(B\)．}
\end{ansblock}

\tuhao{15} \tieside{中}定义离心率是 \((\sqrt{5}-1)/2\) 的椭圆为\("\)黄金椭圆\("\)．已知椭圆 \(E\)：\(x^{2}/10+y^{2}/m=1 (10>m>0)\) 是\("\)黄金椭圆\("\)，则 \(m=\)\kongda{\(5\sqrt{5}-5\)；\((\sqrt{5}+1)/2\)}；若\("\)黄金椭圆\(" C\)：\(x^{2}/a^{2}+y^{2}/b^{2}=1 (a>b>0)\) 两个焦点分别为 \(F_{1}(-c,0)\)、\(F_{2}(c,0) (c>0)\)，\(P\) 为椭圆 \(C\) 上的异于顶点的任意一点，点 \(M\) 是 \(\triangle PF_{1}F_{2}\) 的内心，连接 \(PM\) 并延长交 \(F_{1}F_{2}\) 于点 \(N\)，则 \(|PM|/|MN|=\kongbai{}\)．
\begin{ansblock}[2章-拓-课时12-15]
% ans:2章-拓-课时12-15
\ansitem{15}{\(5\sqrt{5}-5\)；\((\sqrt{5}+1)/2\)}
\ansnote{详解}{第一空：\(10>m>0\) 时焦点在 \(x\) 轴上，\(e^{2}=1-m/10\)．又 \(e=(\sqrt{5}-1)/2 \Rightarrow  e^{2}=(6-2\sqrt{5})/4=(3-\sqrt{5})/2\)， 故 \(m/10=1-(3-\sqrt{5})/2=(\sqrt{5}-1)/2\)，\(m=5(\sqrt{5}-1)=5\sqrt{5}-5\)． 第二空：设 \(\triangle PF_{1}F_{2}\) 的内切圆半径为 \(r\)，则 \(M\) 到三边距离均为 \(r\)． \(S\triangle PF_{1}F_{2}=(1/2)r(|PF_{1}|+|PF_{2}|+|F_{1}F_{2}|)=(1/2)r(2a+2c)=(a+c)r\)； \(S\triangle MF_{1}F_{2}=(1/2)r|F_{1}F_{2}|=cr\)． \(\triangle MF_{1}F_{2}\) 与 \(\triangle PF_{1}F_{2}\) 同底 \(F_{1}F_{2}\)，面积比等于高之比；又 \(P\)、\(M\)、\(N\) 共线且 \(N\) 在 \(F_{1}F_{2}\) 上，故高之比\(=|MN|/|PN|\)， 于是 \(|MN|/|PN|=cr/[(a+c)r]=c/(a+c)\)， \(|PM|=|PN|-|MN|=[1-c/(a+c)]|PN|=a/(a+c)\cdot |PN|\)， 所以 \(|PM|/|MN|=a/c=1/e=2/(\sqrt{5}-1)=(\sqrt{5}+1)/2\)．}
\end{ansblock}

\tuhao{16}\duoxuan\hspace{0.6em} \tieside{难}曲率半径是用来描述曲线上某点处曲线弯曲变化程度的量，已知对于曲线 \(x^{2}/a^{2}+y^{2}/b^{2}=1 (a>0,b>0)\) 上点 \(P(x_{0},y_{0})\) 处的曲率半径公式为 \(R=a^{2}b^{2}(x_{0}^{2}/a^{4}+y_{0}^{2}/b^{4})^(3/2)\)，则下列说法正确的是（　　）
\lxopt{\(A\)．对于半径为 \(R\) 的圆，其圆上任一点的曲率半径均为 \(R\)}
\lxopt{\(B\)．椭圆 \(x^{2}/a^{2}+y^{2}/b^{2}=1 (a>b>0)\) 上一点处的曲率半径的最大值为 \(a\)}
\lxopt{\(C\)．椭圆 \(x^{2}/a^{2}+y^{2}/b^{2}=1 (a>b>0)\) 上一点处的曲率半径的最小值为 \(b^{2}/a\)}
\lxopt{\(D\)．对于椭圆 \(x^{2}/a^{2}+y^{2}=1 (a>1)\) 上点 \((1/2, y_{0})\) 处的曲率半径随着 \(a\) 的增大而减小}
\begin{ansblock}[2章-拓-课时12-16]
% ans:2章-拓-课时12-16
\ansitem{16}{\(AC\)}
\ansnote{详解}{\(A\)：圆写作 \(x^{2}/R^{2}+y^{2}/R^{2}=1\)，即 \(a^{2}=b^{2}=R^{2}\)，圆上点满足 \(x_{0}^{2}+y_{0}^{2}=R^{2}\)， \(R'=R^{4}((x_{0}^{2}+y_{0}^{2})/R^{4})^(3/2)=R^{4}(R^{2}/R^{4})^(3/2)=R^{4}/R^{3}=R\)，与点的位置无关，\(A\) 正确． \(B\)、\(C\)：记 \(f=x_{0}^{2}/a^{4}+y_{0}^{2}/b^{4}\)，令 \(u=x_{0}^{2}/a^{2}\in [0,1]\)，则 \(y_{0}^{2}/b^{2}=1-u\)，\(f=u/a^{2}+(1-u)/b^{2}=1/b^{2}+u(1/a^{2}-1/b^{2})\)， 因 \(a>b>0\) 有 \(1/a^{2}-1/b^{2}<0\)，故 \(f\) 关于 \(u\) 单调递减：\(u=0\)（短轴端点 \((0,\pm b)\)）时 \(f\) 最大\(=1/b^{2}\)，\(u=1\)（长轴端点 \((\pm a,0)\)）时 \(f\) 最小\(=1/a^{2}\)． 于是 \(R\_max=a^{2}b^{2}(1/b^{2})^(3/2)=a^{2}/b\)（在 \((0,\pm b)\) 处），\(R\_min=a^{2}b^{2}(1/a^{2})^(3/2)=b^{2}/a\)（在 \((\pm a,0)\) 处），\(R\in [b^{2}/a, a^{2}/b]\)． 故 \(B\)（最大值 \(a\)）错误，\(C\) 正确． \(D\)：\(b^{2}=1\)，\(R=a^{2}(1/(4a^{4})+y_{0}^{2})^(3/2)\)，其中 \(y_{0}^{2}=1-1/(4a^{2})\)．令 \(w=1/a^{2}\in (0,1)\)，则 \(1/(4a^{4})=w^{2}/4\)， \(R=(1/w)(1-w/4+w^{2}/4)^(3/2)\)，取对数 \(ln R=-ln w+(3/2)ln(1-w/4+w^{2}/4)\)，记 \(g(w)=1-w/4+w^{2}/4\)（\(g'=-1/4+w/2\)；\(g\) 的判别式 \(1/16-1<0\)，故 \(g>0\) 恒成立）， \(d(ln R)/dw=-1/w+(3/2)g'/g=[-g(w)+(3/2)w\cdot g'(w)]/[w\cdot g(w)]=(w^{2}/2-w/8-1)/[w\cdot g(w)]\)， \(w\in (0,1)\) 时 \(w^{2}/2<1/2\)，分子\(<-1/2-w/8<0\)，分母\(>0\)，故 \(d(ln R)/dw<0\)， 即 \(R\) 关于 \(w\) 单调递减，而 \(w=1/a^{2}\) 随 \(a\) 增大而减小，所以曲率半径 \(R\) 随 \(a\) 增大而增大，\(D\) 错误． 故选：\(AC\)．}
\end{ansblock}

% 图债注：「如图」指源件配图，印面无图（冻片【注】裁定：去装饰图/条件全在文字/尺寸随注——图债登记汇编报告）
\tuhao{17}\duoxuan\hspace{0.6em} \tieside{难}\("\)脸谱\("\)是戏曲舞台演出时的化妆造型艺术，更是中国传统戏曲文化的重要载体．如图，\("\)脸谱\("\)图形可近似看作由半圆和半椭圆组成的曲线 \(C\)，其方程为 \(C\)：\(x^{2}+y^{2}=4 (y>0)\)；\(x^{2}/4+y^{2}/9=1 (y\leqslant 0)\)．则下列说法正确的是（　　）
\lxopt{\(A\)．曲线 \(C\) 包含的封闭图形内部（不含边界）有 \(11\) 个整数点（横、纵坐标均为整数）}
\lxopt{\(B\)．曲线 \(C\) 上任意一点到原点距离的最大值与最小值之和为 \(5\)}
\lxopt{\(C\)．若 \(A(0,-\sqrt{5})\)、\(B(0,\sqrt{5})\)，\(P\) 是曲线 \(C\) 下半部分中半椭圆上的一个动点，则 \(cos\angle APB\) 的最小值为 \(-1/9\)}
\lxopt{\(D\)．画法几何的创始人加斯帕尔\(\cdot \)蒙日发现：椭圆中任意两条互相垂直的切线，其交点都在与椭圆同中心的圆上，称该圆为椭圆的蒙日圆；那么曲线 \(C\) 中下半部分半椭圆扩充为整个椭圆 \(C'\)：\(x^{2}/4+y^{2}/9=1 (-3\leqslant y\leqslant 3)\) 后，椭圆 \(C'\) 的蒙日圆方程为 \(x^{2}+y^{2}=13\)}
\begin{ansblock}[2章-拓-课时12-17]
% ans:2章-拓-课时12-17
\ansitem{17}{\(BCD\)}
\ansnote{详解}{\(A\)：封闭图形由上半圆（半径 \(2\)）与下半椭圆（横半轴 \(2\)、竖半轴 \(3\)）拼成，横坐标范围 \(-2\leqslant x\leqslant 2\)．按 \(x\in \{-2,-1,0,1,2\}\) 逐列数严格内部整数点： \(x=\pm 2\) 时 \(y=0\) 在边界上，无内部整数点； \(x=\pm 1\) 时：上半需 \(y>0\) 且 \(1+y^{2}<4 \Rightarrow  y=1\)；下半需 \(y\leqslant 0\) 且 \(1/4+y^{2}/9<1 \Rightarrow  y^{2}<27/4 \Rightarrow  y=0,-1,-2\)（\(y=-3\) 时 \(1/4+1>1\) 在形外）； 即 \((-1,1)\)、\((-1,0)\)、\((-1,-1)\)、\((-1,-2)\) 与 \((1,1)\)、\((1,0)\)、\((1,-1)\)、\((1,-2)\)，共 \(8\) 个； \(x=0\) 时：上需 \(y>0\)、\(y^{2}<4 \Rightarrow  y=1\)；下需 \(y^{2}<9 \Rightarrow  y=0,-1,-2\)，共 \((0,1)\)、\((0,0)\)、\((0,-1)\)、\((0,-2)\)，\(4\) 个． 合计 \(8+4=12\) 个整数点，\(A\) 说 \(11\) 个，错误． \(B\)：\(y>0\) 段在半圆上，到原点距离恒为 \(2\)；\(y\leqslant 0\) 段设 \(P(2cos\theta , 3sin\theta )\)，\(\theta \in [\pi ,2\pi ]\)， \(|OP|^{2}=4cos^{2}\theta +9sin^{2}\theta =9-5cos^{2}\theta \in [4,9]\)（\(\theta =\pi ,2\pi \) 时取 \(4\)，\(\theta =3\pi /2\) 时取 \(9\)），即 \(|OP|\in [2,3]\)； 故曲线上到原点距离的最大值 \(3\)、最小值 \(2\)，和为 \(5\)，\(B\) 正确． \(C\)：椭圆 \(x^{2}/4+y^{2}/9=1\) 中 \(a^{2}=9\)（在 \(y\) 项）、\(b^{2}=4\)，\(c^{2}=5\)，焦点恰为 \(A(0,-\sqrt{5})\)、\(B(0,\sqrt{5})\)，故 \(P\) 在半椭圆上时 \(|PA|+|PB|=2a=6\)，\(|AB|=2\sqrt{5}\)． \(|PA|\cdot |PB|\leqslant [(|PA|+|PB|)/2]^{2}=9\)，等号当 \(|PA|=|PB|\)，即 \(P\) 在 \(x\) 轴上（\(P(\pm 2,0)\)，属 \(y\leqslant 0\) 段）取到． \(cos\angle APB=(|PA|^{2}+|PB|^{2}-|AB|^{2})/(2|PA||PB|)=[(|PA|+|PB|)^{2}-2|PA||PB|-20]/(2|PA||PB|)=16/(2|PA||PB|)-1=8/(|PA||PB|)-1\geqslant 8/9-1=-1/9\)， 故 \(cos\angle APB\) 的最小值为 \(-1/9\)，\(C\) 正确． \(D\)：由蒙日圆半径公式 \(r^{2}=a^{2}+b^{2}=9+4=13\)（亦如源取两条特殊切线 \(x=2\)、\(y=3\)，交点 \((2,3)\) 在蒙日圆上，\(r^{2}=4+9=13\)），蒙日圆为 \(x^{2}+y^{2}=13\)，\(D\) 正确． 故选：\(BCD\)．}
\end{ansblock}

\tuhao{18} \tieside{难}（本小题满分 \(16\) 分）已知圆 \(C\)：\(x^{2}+y^{2}=r^{2}\) 有以下性质：
①过圆 \(C\) 上一点 \(M(x_{0},y_{0})\) 的圆的切线方程是 \(x_{0}x+y_{0}y=r^{2}\)；
②若 \(M(x_{0},y_{0})\) 为圆 \(C\) 外一点，过 \(M\) 作圆 \(C\) 的两条切线，切点分别为 \(A\)，\(B\)，则直线 \(AB\) 的方程为 \(x_{0}x+y_{0}y=r^{2}\)；
③若不在坐标轴上的点 \(M(x_{0},y_{0})\) 为圆 \(C\) 外一点，过 \(M\) 作圆 \(C\) 的两条切线，切点分别为 \(A\)，\(B\)，则 \(OM\) 垂直 \(AB\)，即 \(k_{AB}\cdot k_{OM}=-1\)，且 \(OM\) 平分线段 \(AB\)．
\((1)\) 类比上述有关结论，猜想过椭圆 \(C'\)：\(x^{2}/a^{2}+y^{2}/b^{2}=1\) 上一点 \(M(x_{0},y_{0})\) 的切线方程（不要求证明）；
\((2)\) 过椭圆 \(C'\)：\(x^{2}/a^{2}+y^{2}/b^{2}=1\) 外一点 \(M(x_{0},y_{0})\) 作两直线，与椭圆相切于 \(A\)，\(B\) 两点，求过 \(A\)，\(B\) 两点的直线方程；
\((3)\) 若过椭圆 \(C'\)：\(x^{2}/a^{2}+y^{2}/b^{2}=1\) 外一点 \(M(x_{0},y_{0})\)（\(M\) 不在坐标轴上）作两直线，与椭圆相切于 \(A\)，\(B\) 两点，求证：\(k_{AB}\cdot k_{OM}\) 为定值，且 \(OM\) 平分线段 \(AB\)．
\liubai[32mm]{解答书写区}
\begin{ansblock}[2章-拓-课时12-18]
% ans:2章-拓-课时12-18
\ansitem{18}{\((1) x_{0}x/a^{2}+y_{0}y/b^{2}=1\)；\((2) x_{0}x/a^{2}+y_{0}y/b^{2}=1\)；\((3)\) 定值为 \(-b^{2}/a^{2}\)，证明见详解}
\ansnote{详解}{\((1)\) 类比①：圆方程 \(x^{2}/R^{2}+y^{2}/R^{2}=1\) 中 \(r^{2}\) 同时充当 \(a^{2}\)、\(b^{2}\)，切线 \(x_{0}x+y_{0}y=r^{2}\) 即 \(x_{0}x/R^{2}+y_{0}y/R^{2}=1\)，故把 \(r^{2}\) 对应替换为 \(a^{2}\)、\(b^{2}\)，猜想过 \(M(x_{0},y_{0})\) 的切线方程为 \(x_{0}x/a^{2}+y_{0}y/b^{2}=1\)． \((2)\) 设 \(A(x_{1},y_{1})\)、\(B(x_{2},y_{2})\)．由 \((1)\)，过椭圆上点 \(A\) 的切线 \(l_{1}\) 为 \(x_{1}x/a^{2}+y_{1}y/b^{2}=1\)；\(l_{1}\) 过 \(M(x_{0},y_{0})\)，故 \(x_{1}x_{0}/a^{2}+y_{1}y_{0}/b^{2}=1\)．同理 \(x_{2}x_{0}/a^{2}+y_{2}y_{0}/b^{2}=1\)． 两式说明 \(A\)、\(B\) 的坐标都满足同一方程 \(x_{0}x/a^{2}+y_{0}y/b^{2}=1\)，而两点确定一条直线，故直线 \(AB\) 的方程为 \(x_{0}x/a^{2}+y_{0}y/b^{2}=1\)． \((3)\) 由 \((2)\)，\(AB\)：\(x_{0}x/a^{2}+y_{0}y/b^{2}=1\)，因 \(M\) 不在坐标轴上，\(x_{0}y_{0}\neq 0\)，故 \(k_{AB}=-(b^{2}x_{0})/(a^{2}y_{0})\)；又 \(k_{OM}=y_{0}/x_{0}\)， 所以 \(k_{AB}\cdot k_{OM}=-(b^{2}x_{0})/(a^{2}y_{0})\cdot y_{0}/x_{0}=-b^{2}/a^{2}\)，为定值（与切点位置无关）． 设 \(AB\) 的中点为 \(P((x_{1}+x_{2})/2, (y_{1}+y_{2})/2)\)．\(A\)、\(B\) 均在椭圆上：\(x_{1}^{2}/a^{2}+y_{1}^{2}/b^{2}=1\)，\(x_{2}^{2}/a^{2}+y_{2}^{2}/b^{2}=1\)，两式相减得 \((x_{2}^{2}-x_{1}^{2})/a^{2}+(y_{2}^{2}-y_{1}^{2})/b^{2}=0\)，即 \((x_{2}+x_{1})(x_{2}-x_{1})/a^{2}+(y_{2}+y_{1})(y_{2}-y_{1})/b^{2}=0\)， 于是 \(k_{AB}\cdot k_{OP}=[(y_{2}-y_{1})/(x_{2}-x_{1})]\cdot [(y_{2}+y_{1})/(x_{2}+x_{1})]=-b^{2}/a^{2}\)． 结合已证的 \(k_{AB}\cdot k_{OM}=-b^{2}/a^{2}\)，且 \(k_{AB}\neq 0\)（\(M\) 不在坐标轴上使 \(x_{0}y_{0}\neq 0\)），得 \(k_{OP}=k_{OM}\)， 故 \(O\)、\(P\)、\(M\) 三点共线，即直线 \(OM\) 过弦 \(AB\) 的中点 \(P\)，所以 \(OM\) 平分线段 \(AB\)．}
\end{ansblock}

\tuhao{19} \tieside{难}设椭圆 \(C\)：\(x^{2}/a^{2}+y^{2}/b^{2}=1 (a>b>0)\) 的离心率为 \(\sqrt{2}/2\)，且经过点 \((-2,1)\)．
\((1)\) 求椭圆 \(C\) 的标准方程；
\((2)\) 设 \(P(m,n)\) 为椭圆 \(C\) 外一动点，过点 \(P\) 作椭圆 \(C\) 的两条互相垂直的切线 \(PA\)，\(PB\)，\(A\)、\(B\) 为切点，求动点 \(P\) 的轨迹方程．
\liubai[24mm]{解答书写区}
\begin{ansblock}[2章-拓-课时12-19]
% ans:2章-拓-课时12-19
\ansitem{19}{\((1) x^{2}/6+y^{2}/3=1\)；\((2) x^{2}+y^{2}=9\)}
\ansnote{详解}{\((1) e^{2}=c^{2}/a^{2}=1/2\)，又 \(a^{2}=b^{2}+c^{2}\)，故 \(a^{2}=2b^{2}\)．点 \((-2,1)\) 在椭圆上：\(4/a^{2}+1/b^{2}=1\)，即 \(4/(2b^{2})+1/b^{2}=3/b^{2}=1\)，得 \(b^{2}=3\)、\(a^{2}=6\)， 椭圆 \(C\)：\(x^{2}/6+y^{2}/3=1\)． \((2)\)（ⅰ）两切线中有一条斜率不存在时：该切线为 \(x=\pm \sqrt{6}\)，另一条与之垂直且与椭圆相切，只能是 \(y=\pm \sqrt{3}\)，交点 \(P\) 为 \((\pm \sqrt{6}, \pm \sqrt{3})\)（四组），均满足 \(6+3=9\)． （ⅱ）两切线斜率均存在时，设切线为 \(y-n=k(x-m)\)，与 \(x^{2}/6+y^{2}/3=1\)（即 \(x^{2}+2y^{2}=6\)）联立，消 \(y\) 得 \((1+2k^{2})x^{2}-4k(km-n)x+2(km-n)^{2}-6=0\)．相切 \(\Rightarrow  \Delta =16k^{2}(km-n)^{2}-4(1+2k^{2})[2(km-n)^{2}-6]=0\)， 展开约去同类项后得 \((m^{2}-6)k^{2}-2mnk+n^{2}-3=0\)（\(m^{2}\neq 6\) 时为关于 \(k\) 的二次方程）． 设两切线斜率为 \(k_{1}\)、\(k_{2}\)，则 \(k_{1}k_{2}=(n^{2}-3)/(m^{2}-6)\)，垂直 \(\Rightarrow  (n^{2}-3)/(m^{2}-6)=-1\)，即 \(m^{2}+n^{2}=9\)（\(m\neq \pm \sqrt{6}\)）． 结合（ⅰ）知斜率不存在的情形也在同一圆上，且该圆上每一点到原点的距离 \(3\) 大于椭圆上点的最大距离 \(\sqrt{6}\)，故圆上点均在椭圆外、两切线总存在， 所以动点 \(P\) 的轨迹方程为 \(x^{2}+y^{2}=9\)（即以 \(a^{2}+b^{2}=9\) 为半径平方的蒙日圆）．}
\end{ansblock}

\tuhao{20} \tieside{中}已知 \(A\)，\(B\) 在椭圆 \(C\)：\(x^{2}/a^{2}+y^{2}=1 (a>1)\) 上，动点 \(P\) 在直线 \(3x+4y-10=0\) 上，若 \(\angle APB\) 始终为锐角，则椭圆离心率的取值范围是\kongda{\((0, \sqrt{6}/3)\)}．
\begin{ansblock}[2章-拓-课时12-20]
% ans:2章-拓-课时12-20
\ansitem{20}{\((0, \sqrt{6}/3)\)}
\ansnote{详解}{先定临界位置：对平面内一点 \(P\)（在椭圆外），过 \(P\) 的两条切线的夹角是 \(P\) 对椭圆上两点的张角 \(\angle APB\) 的最大值（\(A\)、\(B\) 沿切线方向时张角最大），而两切线夹角为直角 \(\Leftrightarrow  P\) 在蒙日圆上、为钝角 \(\Leftrightarrow  P\) 在蒙日圆内、为锐角 \(\Leftrightarrow  P\) 在蒙日圆外． 因此\("\)对椭圆上任意 \(A\)、\(B\) 与直线上任意 \(P\)，\(\angle APB\) 恒为锐角\("\Leftrightarrow \) 直线 \(3x+4y-10=0\) 上每一点都在蒙日圆外，即直线与蒙日圆相离． 椭圆 \(x^{2}/a^{2}+y^{2}=1\) 的 \(a^{2}=a^{2}\)、\(b^{2}=1\)，蒙日圆为 \(x^{2}+y^{2}=a^{2}+1\)，圆心 \(O(0,0)\)、半径 \(\sqrt{a^{2}+1}\)． 相离 \(\Leftrightarrow  d=|-10|/\sqrt{9+16}=2>\sqrt{a^{2}+1} \Leftrightarrow  a^{2}+1<4 \Leftrightarrow  a^{2}<3\)，结合 \(a>1\) 得 \(1<a<\sqrt{3}\)，即 \(a^{2}\in (1,3)\)． \(e^{2}=1-b^{2}/a^{2}=1-1/a^{2}\) 在 \(a^{2}\in (1,3)\) 上随 \(a^{2}\) 递增：\(a^{2}\rightarrow 1^{+}\) 时 \(e^{2}\rightarrow 0^{+}\)，\(a^{2}\rightarrow 3^{-}\) 时 \(e^{2}\rightarrow 2/3\)， 故 \(e^{2}\in (0, 2/3)\)，\(e\in (0, \sqrt{6}/3)\)．}
\end{ansblock}

\tuhao{21} \tieside{中}已知点 \(P\) 为直线 \(ax+y-4=0\) 上一点，\(PA\)，\(PB\) 是椭圆 \(C\)：\(x^{2}/a^{2}+y^{2}=1 (a>1)\) 的两条切线，若恰好存在一点 \(P\) 使得 \(PA\perp PB\)，则椭圆 \(C\) 的离心率为\kongda{\(\sqrt{6}/3\)}．
\begin{ansblock}[2章-拓-课时12-21]
% ans:2章-拓-课时12-21
\ansitem{21}{\(\sqrt{6}/3\)}
\ansnote{详解}{设 \(P(m,n)\)，过 \(P\) 的切线 \(y-n=k(x-m)\) 与 \(x^{2}/a^{2}+y^{2}=1\)（即 \(x^{2}+a^{2}y^{2}=a^{2}\)）联立，消 \(y\) 得 \((1+a^{2}k^{2})x^{2}+2a^{2}k(n-km)x+a^{2}(n-km)^{2}-a^{2}=0\)，相切 \(\Rightarrow  \Delta =0\)，化简得 \((a^{2}-m^{2})k^{2}+2mnk+(1-n^{2})=0\)． 当两切线斜率均存在时，\(k_{1}k_{2}=(1-n^{2})/(a^{2}-m^{2})\)，\(PA\perp PB \Leftrightarrow  (1-n^{2})/(a^{2}-m^{2})=-1 \Leftrightarrow  m^{2}+n^{2}=a^{2}+1\)， 即满足 \(PA\perp PB\) 的点落在以 \(O\) 为圆心、\(\sqrt{a^{2}+1}\) 为半径的圆上（蒙日圆）． \("\)恰好存在一点 \(P\) 使 \(PA\perp PB" \Leftrightarrow \) 直线 \(ax+y-4=0\) 与该圆恰有一个公共点，即相切： \(d=|-4|/\sqrt{a^{2}+1}=\sqrt{a^{2}+1} \Rightarrow  a^{2}+1=4 \Rightarrow  a^{2}=3\)（\(a>1\) 故 \(a=\sqrt{3}\)），此时唯一点 \(P\) 为 \(O\) 到直线的垂足 \((\sqrt{3}, 1)\)（\(\sqrt{3}\cdot \sqrt{3}+1=4\) ，\(3+1=4=a^{2}+1\) ）． 该点 \(m^{2}=a^{2}\) 使上面关于 \(k\) 的方程退化，其两切线为 \(x=\sqrt{3}\)（竖直，切于右顶点）与 \(y=1\)（水平，切于上顶点），确实互相垂直，故取等情形成立、结论不受影响． 于是 \(a^{2}=3\)、\(b^{2}=1\)，\(c^{2}=2\)，\(e=c/a=\sqrt{2}/\sqrt{3}=\sqrt{6}/3\)．}
\end{ansblock}

\tuhao{22} \tieside{中}已知椭圆 \(C\)：\(x^{2}/16+y^{2}/9=1\)，直线 \(l\)：\(x+y=6\sqrt{2}\)，\(M\) 是平面上一动点，若由 \(M\) 向椭圆引两条切线 \(l_{1}\)，\(l_{2}\)，且 \(l_{1}\) 与 \(l_{2}\) 的夹角为钝角，则动点 \(M\) 到直线 \(l\) 的距离的取值范围是\kongda{\((1,11)\)}．
\begin{ansblock}[2章-拓-课时12-22]
% ans:2章-拓-课时12-22
\ansitem{22}{\((1,11)\)}
\ansnote{详解}{两切线夹角为钝角 \(\Leftrightarrow  M\) 在蒙日圆内部（不含边界）；又要能引出两条切线，须 \(M\) 在椭圆外． 蒙日圆半径平方 \(r^{2}=a^{2}+b^{2}=16+9=25\)，即 \(M\) 在圆 \(x^{2}+y^{2}=25\) 内、椭圆外（开区域，两端边界均不含）． 圆心 \(O\) 到 \(l\)：\(x+y-6\sqrt{2}=0\) 的距离 \(d_{0}=6\sqrt{2}/\sqrt{2}=6>5\)，故直线 \(l\) 在蒙日圆外，圆上点到 \(l\) 的距离取值恰为 \((6-5, 6+5)=(1,11)\) 的上下确界． 对任一 \(v\in (1,11)\)，与 \(l\) 平行且相距 \(v\) 的直线与开圆盘相交成一条弦，弦的端点趋近圆上点；因椭圆上点到原点距离最大为 \(a=4<5\)，圆附近、弦上的点必在椭圆外，故该直线上的可取点确实存在，\(v\) 可达； 反过来，区域内的点到 \(l\) 的距离不可能达到 \(1\) 或 \(11\)（那需在圆外或圆上）． 故动点 \(M\) 到直线 \(l\) 的距离的取值范围是 \((1,11)\)（开区间，不含端点）。}
\end{ansblock}

\tuhao{23} \tieside{难}已知椭圆 \(C\)：\(x^{2}/2+y^{2}=1\)，\(M\) 是圆 \(x^{2}+y^{2}=3\) 上的任意一点，\(MA\)，\(MB\) 分别与椭圆切于 \(A\)，\(B\)．求 \(\triangle AOB\) 面积的取值范围．
\liubai[16mm]{解答书写区}
\begin{ansblock}[2章-拓-课时12-23]
% ans:2章-拓-课时12-23
\ansitem{23}{\([2/3, \sqrt{2}/2]\)}
\ansnote{详解}{设 \(M(x_{0},y_{0})\)（\(x_{0}^{2}+y_{0}^{2}=3\)）、\(A(x_{1},y_{1})\)、\(B(x_{2},y_{2})\)． 由切线代换式（见 \(12-\)拓\(18\)）：过椭圆 \(x^{2}/2+y^{2}=1\) 上点 \((x_{1},y_{1})\) 的切线为 \(x_{1}x/2+y_{1}y=1\)，切线 \(MA\) 过 \(M \Rightarrow  x_{1}x_{0}/2+y_{1}y_{0}=1\)，同理 \(x_{2}x_{0}/2+y_{2}y_{0}=1\)， 故切点弦 \(AB\)：\(x_{0}x/2+y_{0}y=1\)． （ⅰ）\(y_{0}\neq 0\) 时，由直线解出 \(y=(1-x_{0}x/2)/y_{0}\) 代回 \(x^{2}+2y^{2}=2\)，整理并用 \(x_{0}^{2}=3-y_{0}^{2}\) 化简得 \((y_{0}^{2}+3)x^{2}-4x_{0}x+(4-4y_{0}^{2})=0\)，\(\therefore  x_{1}+x_{2}=4x_{0}/(y_{0}^{2}+3)\)，\(x_{1}x_{2}=(4-4y_{0}^{2})/(y_{0}^{2}+3)\)， \((x_{1}-x_{2})^{2}=(x_{1}+x_{2})^{2}-4x_{1}x_{2}=[16x_{0}^{2}-16(1-y_{0}^{2})(y_{0}^{2}+3)]/(y_{0}^{2}+3)^{2}=16y_{0}^{2}(y_{0}^{2}+1)/(y_{0}^{2}+3)^{2}\)， 又 \(y_{1}-y_{2}=-[x_{0}/(2y_{0})](x_{1}-x_{2})\)，故 \(|AB|=|x_{1}-x_{2}|\cdot \sqrt{}(1+x_{0}^{2}/(4y_{0}^{2}))=[4|y_{0}|\sqrt{y_{0}^{2}+1}/(y_{0}^{2}+3)]\cdot [\sqrt{4y_{0}^{2}+x_{0}^{2}}/(2|y_{0}|)]\)， 而 \(4y_{0}^{2}+x_{0}^{2}=4y_{0}^{2}+3-y_{0}^{2}=3(y_{0}^{2}+1)\)，所以 \(|AB|=2\sqrt{3}(y_{0}^{2}+1)/(y_{0}^{2}+3)\)． 原点 \(O\) 到 \(AB\) 的距离 \(d=1/\sqrt{x_{0}^{2}/4+y_{0}^{2}}=1/\sqrt{}(3(1+y_{0}^{2})/4)=2/[\sqrt{3}\cdot \sqrt{y_{0}^{2}+1}]\)， \(S=(1/2)|AB|\cdot d=(1/2)\cdot [2\sqrt{3}(y_{0}^{2}+1)/(y_{0}^{2}+3)]\cdot [2/(\sqrt{3}\sqrt{y_{0}^{2}+1})]=2\sqrt{y_{0}^{2}+1}/(y_{0}^{2}+3)\)． （ⅱ）\(y_{0}=0\) 时 \(x_{0}^{2}=3\)，\(AB\) 为 \(x=2/x_{0}=\pm 2\sqrt{3}/3\)，切点纵坐标满足 \(y^{2}=1-x^{2}/2=1-(2/\sqrt{3})^{2}/2=1-2/3=1/3\)， \(|AB|=2\sqrt{3}/3\)，\(d=2\sqrt{3}/3\)，\(S=(1/2)\cdot (2\sqrt{3}/3)\cdot (2\sqrt{3}/3)=2/3\)，与 \((\)ⅰ\()\) 式在 \(y_{0}=0\) 时的值一致，故通式覆盖全部情形． 令 \(t=\sqrt{y_{0}^{2}+1}\)，由 \(y_{0}^{2}\in [0,3]\) 得 \(t\in [1,2]\)，则 \(y_{0}^{2}+3=t^{2}+2\)，\(S=2t/(t^{2}+2)=2/(t+2/t)\)． \(t+2/t\) 在 \([1,2]\) 上：由对勾函数性质，在 \(t=\sqrt{2}\)（\(\in [1,2]\)）处取最小 \(2\sqrt{2}\)，在端点 \(t=1\)、\(t=2\) 处取最大 \(3\)， 故 \(t+2/t\in [2\sqrt{2}, 3]\)，\(S\in [2/3, 2/(2\sqrt{2})]=[2/3, \sqrt{2}/2]\)． 所以 \(\triangle AOB\) 面积的取值范围为 \([2/3, \sqrt{2}/2]\)．}
\end{ansblock}

\tuhao{24} \tieside{中}已知从圆 \(C\)：\(x^{2}+y^{2}=r^{2} (r>0)\) 上一点 \(Q(0,r)\) 作两条互相垂直的直线与椭圆 \(\tau \)：\(x^{2}/12+y^{2}/4=1\) 相切，同时圆 \(C\) 与直线 \(l\)：\(mx+y-\sqrt{3}m-1=0\) 交于 \(A\)，\(B\) 两点，则 \(|AB|\) 的最小值为（　　）．
\bindopt {\fontsize{10.5pt}{\qpTJgLeadLiti}\selectfont \makebox[\dimexpr0.25\linewidth-0.25\lxhang-0.25em\relax][l]{\(A\)．\(2\sqrt{3}\)}\makebox[\dimexpr0.25\linewidth-0.25\lxhang-0.25em\relax][l]{\(B\)．\(4\)}\makebox[\dimexpr0.25\linewidth-0.25\lxhang-0.25em\relax][l]{\(C\)．\(4\sqrt{3}\)}\makebox[\dimexpr0.25\linewidth-0.25\lxhang-0.25em\relax][l]{\(D\)．\(8\)}\par}
\begin{ansblock}[2章-拓-课时12-24]
% ans:2章-拓-课时12-24
\ansitem{24}{\(C\)}
\ansnote{详解}{设过 \(Q(0,r)\) 的切线为 \(y=kx+r\)，与 \(x^{2}/12+y^{2}/4=1\)（即 \(x^{2}+3y^{2}=12\)）联立：\((1+3k^{2})x^{2}+6krx+3r^{2}-12=0\)， \(\Delta =36k^{2}r^{2}-4(1+3k^{2})(3r^{2}-12)=-12r^{2}+48+144k^{2}=0 \Rightarrow  k^{2}=(r^{2}-4)/12\)． 两条切线的斜率为该式的两根 \(k=\pm \sqrt{}((r^{2}-4)/12)\)，互相垂直 \(\Rightarrow \) 其积 \(-(r^{2}-4)/12=-1 \Rightarrow  r^{2}=16\)，圆 \(C\)：\(x^{2}+y^{2}=16\)（半径 \(4\)）． 直线 \(l\)：\(m(x-\sqrt{3})+(y-1)=0\) 恒过定点 \(P(\sqrt{3},1)\)，且 \(|OP|=\sqrt{3+1}=2<4\)，\(P\) 在圆内． 弦长 \(|AB|=2\sqrt{16-d^{2}}\)（\(d\) 为圆心 \(O\) 到 \(l\) 的距离），要使 \(|AB|\) 最小须 \(d\) 最大；\(l\) 恒过定点 \(P\)，故 \(d\leqslant |OP|=2\)，等号当 \(l\perp OP\) 时取得（此时 \(m\) 使 \(l\) 的方向与 \(OP\) 垂直，可取）． 所以 \(|AB|min=2\sqrt{16-4}=2\sqrt{12}=4\sqrt{3}\)．故选：\(C\)．}
\end{ansblock}

\tuhao{25} \tieside{难}\("\)蒙日圆\("\)涉及几何学中的一个著名定理，该定理的内容为：椭圆上任意两条互相垂直的切线的交点，必在一个与椭圆同心的圆上．称此圆为该椭圆的\("\)蒙日圆\("\)，该圆由法国数学家加斯帕尔\(\cdot \)蒙日（\(1746\)—\(1818\)）最先发现．若椭圆 \(C\)：\(x^{2}/4+y^{2}=1\) 的左、右焦点分别为 \(F_{1}\)、\(F_{2}\)，\(P\) 为椭圆 \(C\) 上一动点，过 \(P\) 和原点作直线 \(l\) 与椭圆 \(C\) 的蒙日圆相交于 \(M\)，\(N\)，则 \((|PM|\cdot |PN|)/(|PF_{1}|\cdot |PF_{2}|)=\)\kongda{\(1\)}．
\begin{ansblock}[2章-拓-课时12-25]
% ans:2章-拓-课时12-25
\ansitem{25}{\(1\)}
\ansnote{详解}{\(a^{2}=4\)、\(b^{2}=1\)、\(c^{2}=3\)；蒙日圆半径平方 \(r^{2}=a^{2}+b^{2}=5\)（过四顶点作切线所得矩形的外接圆，对角线一半即 \(\sqrt{5}\)），圆方程 \(x^{2}+y^{2}=5\)，记圆心 \(O\)、半径 \(r=\sqrt{5}\)． 设 \(m=|PF_{1}|\cdot |PF_{2}|\)．由 \(|PF_{1}|+|PF_{2}|=2a=4\) 得 \(|PF_{1}|^{2}+|PF_{2}|^{2}=16-2m\)； 在 \(\triangle PF_{1}F_{2}\) 中 \(O\) 为 \(F_{1}F_{2}\) 的中点，由中线（阿波罗尼奥斯）定理 \(|PF_{1}|^{2}+|PF_{2}|^{2}=2(|PO|^{2}+|OF_{1}|^{2})=2|PO|^{2}+2c^{2}=2|PO|^{2}+6\)， 故 \(16-2m=2|PO|^{2}+6 \Rightarrow  |PO|^{2}=5-m\)． 又 \(|PF_{1}|\cdot |PF_{2}|=a^{2}-e^{2}x\_P^{2}=4-(3/4)x\_P^{2}\)，\(x\_P^{2}\in [0,4]\)，故 \(m\in [1,4]\)，从而 \(|PO|^{2}=5-m\in [1,4]<r^{2}=5\)，即 \(P\) 在蒙日圆内（与图相符，直线 \(l\) 与圆总交于两点）． 过圆内一点 \(P\) 的弦 \(MN\) 被 \(P\) 分成两段，由相交弦定理（或直接 \((r-|PO|)(r+|PO|)\)）得 \(|PM|\cdot |PN|=r^{2}-|PO|^{2}=5-(5-m)=m\)， 所以 \((|PM|\cdot |PN|)/(|PF_{1}|\cdot |PF_{2}|)=m/m=1\)．}
\end{ansblock}

\tuhao{26} \tieside{难}已知圆 \(O\)：\(x^{2}+y^{2}=34\)，椭圆 \(C\)：\(x^{2}/25+y^{2}/9=1\)．
\((1)\) 若点 \(P\) 在圆 \(O\) 上，线段 \(OP\) 的垂直平分线经过椭圆的右焦点，求点 \(P\) 的横坐标；
\((2)\) 现有如下真命题：①过圆 \(x^{2}+y^{2}=5^{2}+3^{2}\) 上任意一点 \(Q(m,n)\) 作椭圆 \(x^{2}/5^{2}+y^{2}/3^{2}=1\) 的两条切线，则这两条切线互相垂直；②过圆 \(x^{2}+y^{2}=4^{2}+7^{2}\) 上任意一点 \(Q(m,n)\) 作椭圆 \(x^{2}/4^{2}+y^{2}/7^{2}=1\) 的两条切线，则这两条切线互相垂直．据此写出一般结论，并加以证明．
\liubai[24mm]{解答书写区}
\begin{ansblock}[2章-拓-课时12-26]
% ans:2章-拓-课时12-26
\ansitem{26}{\((1) 17/4\)；\((2)\) 一般结论：过圆 \(x^{2}+y^{2}=a^{2}+b^{2}\) 上任意一点 \(Q(m,n)\) 作椭圆 \(x^{2}/a^{2}+y^{2}/b^{2}=1 (a>b>0)\) 的两条切线，则这两条切线互相垂直；证明见详解}
\ansnote{详解}{\((1)\) 设 \(P(x_{0},y_{0})\)，则 \(x_{0}^{2}+y_{0}^{2}=34\) ①．椭圆 \(a^{2}=25\)、\(b^{2}=9 \Rightarrow  c^{2}=16\)，右焦点 \(F(4,0)\)． \(F\) 在线段 \(OP\) 的垂直平分线上 \(\Leftrightarrow  |FP|=|FO|=4\)，即 \((x_{0}-4)^{2}+y_{0}^{2}=16\) ②， ②展开为 \(x_{0}^{2}+y_{0}^{2}-8x_{0}+16=16\)，即 \(x_{0}^{2}+y_{0}^{2}-8x_{0}=0\)，代①：\(34-8x_{0}=0 \Rightarrow  x_{0}=34/8=17/4\)． （存在性核验：\(x_{0}=17/4<\sqrt{34}\)，\(y_{0}^{2}=34-289/16=255/16>0\) ）故点 \(P\) 的横坐标为 \(17/4\)． \((2)\) 一般结论：过圆 \(x^{2}+y^{2}=a^{2}+b^{2}\) 上任意一点 \(Q(m,n)\) 作椭圆 \(x^{2}/a^{2}+y^{2}/b^{2}=1 (a>b>0)\) 的两条切线，则这两条切线互相垂直（即题中①②所述\("\)圆 \(x^{2}+y^{2}=a^{2}+b^{2}"\)正是该椭圆的蒙日圆，①中 \(5^{2}+3^{2}=34\) 与 \((1)\) 的圆 \(O\) 同）。 证明：设 \(Q(m,n)\) 满足 \(m^{2}+n^{2}=a^{2}+b^{2}\)． （ⅰ）若过 \(Q\) 的切线中有一条斜率不存在，则该切线为 \(x=\pm a\)，于是 \(m=\pm a\)，由 \(m^{2}+n^{2}=a^{2}+b^{2}\) 得 \(n^{2}=b^{2}\)，即 \(n=\pm b\)，\(Q(\pm a,\pm b)\)； 此时直线 \(y=\pm b\) 也过 \(Q\) 且与椭圆相切（切于上、下顶点），它与 \(x=\pm a\) 垂直，故两切线互相垂直（再过 \(Q\) 不能有其他切线，因椭圆的切线必在 \(x=\pm a\)、\(y=\pm b\) 围成的矩形外或为其边）． （ⅱ）若过 \(Q\) 的两条切线斜率均存在，设切线为 \(y-n=k(x-m)\)，与 \(b^{2}x^{2}+a^{2}y^{2}=a^{2}b^{2}\) 联立，消 \(y\) 得 \((b^{2}+a^{2}k^{2})x^{2}+2a^{2}k(n-km)x+a^{2}(n-km)^{2}-a^{2}b^{2}=0\)， 相切 \(\Rightarrow  \Delta =4a^{4}k^{2}(n-km)^{2}-4(b^{2}+a^{2}k^{2})[a^{2}(n-km)^{2}-a^{2}b^{2}]=0\)，两边除 \(4a^{2}\) 并整理得 \((m^{2}-a^{2})k^{2}-2mnk+(n^{2}-b^{2})=0\)（ⅱ 中两根即两切线斜率；若 \(m^{2}=a^{2}\) 则由 \(a^{2}+b^{2}=m^{2}+n^{2}\) 得 \(n^{2}=b^{2}\)，即情形（ⅰ））． 故 \(k_{1}k_{2}=(n^{2}-b^{2})/(m^{2}-a^{2})\)；又 \(m^{2}+n^{2}=a^{2}+b^{2} \Rightarrow  m^{2}-a^{2}=b^{2}-n^{2}=-(n^{2}-b^{2})\)， 所以 \(k_{1}k_{2}=(n^{2}-b^{2})/[-(n^{2}-b^{2})]=-1\)，两切线互相垂直． 综合（ⅰ）（ⅱ），一般结论成立．}
\end{ansblock}

\zux{蒙日圆族}
\tuhao{28} \tieside{简}已知 \(F_{1}\)，\(F_{2}\) 分别是椭圆 \(x^{2}/a^{2}+y^{2}/b^{2}=1 (a>b>0)\) 的左、右焦点，若椭圆上存在点 \(P\)，使 \(\angle F_{1}PF_{2}=90^{\circ}\)，则椭圆的离心率 \(e\) 的取值范围为（　　）
\bindopt {\fontsize{10.5pt}{\qpTJgLeadLiti}\selectfont \makebox[\dimexpr0.5\linewidth-0.5\lxhang-0.25em\relax][l]{\(A\)．\((0, \sqrt{2}/2]\)}\makebox[\dimexpr0.5\linewidth-0.5\lxhang-0.25em\relax][l]{\(B\)．\([\sqrt{2}/2, 1)\)}\par}
\bindopt {\fontsize{10.5pt}{\qpTJgLeadLiti}\selectfont \makebox[\dimexpr0.5\linewidth-0.5\lxhang-0.25em\relax][l]{\(C\)．\((0, \sqrt{3}/2]\)}\makebox[\dimexpr0.5\linewidth-0.5\lxhang-0.25em\relax][l]{\(D\)．\([\sqrt{5}/2, 1)\)}\par}
\begin{ansblock}[2章-拓-课时12-28]
% ans:2章-拓-课时12-28
\ansitem{28}{\(B\)}
\ansnote{详解}{\(\angle F_{1}PF_{2}=90^{\circ} \Leftrightarrow  P\) 在以 \(F_{1}F_{2}\) 为直径、即以原点 \(O\) 为圆心、\(c\) 为半径的圆上，故题设等价于该圆与椭圆有公共点． 椭圆上点到中心 \(O\) 的距离：\(|OP|^{2}=x^{2}+y^{2}=a^{2}-(a^{2}-b^{2})x^{2}/a^{2}\)（用 \(y^{2}=b^{2}(1-x^{2}/a^{2})\)），\(x^{2}\in [0,a^{2}]\) 时随 \(x^{2}\) 递减，故 \(|OP|^{2}\in [b^{2},a^{2}]\)，即 \(|OP|\in [b,a]\)． 于是圆 \(x^{2}+y^{2}=c^{2}\) 与椭圆有交点 \(\Leftrightarrow  b\leqslant c\leqslant a\)；而 \(c<a\) 恒成立（\(c^{2}=a^{2}-b^{2}<a^{2}\)），故条件即 \(b\leqslant c\)． \(b^{2}\leqslant c^{2} \Rightarrow  a^{2}-c^{2}\leqslant c^{2} \Rightarrow  a^{2}\leqslant 2c^{2} \Rightarrow  e^{2}\geqslant 1/2\)，结合 \(0<e<1\) 得 \(e\in [\sqrt{2}/2, 1)\)．故选：\(B\)．}
\end{ansblock}

\tuhao{29} \tieside{简}离心率相同的椭圆是同一个椭圆．（　　）（判断正误）
\liubai[16mm]{解答书写区}
\begin{ansblock}[2章-拓-课时12-29]
% ans:2章-拓-课时12-29
\ansitem{29}{错误}
\ansnote{详解}{椭圆的离心率 \(e=\sqrt{1-b^{2}/a^{2}}\) 只含 \(a\) 与 \(b\) 的比值这一个信息，而确定一个椭圆需要 \(a\)、\(b\) 两个量（还需定位焦点在哪个轴上），故 \(e\) 相同只保证形状相似，不保证大小与方位相同． 反例：\(x^{2}/4+y^{2}/3=1\) 与 \(x^{2}/8+y^{2}/6=1\)，两者的 \(b^{2}/a^{2}\) 同为 \(3/4\)，\(e\) 均为 \(1/2\)，但长轴长分别为 \(4\) 与 \(4\sqrt{2}\)，显然不是同一个椭圆；再把第二个改写为 \(y^{2}/8+x^{2}/6=1\)，则离心率仍为 \(1/2\) 而焦点改在 \(y\) 轴上． 故\("\)离心率相同的椭圆是同一个椭圆\("\)错误（正确说法是：离心率相同的椭圆形状相同\(/\)相似，长短轴之比相同）．}
\end{ansblock}

\tuhao{30} \tieside{中}已知点 \(M\) 为椭圆 \(x^{2}/a^{2}+y^{2}/b^{2}=1 (a>b>0)\) 上一点，椭圆的长轴长为 \(8\sqrt{2}\)，离心率 \(e=\sqrt{2}/2\)，左、右焦点分别为 \(F_{1}\)，\(F_{2}\)，其中 \(B(3,2)\)，则 \(|MF_{1}|+|MB|\) 的最小值为（　　）
\bindopt {\fontsize{10.5pt}{\qpTJgLeadLiti}\selectfont \makebox[\dimexpr0.5\linewidth-0.5\lxhang-0.25em\relax][l]{\(A\)．\(5\sqrt{2}\)}\makebox[\dimexpr0.5\linewidth-0.5\lxhang-0.25em\relax][l]{\(B\)．\(6\sqrt{2}-\sqrt{3}\)}\par}
\bindopt {\fontsize{10.5pt}{\qpTJgLeadLiti}\selectfont \makebox[\dimexpr0.5\linewidth-0.5\lxhang-0.25em\relax][l]{\(C\)．\(4\sqrt{2}\)}\makebox[\dimexpr0.5\linewidth-0.5\lxhang-0.25em\relax][l]{\(D\)．\(8\sqrt{2}-\sqrt{5}\)}\par}
\begin{ansblock}[2章-拓-课时12-30]
% ans:2章-拓-课时12-30
\ansitem{30}{\(D\)}
\ansnote{详解}{\(2a=8\sqrt{2} \Rightarrow  a=4\sqrt{2}\)、\(a^{2}=32\)；\(e=\sqrt{2}/2 \Rightarrow  c=a\cdot e=4\)、\(c^{2}=16\)，\(b^{2}=32-16=16\)，椭圆为 \(x^{2}/32+y^{2}/16=1\)，\(F_{2}(4,0)\)． 由定义 \(|MF_{1}|=2a-|MF_{2}|\)，故 \(|MF_{1}|+|MB|=2a+(|MB|-|MF_{2}|)=8\sqrt{2}+(|MB|-|MF_{2}|)\)． 对任意 \(M\)，由三角形两边之差：\(|MB|-|MF_{2}|\geqslant -|MF_{2}|+|MB|\geqslant -|BF_{2}|\)（即 \(||MB|-|MF_{2}||\leqslant |BF_{2}|\)）， 而 \(|BF_{2}|=\sqrt{}((3-4)^{2}+(2-0)^{2})=\sqrt{5}\)，故 \(|MF_{1}|+|MB|\geqslant 8\sqrt{2}-\sqrt{5}\)． 等号条件：\(|MF_{2}|-|MB|=|BF_{2}|\)，即 \(B\) 在线段 \(MF_{2}\) 上，亦即 \(M\) 在射线 \(F_{2}B\)（自 \(F_{2}\) 过 \(B\) 向外）上；因 \(B(3,2)\) 满足 \(9/32+4/16=17/32<1\) 在椭圆内，该射线必与椭圆交于一点 \(M\)，故等号可取到． 所以最小值为 \(8\sqrt{2}-\sqrt{5}\)．故选：\(D\)．}
\end{ansblock}

\tuhao{31} \tieside{简}定义：若点 \(P(x_{0},y_{0})\) 在椭圆 \(x^{2}/a^{2}+y^{2}/b^{2}=1 (a>b>0)\) 上，则以 \(P\) 为切点的切线方程为 \(x_{0}x/a^{2}+y_{0}y/b^{2}=1\)．已知椭圆 \(C\)：\(x^{2}/3+y^{2}/2=1\)，点 \(M\) 为直线 \(x-2y-6=0\) 上一个动点，过点 \(M\) 作椭圆 \(C\) 的两条切线 \(MA\)，\(MB\)，切点分别为 \(A\)，\(B\)，则直线 \(AB\) 恒过定点（　　）
\bindopt {\fontsize{10.5pt}{\qpTJgLeadLiti}\selectfont \makebox[\dimexpr0.5\linewidth-0.5\lxhang-0.25em\relax][l]{\(A\)．\((1/2, -1/3)\)}\makebox[\dimexpr0.5\linewidth-0.5\lxhang-0.25em\relax][l]{\(B\)．\((-1/2, 1/3)\)}\par}
\bindopt {\fontsize{10.5pt}{\qpTJgLeadLiti}\selectfont \makebox[\dimexpr0.5\linewidth-0.5\lxhang-0.25em\relax][l]{\(C\)．\((1/2, -2/3)\)}\makebox[\dimexpr0.5\linewidth-0.5\lxhang-0.25em\relax][l]{\(D\)．\((-1/2, 2/3)\)}\par}
\begin{ansblock}[2章-拓-课时12-31]
% ans:2章-拓-课时12-31
\ansitem{31}{\(C\)}
\ansnote{详解}{设 \(M(2t+6, t)\)（在直线上，\(x_{0}=2y_{0}+6\)），\(A(x_{1},y_{1})\)、\(B(x_{2},y_{2})\)． 由题给切线公式，切线 \(MA\)：\(x_{1}x/3+y_{1}y/2=1\)，它过 \(M \Rightarrow  x_{1}(2t+6)/3+y_{1}t/2=1\)；同理 \(x_{2}(2t+6)/3+y_{2}t/2=1\)． 故 \(A\)、\(B\) 的坐标同满足方程 \((2t+6)x/3+ty/2=1\)，此即直线 \(AB\) 的方程． 两边乘 \(6\)：\(2(2t+6)x+3ty=6 \Rightarrow  (4x+3y)t+(12x-6)=0\) 对一切 \(t\) 成立 \(\Rightarrow  4x+3y=0\) 且 \(12x-6=0\)， 解得 \(x=1/2\)、\(y=-2/3\)（满足 \(4\cdot (1/2)+3\cdot (-2/3)=2-2=0\) ）． （可行性核验：直线上点 \(M(2t+6,t)\) 到椭圆的定位 \((2t+6)^{2}/3+t^{2}/2=(11/6)t^{2}+8t+12\)，其最小值在 \(t=-24/11\) 处为 \(36/11>1\)，故 \(M\) 恒在椭圆外，两切线总存在．） 所以直线 \(AB\) 恒过定点 \((1/2, -2/3)\)．故选：\(C\)．}
\end{ansblock}

\tuhao{32} \tieside{中}已知椭圆 \(E\) 的中心在原点，焦点在 \(x\) 轴上，长轴长为 \(4\)，离心率等于 \(\sqrt{3}/2\)．
\((1)\) 求椭圆 \(E\) 的方程；
\((2)\) 设 \(N(0,1)\)，若椭圆 \(E\) 上存在两个不同点 \(P\)、\(Q\) 满足 \(\angle PNQ=90^{\circ}\)，证明：直线 \(PQ\) 过定点，并求该定点的坐标．
\liubai[24mm]{解答书写区}
\begin{ansblock}[2章-拓-课时12-32]
% ans:2章-拓-课时12-32
\ansitem{32}{\((1) x^{2}/4+y^{2}=1\)；\((2)\) 证明见详解，定点为 \((0, -3/5)\)}
\ansnote{详解}{\((1) 2a=4 \Rightarrow  a=2\)；\(e=c/a=\sqrt{3}/2 \Rightarrow  c=\sqrt{3}\)，\(b^{2}=a^{2}-c^{2}=1\)，椭圆 \(E\)：\(x^{2}/4+y^{2}=1\)． \((2)\) 先设直线 \(PQ\) 斜率存在，方程 \(y=kx+m\)（\(m\neq 1\)，否则直线过 \(N(0,1)\)，\(P\)、\(Q\) 中有一点与 \(N\) 重合，\(\angle PNQ\) 无意义）． 与 \(x^{2}+4y^{2}=4\) 联立：\((1+4k^{2})x^{2}+8kmx+4(m^{2}-1)=0\)，须 \(\Delta =64k^{2}m^{2}-16(1+4k^{2})(m^{2}-1)=16[(1+4k^{2})-m^{2}]>0\)， 且 \(x_{1}+x_{2}=-8km/(1+4k^{2})\)，\(x_{1}x_{2}=4(m^{2}-1)/(1+4k^{2})\)． \(\angle PNQ=90^{\circ} \Leftrightarrow \) 向量\(NP\cdot NQ=0\)，其中 \(NP=(x_{1}, y_{1}-1)\)、\(NQ=(x_{2}, y_{2}-1)\)： \(x_{1}x_{2}+(kx_{1}+m-1)(kx_{2}+m-1)=(1+k^{2})x_{1}x_{2}+k(m-1)(x_{1}+x_{2})+(m-1)^{2}=0\)， 代韦达并乘 \((1+4k^{2})\)：\((1+k^{2})\cdot 4(m^{2}-1)+k(m-1)(-8km)+(m-1)^{2}(1+4k^{2})=0\)， 提取 \((m-1)\)（\(m\neq 1\)）：\(4(1+k^{2})(m+1)-8k^{2}m+(m-1)(1+4k^{2})=0\)， 展开：\(4+4k^{2}+4m+4k^{2}m-8k^{2}m+m+4k^{2}m-1-4k^{2}=0 \Rightarrow  3+5m=0 \Rightarrow  m=-3/5\)（满足 \(\Delta >0\)：\(1+4k^{2}>9/25\) 恒成立）． 故 \(PQ\)：\(y=kx-3/5\) 过定点 \((0, -3/5)\)． 再设直线 \(PQ\) 斜率不存在（\(x=s\)）：\(P(s,u)\)、\(Q(s,-u)\)（\(u^{2}=1-s^{2}/4\)），则 \(NP\cdot NQ=s^{2}+(u-1)(-u-1)=s^{2}+1-u^{2}=s^{2}+s^{2}/4=(5/4)s^{2}\)， 要为 \(0\) 须 \(s=0\)，此时 \(P(0,\pm 1)\)，其中 \((0,1)\) 即 \(N\)，\(\angle PNQ\) 不存在，故此情形不合题意． 综上，直线 \(PQ\) 恒过定点 \((0, -3/5)\)．}
\end{ansblock}

\tuhao{33} \tieside{中}已知椭圆 \(E\)：\(x^{2}/a^{2}+y^{2}/b^{2}=1 (a>b>0)\) 的离心率为 \(\sqrt{3}/2\)，且过点 \((1, \sqrt{3}/2)\)，\(A\)，\(B\) 分别为椭圆 \(E\) 的左、右顶点，\(P\) 为直线 \(x=3\) 上的动点（不在 \(x\) 轴上），\(PA\) 与椭圆 \(E\) 的另一交点为 \(C\)，\(PB\) 与椭圆 \(E\) 的另一交点为 \(D\)，记直线 \(PA\) 与 \(PB\) 的斜率分别为 \(k_{1}\)，\(k_{2}\)．
\((\)Ⅰ\()\) 求椭圆 \(E\) 的方程；
\((\)Ⅱ\()\) 求 \(k_{1}/k_{2}\) 的值；
\((\)Ⅲ\()\) 证明：直线 \(CD\) 过一个定点，并求出此定点的坐标．
\liubai[16mm]{解答书写区}
\begin{ansblock}[2章-拓-课时12-33]
% ans:2章-拓-课时12-33
\ansitem{33}{\((\)Ⅰ\() x^{2}/4+y^{2}=1\)；\((\)Ⅱ\() 1/5\)；\((\)Ⅲ\()\) 证明见详解，定点为 \((4/3, 0)\)}
\ansnote{详解}{\((\)Ⅰ\() e^{2}=c^{2}/a^{2}=3/4 \Rightarrow  a^{2}=4b^{2}\)（由 \(c^{2}=a^{2}-b^{2}\)）；点 \((1, \sqrt{3}/2)\) 代入：\(1/a^{2}+3/(4b^{2})=1\)，即 \(1/(4b^{2})+3/(4b^{2})=1/b^{2}=1 \Rightarrow  b^{2}=1\)、\(a^{2}=4\)，\(E\)：\(x^{2}/4+y^{2}=1\)． \((\)Ⅱ\() A(-2,0)\)、\(B(2,0)\)，设 \(P(3,t)\)（\(t\neq 0\)），\(k_{1}=t/(3+2)=t/5\)、\(k_{2}=t/(3-2)=t\)，故 \(k_{1}/k_{2}=1/5\)． \((\)Ⅲ\() PA\)：\(y=(t/5)(x+2)\)，与 \(x^{2}+4y^{2}=4\) 联立得 \((4t^{2}+25)x^{2}+16t^{2}x+16t^{2}-100=0\)． 其一根为 \(x\_A=-2\)，由韦达 \(x\_A+x\_C=-16t^{2}/(4t^{2}+25) \Rightarrow  x\_C=-16t^{2}/(4t^{2}+25)+2=(50-8t^{2})/(4t^{2}+25)\)，\(y\_C=(t/5)(x\_C+2)=20t/(4t^{2}+25)\)． \(PB\)：\(y=t(x-2)\)，联立得 \((1+4t^{2})x^{2}-16t^{2}x+16t^{2}-4=0\)，一根为 \(x\_B=2\)，\(x\_B+x\_D=16t^{2}/(1+4t^{2}) \Rightarrow  x\_D=(8t^{2}-2)/(1+4t^{2})\)，\(y\_D=t(x\_D-2)=-4t/(1+4t^{2})\)． 直线 \(CD\) 令 \(y=0\) 求横截距：\(x_{0}=x\_C-y\_C(x\_D-x\_C)/(y\_D-y\_C)\)，代上述四式化简得 \(x_{0}=4/3\)（与 \(t\) 无关）； （独立核验：直接取 \(t=1\)、\(t=2\) 分别算得 \(C\)、\(D\) 坐标，两点连线与 \(x\) 轴交点均为 \((4/3,0)\) ） 故直线 \(CD\) 恒过定点 \((4/3, 0)\)．}
\end{ansblock}

% 图债注：「如图」指源件配图，印面无图（冻片【注】裁定：去装饰图/条件全在文字/尺寸随注——图债登记汇编报告）
\tuhao{34} \tieside{中}如图，已知椭圆 \(\Gamma \)：\(x^{2}/b^{2}+y^{2}/a^{2}=1 (a>b>0)\) 的离心率 \(e=\sqrt{2}/2\)，短轴右端点为 \(A\)，\(M(1,0)\) 为线段 \(OA\) 的中点．
\((1)\) 求椭圆 \(\Gamma \) 的方程；
\((2)\) 过点 \(M\) 任作一条直线与椭圆 \(\Gamma \) 相交于两点 \(P\)，\(Q\)，试问在 \(x\) 轴上是否存在定点 \(N\)，使得 \(\angle PNM=\angle QNM\)，若存在，求出点 \(N\) 的坐标；若不存在，说明理由．
\liubai[24mm]{解答书写区}
\begin{ansblock}[2章-拓-课时12-34]
% ans:2章-拓-课时12-34
\ansitem{34}{\((1) x^{2}/4+y^{2}/8=1\)；\((2)\) 存在，\(N(4,0)\)}
\ansnote{详解}{\((1) a^{2}\) 在 \(y\) 项下，故长轴在 \(y\) 轴、短轴在 \(x\) 轴，短轴右端点 \(A(b,0)\)．\(M(1,0)\) 为 \(OA\) 中点 \(\Rightarrow  b=2\)． \(e^{2}=c^{2}/a^{2}=1/2\) 且 \(c^{2}=a^{2}-b^{2} \Rightarrow  a^{2}-4=a^{2}/2 \Rightarrow  a^{2}=8\)，故 \(\Gamma \)：\(x^{2}/4+y^{2}/8=1\)． \((2)\) 因 \(M\)、\(N\) 均在 \(x\) 轴上，\(\angle PNM=\angle QNM\) 表示射线 \(NP\)、\(NQ\) 关于 \(x\) 轴对称，即 \(k_{PN}+k_{QN}=0\)（\(k_{PN}=k_{QN}\) 会使 \(P\)、\(Q\)、\(N\) 共线，此时两角互补不相等，舍）． 设 \(N(x_{0},0)\)．当 \(PQ\perp x\) 轴时，\(PQ\)：\(x=1\)，\(P\)、\(Q\) 关于 \(x\) 轴对称，对任意 \(x_{0}\) 恒有 \(\angle PNM=\angle QNM\)，不产生约束． 当 \(PQ\) 不垂直 \(x\) 轴时，设 \(PQ\)：\(y=k(x-1)\)，代入 \(2x^{2}+y^{2}=8\) 得 \((k^{2}+2)x^{2}-2k^{2}x+k^{2}-8=0\)（\(\Delta =4k^{4}-4(k^{2}+2)(k^{2}-8)=24k^{2}+64>0\) 恒成立）， \(x_{1}+x_{2}=2k^{2}/(k^{2}+2)\)，\(x_{1}x_{2}=(k^{2}-8)/(k^{2}+2)\)． \(k_{PN}+k_{QN}=k(x_{1}-1)/(x_{1}-x_{0})+k(x_{2}-1)/(x_{2}-x_{0})=k[2x_{1}x_{2}-(x_{0}+1)(x_{1}+x_{2})+2x_{0}]/[(x_{1}-x_{0})(x_{2}-x_{0})]\)， 其中分子方括号\(=[2(k^{2}-8)-(x_{0}+1)\cdot 2k^{2}+2x_{0}(k^{2}+2)]/(k^{2}+2)=(4x_{0}-16)/(k^{2}+2)=4(x_{0}-4)/(k^{2}+2)\)， 故 \(k_{PN}+k_{QN}=4k(x_{0}-4)/\{(k^{2}+2)(x_{1}-x_{0})(x_{2}-x_{0})\}\)． （分母不为 \(0\)：椭圆上点横坐标 \(|x|\leqslant 2\)，取 \(x_{0}=4\) 时 \(x\_i-x_{0}\neq 0\) ） 要对一切实数 \(k\) 成立，须 \(x_{0}-4=0\)，即 \(x_{0}=4\)． 综上，\(x\) 轴上存在定点 \(N(4,0)\)，使 \(\angle PNM=\angle QNM\) 恒成立．}
\end{ansblock}

\zux{蒙日圆族}
\tuhao{35} \tieside{中}已知椭圆 \(C\)：\(x^{2}/a^{2}+y^{2}/b^{2}=1 (a>b>0)\) 的短轴长是 \(2\)，且离心率为 \(\sqrt{3}/2\)．
\((1)\) 求椭圆 \(C\) 的方程；
\((2)\) 设直线 \(y=kx+\sqrt{3}\) 与椭圆 \(C\) 交于 \(M\)，\(N\) 两点，点 \(A(2,0)\)．问在直线 \(x=3\) 上是否存在点 \(P\)，使得四边形 \(PAMN\) 是平行四边形，若存在，求出 \(k\) 的值；若不存在，说明理由．
\liubai[24mm]{解答书写区}
\begin{ansblock}[2章-拓-课时12-35]
% ans:2章-拓-课时12-35
\ansitem{35}{\((1) x^{2}/4+y^{2}=1\)；\((2)\) 存在，\(k=\sqrt{3}/2\) 或 \(k=\pm \sqrt{11}/2\)}
\ansnote{详解}{\((1) 2b=2 \Rightarrow  b=1\)；\(e=c/a=\sqrt{3}/2\) 且 \(a^{2}=b^{2}+c^{2} \Rightarrow  c^{2}=3a^{2}/4 \Rightarrow  a^{2}/4=1 \Rightarrow  a^{2}=4\)、\(c=\sqrt{3}\)，\(C\)：\(x^{2}/4+y^{2}=1\)． \((2)\) 若 \(PAMN\) 为平行四边形，则 \(PA\parallel MN\) 且 \(|PA|=|MN|\)．由 \(PA\parallel MN\) 且过 \(A(2,0)\)，\(PA\)：\(y=k(x-2)\)，与 \(x=3\) 交于 \(P(3,k)\)，故 \(|PA|=\sqrt{1+k^{2}}\)． 联立 \(y=kx+\sqrt{3}\) 与 \(x^{2}+4y^{2}=4\)：\((4k^{2}+1)x^{2}+8\sqrt{3}kx+8=0\)，\(\Delta =192k^{2}-32(4k^{2}+1)=64k^{2}-32>0 \Leftrightarrow  k^{2}>1/2\)， \(x_{1}+x_{2}=-8\sqrt{3}k/(4k^{2}+1)\)，\(x_{1}x_{2}=8/(4k^{2}+1)\)， \(|MN|=\sqrt{1+k^{2}}|x_{1}-x_{2}|=\sqrt{1+k^{2}}\cdot \sqrt{64k^{2}-32}/(4k^{2}+1)\)． \(|PA|=|MN| \Rightarrow  \sqrt{64k^{2}-32}=4k^{2}+1\)，平方：\(64k^{2}-32=16k^{4}+8k^{2}+1 \Rightarrow  16k^{4}-56k^{2}+33=0\)， 即 \((4k^{2}-3)(4k^{2}-11)=0 \Rightarrow  k^{2}=3/4\) 或 \(k^{2}=11/4\)，得 \(k=\pm \sqrt{3}/2\)、\(k=\pm \sqrt{11}/2\)（均满足 \(k^{2}>1/2\) ）． 逐一检验四边形是否退化：\(k=-\sqrt{3}/2\) 时直线 \(MN\) 为 \(y=-(\sqrt{3}/2)x+\sqrt{3}\)，令 \(x=2\) 得 \(y=0\)，即直线 \(MN\) 过 \(A(2,0)\)，此时 \(M\)、\(N\) 中有一点与 \(A\) 重合（解 \((4k^{2}+1)x^{2}+8\sqrt{3}kx+8=0\) 得 \(x=1\) 或 \(x=2\)），\(A\)、\(M\)、\(N\) 三点共线，不能构成平行四边形，舍去； 其余三值：\(k=\sqrt{3}/2\) 时 \(MN\) 不过 \(A\)（\(x=2\) 处 \(y=2\sqrt{3}\neq 0\)），\(k=\pm \sqrt{11}/2\) 时 \(x=2\) 处 \(y=\pm \sqrt{11}+\sqrt{3}\neq 0\)，且 \(|PA|=|MN|\)、\(PA\parallel MN\) 同时成立，命名 \(M\)、\(N\) 使 向量\(PA=\)向量\(NM\) 即得平行四边形，均符合． 故 \(k=\sqrt{3}/2\) 或 \(k=\pm \sqrt{11}/2\)．}
\end{ansblock}

\tuhao{36} \tieside{中}已知点 \(Q\) 是圆 \(M\)：\((x+1)^{2}+y^{2}=64\)（圆心为 \(M\)）上的动点，点 \(N(1,0)\)，线段 \(QN\) 的垂直平分线交 \(MQ\) 于点 \(P\)．
\((1)\) 求动点 \(P\) 的轨迹 \(E\) 的方程；
\((2)\) 已知点 \(B(2,2)\)，\(S\) 是轨迹 \(E\) 上一动点，求 \(|SN|+|SB|\) 的最大值；
\((3)\) 在轨迹 \(E\) 上是否存在点 \(T\)，使 \(2/|MN|=1/|TM|+1/|TN|\)？若存在，求出 \(|TM|\) 与 \(|TN|\) 的值；若不存在，请说明理由．
\liubai[32mm]{解答书写区}
\begin{ansblock}[2章-拓-课时12-36]
% ans:2章-拓-课时12-36
\ansitem{36}{\((1) x^{2}/16+y^{2}/15=1\)；\((2) 8+\sqrt{13}\)；\((3)\) 不存在}
\ansnote{详解}{\((1)\) 圆心 \(M(-1,0)\)、半径 \(8\)，\(N(1,0)\) 在圆内（\(|MN|=2<8\)）．\(P\) 在 \(QN\) 的垂直平分线上 \(\Rightarrow  |PN|=|PQ|\)；又 \(P\) 在半径 \(MQ\) 上 \(\Rightarrow  |PM|+|PQ|=8\)， 故 \(|PM|+|PN|=|PM|+|PQ|=8>|MN|=2\)，由椭圆定义，\(P\) 的轨迹是以 \(M\)、\(N\) 为焦点、长轴长 \(8\) 的椭圆：\(2a=8\)、\(2c=2 \Rightarrow  a=4\)、\(c=1\)、\(b^{2}=15\)， \(E\)：\(x^{2}/16+y^{2}/15=1\)． \((2) N\) 为右焦点，\(M\) 为左焦点，\(|SN|+|SM|=2a=8 \Rightarrow  |SN|=8-|SM|\)， \(|SN|+|SB|=8+(|SB|-|SM|)\leqslant 8+|MB|\)（三角形两边之差小于第三边）， \(|MB|=\sqrt{}((2+1)^{2}+2^{2})=\sqrt{13}\)； 等号当 \(M\) 在线段 \(SB\) 上，即 \(S\) 在射线 \(BM\) 自 \(M\) 向外的部分上——\(M(-1,0)\) 在椭圆内（\(1/16<1\)），该射线与椭圆有唯一交点 \(S\)，故等号可取到，最大值为 \(8+\sqrt{13}\)． \((3)\) 不存在．理由：\(|MN|=2\)，若存在 \(T\) 使 \(2/|MN|=1/|TM|+1/|TN|\)，即 \(1=(|TM|+|TN|)/(|TM|\cdot |TN|)\)， 由定义 \(|TM|+|TN|=8\)，得 \(|TM|\cdot |TN|=8\)，故 \(|TM|\)、\(|TN|\) 是 \(u^{2}-8u+8=0\) 的两根，\(u=4\pm 2\sqrt{2}\)； 另一方面，\(T\) 在椭圆上，焦半径 \(|TM|=a+ex\_T\)（\(e=1/4\)）\(\in [a-c, a+c]=[3,5]\)，同理 \(|TN|\in [3,5]\)； 但 \(4-2\sqrt{2}\approx 1.17<3\)，\(4+2\sqrt{2}\approx 6.83>5\)，均落在 \([3,5]\) 之外（亦可从乘积看：\(u\in [3,5]\) 时 \(u(8-u)\) 在 \([3,4]\) 上递增、\([4,5]\) 上递减，取值 \([15,16]\)，恒大于 \(8\)）， 故满足条件的点 \(T\) 不存在．}
\end{ansblock}

% 图债注：「如图」指源件配图，印面无图（冻片【注】裁定：去装饰图/条件全在文字/尺寸随注——图债登记汇编报告）
\tuhao{37} \tieside{简}如图，过原点 \(O\) 的直线交椭圆 \(x^{2}/4+y^{2}/2=1\) 于 \(P\)，\(A\) 两点，其中点 \(P\) 在第一象限，过点 \(P\) 作 \(x\) 轴的垂线，垂足为 \(C\)，连接 \(AC\) 并延长，交椭圆于另一点 \(B\)，求证：\(k_{PA}\cdot k_{PB}\) 为定值．
\liubai[16mm]{解答书写区}
\begin{ansblock}[2章-拓-课时12-37]
% ans:2章-拓-课时12-37
\ansitem{37}{定值为 \(-1\)}
\ansnote{详解}{设 \(P(x_{1},y_{1})\)（\(x_{1}>0\)、\(y_{1}>0\)）、\(B(x_{2},y_{2})\)，由 \(P\)、\(A\) 关于原点对称得 \(A(-x_{1},-y_{1})\)，又 \(C(x_{1},0)\)． \(A\)、\(B\)、\(P\) 均在椭圆上，即 \(x^{2}+2y^{2}=4\)，故 \(y\_i^{2}=2-x\_i^{2}/2\)（\(i=1,2\)）． 因 \(A\)、\(C\)、\(B\) 三点共线（\(B\) 在 \(AC\) 的延长线上），\(k_{AB}=k_{AC}\)，而 \(k_{AB}\cdot k_{PB}=[(y_{2}+y_{1})/(x_{2}+x_{1})]\cdot [(y_{2}-y_{1})/(x_{2}-x_{1})]=(y_{2}^{2}-y_{1}^{2})/(x_{2}^{2}-x_{1}^{2})=[(2-x_{2}^{2}/2)-(2-x_{1}^{2}/2)]/(x_{2}^{2}-x_{1}^{2})=-1/2\)， （此即第三定义：椭圆上两点与关于原点对称的两点连线斜率之积\(=(b^{2}/a^{2})-1=2/4-1=-1/2\) 的变式，此处 \(x^{2}/4+y^{2}/2=1\) 的 \(a^{2}=4\)、\(b^{2}=2\)） 又 \(k_{AC}=(0+y_{1})/(x_{1}+x_{1})=y_{1}/(2x_{1})\)，\(k_{PA}=(y_{1}+y_{1})/(x_{1}+x_{1})=y_{1}/x_{1}\)，故 \(k_{PA}=2k_{AC}=2k_{AB}\)， 所以 \(k_{PA}\cdot k_{PB}=2k_{AB}\cdot k_{PB}=2\times (-1/2)=-1\)，为定值（与 \(P\) 的位置无关）．}
\end{ansblock}

\tuhao{38} \tieside{中}已知椭圆 \(E\)：\(x^{2}/a^{2}+y^{2}/b^{2}=1 (a>b>0)\) 的右焦点坐标为 \(F(3,0)\)，过 \(F\) 的直线 \(l\) 交椭圆于 \(A\)，\(B\) 两点，当 \(A\) 与上顶点重合时，向量\(AF=3\)向量\(FB\)．
\((1)\) 求椭圆 \(E\) 的方程；
\((2)\) 若点 \(P(6,0)\)，记直线 \(PA\)，\(PB\) 的斜率分别为 \(k_{1}\)、\(k_{2}\)，证明：\(k_{1}+k_{2}\) 为定值．
\liubai[24mm]{解答书写区}
\begin{ansblock}[2章-拓-课时12-38]
% ans:2章-拓-课时12-38
\ansitem{38}{\((1) x^{2}/18+y^{2}/9=1\)；\((2) k_{1}+k_{2}\) 为定值 \(0\)，证明见详解}
\ansnote{详解}{\((1) c=3\)，\(A(0,b)\)．向量\(AF=(3,-b)\)，由向量\(AF=3\)向量\(FB\) 得向量\(FB=(1,-b/3)\)，故 \(B=F+\)向量\(FB=(4,-b/3)\)． \(B\) 在椭圆上：\(16/a^{2}+(b^{2}/9)/b^{2}=1 \Rightarrow  16/a^{2}+1/9=1 \Rightarrow  a^{2}=18\)，\(b^{2}=a^{2}-c^{2}=18-9=9\)， \(E\)：\(x^{2}/18+y^{2}/9=1\)（即 \(x^{2}+2y^{2}=18\)）． \((2)\) 当 \(l\) 不为 \(x\) 轴时，设 \(l\)：\(x=my+3\)（此式可表示除 \(x\) 轴外一切过 \(F\) 的直线，其中 \(m=0\) 即垂直于 \(x\) 轴的直线 \(x=3\)），设 \(A(x_{1},y_{1})\)、\(B(x_{2},y_{2})\)， 代入 \(x^{2}+2y^{2}=18\)：\((m^{2}+2)y^{2}+6my-9=0\)（\(\Delta =36m^{2}+36(m^{2}+2)>0\) 恒成立），\(y_{1}+y_{2}=-6m/(m^{2}+2)\)，\(y_{1}y_{2}=-9/(m^{2}+2)\)． \(k_{1}+k_{2}=y_{1}/(x_{1}-6)+y_{2}/(x_{2}-6)=y_{1}/(my_{1}-3)+y_{2}/(my_{2}-3)=[2my_{1}y_{2}-3(y_{1}+y_{2})]/[(my_{1}-3)(my_{2}-3)]\)， 分子\(=2m\cdot (-9/(m^{2}+2))-3\cdot (-6m/(m^{2}+2))=(-18m+18m)/(m^{2}+2)=0\)， 分母\(=m^{2}y_{1}y_{2}-3m(y_{1}+y_{2})+9=-9m^{2}/(m^{2}+2)+18m^{2}/(m^{2}+2)+9=9m^{2}/(m^{2}+2)+9=(18m^{2}+18)/(m^{2}+2)\neq 0\)， 故 \(k_{1}+k_{2}=0\)． 当 \(l\) 为 \(x\) 轴（\(y=0\)）时，\(A\)、\(B\) 为左右顶点 \((\pm 3\sqrt{2},0)\)，\(k_{1}=k_{2}=0\)，仍有 \(k_{1}+k_{2}=0\)． （另核验：\(P(6,0)\) 在椭圆外（\(36/18=2>1\)），且 \(6\notin [-3\sqrt{2},3\sqrt{2}]\)，故 \(x\_i-6\neq 0\)，斜率总有意义．） 所以 \(k_{1}+k_{2}\) 为定值 \(0\)．}
\end{ansblock}

\tuhao{39} \tieside{中}已知椭圆 \(C\)：\(x^{2}/a^{2}+y^{2}/b^{2}=1\) 过 \(A(2,0)\)，\(B(0,1)\) 两点．
\((1)\) 求椭圆 \(C\) 的方程；
\((2)\) 设 \(P\) 为第三象限内一点且在椭圆 \(C\) 上，直线 \(PA\) 与 \(y\) 轴交于点 \(M\)，直线 \(PB\) 与 \(x\) 轴交于点 \(N\)，求证：四边形 \(ABNM\) 的面积为定值．
\liubai[24mm]{解答书写区}
\begin{ansblock}[2章-拓-课时12-39]
% ans:2章-拓-课时12-39
\ansitem{39}{\((1) x^{2}/4+y^{2}=1\)；\((2)\) 定值为 \(2\)，证明见详解}
\ansnote{详解}{\((1) A(2,0)\) 在 \(x\) 轴上、\(B(0,1)\) 在 \(y\) 轴上，故 \(a=2\)、\(b=1\)（\(a>b>0\) 满足），\(C\)：\(x^{2}/4+y^{2}=1\)． \((2)\) 设 \(P(x_{0},y_{0})\)，第三象限故 \(x_{0}\in (-2,0)\)、\(y_{0}\in (-1,0)\)，且 \(x_{0}^{2}+4y_{0}^{2}=4\)． 直线 \(PA\)：\(y=y_{0}/(x_{0}-2)\cdot (x-2)\)，令 \(x=0\) 得 \(M(0, -2y_{0}/(x_{0}-2))\)；直线 \(PB\)：\(y-1=(y_{0}-1)/x_{0}\cdot x\)，令 \(y=0\) 得 \(N(-x_{0}/(y_{0}-1), 0)\)． 因 \(x_{0}-2<0\)、\(-2y_{0}>0\)，\(M\) 的纵坐标\(<0\)；\(y_{0}-1<0\)、\(-x_{0}>0\)，\(N\) 的横坐标\(<0\)，即 \(M\) 在 \(y\) 轴负半轴、\(N\) 在 \(x\) 轴负半轴． 四边形 \(ABNM\) 的顶点依次为 \(A(2,0)\)、\(B(0,1)\)、\(N(\)负\(,0)\)、\(M(0,\)负\()\)，两条对角线 \(AN\) 在 \(x\) 轴上、\(BM\) 在 \(y\) 轴上，互相垂直且相交于原点， 故 \(S=(1/2)|AN|\cdot |BM|=(1/2)\cdot |2+x_{0}/(y_{0}-1)|\cdot |1+2y_{0}/(x_{0}-2)| =(1/2)\cdot |(x_{0}+2y_{0}-2)/(y_{0}-1)|\cdot |(x_{0}+2y_{0}-2)/(x_{0}-2)|=(1/2)\cdot |(x_{0}+2y_{0}-2)^{2}/[(y_{0}-1)(x_{0}-2)]|\)． 分子展开并用 \(x_{0}^{2}+4y_{0}^{2}=4\)：\((x_{0}+2y_{0}-2)^{2}=x_{0}^{2}+4y_{0}^{2}+4+4x_{0}y_{0}-4x_{0}-8y_{0}=8+4x_{0}y_{0}-4x_{0}-8y_{0}=4(x_{0}y_{0}-x_{0}-2y_{0}+2)=4(y_{0}-1)(x_{0}-2)\)， 所以 \(S=(1/2)\cdot |4|=2\)，与 \(P\) 的位置无关，即四边形 \(ABNM\) 的面积为定值 \(2\)．}
\end{ansblock}

\tuhao{40} \tieside{中}设椭圆 \(E\) 的方程为 \(x^{2}/a^{2}+y^{2}/b^{2}=1 (a>b>0)\)，椭圆 \(E\) 的短轴长为 \(2\)，离心率为 \(\sqrt{3}/2\)．
\((1)\) 求椭圆 \(E\) 的方程；
\((2)\) 设 \(A\)、\(B\) 分别为椭圆 \(E\) 的左、右顶点，过定点 \(T(1,0)\) 的直线与椭圆 \(E\) 交于 \(C\)、\(D\) 两点，证明：直线 \(AC\)，\(BD\) 的交点在定直线上．
\liubai[24mm]{解答书写区}
\begin{ansblock}[2章-拓-课时12-40]
% ans:2章-拓-课时12-40
\ansitem{40}{\((1) x^{2}/4+y^{2}=1\)；\((2)\) 证明见详解，定直线为 \(x=4\)}
\ansnote{详解}{\((1) 2b=2 \Rightarrow  b=1\)；\(e=c/a=\sqrt{3}/2\) 且 \(a^{2}=b^{2}+c^{2} \Rightarrow  a^{2}=4\)、\(c=\sqrt{3}\)，\(E\)：\(x^{2}/4+y^{2}=1\)． \((2) A(-2,0)\)、\(B(2,0)\)，设 \(C(x_{1},y_{1})\)、\(D(x_{2},y_{2})\)．过 \(T\) 的直线不能为 \(x\) 轴（否则 \(C\)、\(D\) 即两顶点，\(AC\)、\(BD\) 均退化为 \(x\) 轴），故设为 \(x=my+1\)（\(m=0\) 时为直线 \(x=1\)）． 代入 \(x^{2}+4y^{2}=4\)：\((m^{2}+4)y^{2}+2my-3=0\)，\(\Delta =4m^{2}+12(m^{2}+4)>0\) 恒成立，\(y_{1}+y_{2}=-2m/(m^{2}+4)\)，\(y_{1}y_{2}=-3/(m^{2}+4)\)① 直线 \(AC\)：\(y=y_{1}/(x_{1}+2)\cdot (x+2)\)；直线 \(BD\)：\(y=y_{2}/(x_{2}-2)\cdot (x-2)\)。设交点为 \(Q(x,y)\)，两式相除得 \((x+2)/(x-2)=y_{2}(x_{1}+2)/[y_{1}(x_{2}-2)]=y_{2}(my_{1}+3)/[y_{1}(my_{2}-1)]=(my_{1}y_{2}+3y_{2})/(my_{1}y_{2}-y_{1})\)． 由①：\(my_{1}y_{2}=-3m/(m^{2}+4)=(3/2)\cdot (-2m/(m^{2}+4))=(3/2)(y_{1}+y_{2})\)， 故分子\(=(3/2)(y_{1}+y_{2})+3y_{2}=(3/2)(y_{1}+3y_{2})\)，分母\(=(3/2)(y_{1}+y_{2})-y_{1}=(1/2)(y_{1}+3y_{2})\)， （\(y_{1}+3y_{2}=0\) 不可能：若 \(y_{1}=-3y_{2}\)，则 \(y_{1}+y_{2}=-2y_{2}=-2m/(m^{2}+4)\)、\(y_{1}y_{2}=-3y_{2}^{2}=-3/(m^{2}+4)\)，两式相除得 \(y_{2}=m\)，代回得 \(m^{2}=m^{2}+4\)，矛盾 ） 所以 \((x+2)/(x-2)=3 \Rightarrow  x+2=3x-6 \Rightarrow  x=4\)， 即直线 \(AC\) 与 \(BD\) 的交点横坐标恒为 \(4\)，交点恒在定直线 \(x=4\) 上．}
\end{ansblock}

\tuhao{41} \tieside{中}已知椭圆 \(C\)：\(x^{2}/a^{2}+y^{2}/b^{2}=1\)（\(a>b>0\)）的左、右焦点分别为 \(F_{1}(-c,0)\)、\(F_{2}(c,0)\)，\(M\)，\(N\) 分别为左、右顶点，直线 \(l\)：\(x=ty+1\) 与椭圆 \(C\) 交于 \(A\)，\(B\) 两点，当 \(t=-\sqrt{3}/3\) 时，\(A\) 是椭圆的上顶点，且 \(\triangle AF_{1}F_{2}\) 的周长为 \(6\)．
\((1)\) 求椭圆 \(C\) 的方程；
\((2)\) 设直线 \(AM\)，\(BN\) 交于点 \(Q\)，证明：点 \(Q\) 在定直线上．
\liubai[24mm]{解答书写区}
\begin{ansblock}[2章-拓-课时12-41]
% ans:2章-拓-课时12-41
\ansitem{41}{\((1) x^{2}/4+y^{2}/3=1\)；\((2)\) 证明见详解，定直线为 \(x=4\)}
\ansnote{详解}{\((1) t=-\sqrt{3}/3\) 时 \(l\)：\(x=-(\sqrt{3}/3)y+1\)，令 \(x=0\) 得 \(y=\sqrt{3}\)，即上顶点为 \((0,\sqrt{3})\)，故 \(b=\sqrt{3}\)； \(A\) 为上顶点时 \(\triangle AF_{1}F_{2}\) 的周长\(=|AF_{1}|+|AF_{2}|+|F_{1}F_{2}|=2a+2c=6\)，即 \(a+c=3\)；又 \(a^{2}=b^{2}+c^{2}=3+c^{2}\)， 联立 \(a=3-c\)：\((3-c)^{2}=3+c^{2} \Rightarrow  9-6c=3 \Rightarrow  c=1\)、\(a=2\)，\(C\)：\(x^{2}/4+y^{2}/3=1\)． \((2)\) 直线 \(l\)：\(x=ty+1\) 恒过 \((1,0)=F_{2}\)，且不为 \(x\) 轴（\(x\) 轴与椭圆交于两顶点，此时 \(AM\)、\(BN\) 退化），设 \(A(x_{1},y_{1})\)、\(B(x_{2},y_{2})\)， 代入 \(3x^{2}+4y^{2}=12\)：\((3t^{2}+4)y^{2}+6ty-9=0\)（\(\Delta =36t^{2}+36(3t^{2}+4)>0\) 恒成立），\(y_{1}+y_{2}=-6t/(3t^{2}+4)\)，\(y_{1}y_{2}=-9/(3t^{2}+4)\)． \(M(-2,0)\)、\(N(2,0)\)，直线 \(AM\)：\(y=y_{1}/(x_{1}+2)\cdot (x+2)\)，直线 \(BN\)：\(y=y_{2}/(x_{2}-2)\cdot (x-2)\)，交点 \(Q(x,y)\) 满足 \((x+2)/(x-2)=y_{2}(x_{1}+2)/[y_{1}(x_{2}-2)]=y_{2}(ty_{1}+3)/[y_{1}(ty_{2}-1)]=(ty_{1}y_{2}+3y_{2})/(ty_{1}y_{2}-y_{1})\)， 由 \(ty_{1}y_{2}=-9t/(3t^{2}+4)=(3/2)(y_{1}+y_{2})\)，得分子\(=(3/2)(y_{1}+y_{2})+3y_{2}=(3/2)(y_{1}+3y_{2})\)、分母\(=(3/2)(y_{1}+y_{2})-y_{1}=(1/2)(y_{1}+3y_{2})\)， （\(y_{1}+3y_{2}=0\) 时同 \(12-\)拓\(40\) 由 \(y_{1}y_{2}\) 与 \(y_{1}+y_{2}\) 之比会推出矛盾，此处 \(3t^{2}=3t^{2}+4\) 不可能 ） 故 \((x+2)/(x-2)=3\)，解得 \(x=4\)，即点 \(Q\) 恒在定直线 \(x=4\) 上．}
\end{ansblock}

% 图债注：「如图」指源件配图，印面无图（冻片【注】裁定：去装饰图/条件全在文字/尺寸随注——图债登记汇编报告）
\tuhao{42} \tieside{中}如图，椭圆 \(C\)：\(x^{2}/a^{2}+y^{2}/b^{2}=1 (a>b>0)\) 的左、右焦点分别为 \(F_{1}\)，\(F_{2}\)，上顶点为 \(A\)，过点 \(A\) 与 \(AF_{2}\) 垂直的直线交 \(x\) 轴负半轴于点 \(Q\)，且 \(F_{1}\) 恰是 \(QF_{2}\) 的中点，若过 \(A\)，\(Q\)，\(F_{2}\) 三点的圆与直线 \(l\)：\(x-\sqrt{3}y-3=0\) 相切．
\((1)\) 求椭圆 \(C\) 的方程；
\((2)\) 设 \(M\)，\(N\) 为椭圆 \(C\) 的长轴两端点，直线 \(m\) 过点 \(P(4,0)\) 交 \(C\) 于不同两点 \(G\)，\(H\)，证明：四边形 \(MNHG\) 的对角线交点在定直线上，并求出定直线方程．
\liubai[24mm]{解答书写区}
\begin{ansblock}[2章-拓-课时12-42]
% ans:2章-拓-课时12-42
\ansitem{42}{\((1) x^{2}/4+y^{2}/3=1\)；\((2)\) 证明见详解，定直线为 \(x=1\)}
\ansnote{详解}{\((1)\) 设半焦距为 \(c\)，\(A(0,b)\)、\(F_{2}(c,0)\)、\(F_{1}(-c,0)\)、\(Q(q,0)\)（\(q<0\)）． \(AQ\perp AF_{2} \Rightarrow \) 向量\(AQ\cdot \)向量\(AF_{2}=(q,-b)\cdot (c,-b)=qc+b^{2}=0 \Rightarrow  q=-b^{2}/c\)； \(F_{1}\) 为 \(QF_{2}\) 的中点 \(\Rightarrow  (q+c)/2=-c \Rightarrow  q=-3c\)，故 \(b^{2}=3c^{2}\)，\(a^{2}=b^{2}+c^{2}=4c^{2}\)，即 \(a=2c\)，且 \(|QF_{2}|=q\) 到 \(c\) 的距离\(=4c\)， \(\triangle QAF_{2}\) 中 \(\angle QAF_{2}=90^{\circ}\)，故过 \(A\)、\(Q\)、\(F_{2}\) 三点的圆以 \(QF_{2}\) 为直径，圆心为 \(F_{1}(-c,0)\)、半径 \(2c\)； 该圆与 \(l\)：\(x-\sqrt{3}y-3=0\) 相切 \(\Rightarrow  |-c-3|/\sqrt{1+3}=2c \Rightarrow  (c+3)/2=2c\)（\(c>0\)）\(\Rightarrow  c=1\)， 于是 \(a=2\)、\(b^{2}=3\)，\(C\)：\(x^{2}/4+y^{2}/3=1\)． \((2)\) 直线 \(m\) 过 \(P(4,0)\) 且与椭圆交于两点，不能为 \(x\) 轴（否则 \(G\)、\(H\) 即长轴两端点，四边形退化），设 \(m\)：\(x=my+4\)， 代入 \(3x^{2}+4y^{2}=12\)：\((3m^{2}+4)y^{2}+24my+36=0\)，须 \(\Delta =576m^{2}-144(3m^{2}+4)=144(m^{2}-4)>0\)，即 \(m^{2}>4\)， \(y_{1}+y_{2}=-24m/(3m^{2}+4)\)，\(y_{1}y_{2}=36/(3m^{2}+4)\)，由此得 \(m y_{1}y_{2}=-(3/2)(y_{1}+y_{2})\)． 设 \(H(x_{1},y_{1})\)、\(G(x_{2},y_{2})\)，\(M(-2,0)\)、\(N(2,0)\)，四边形 \(MNHG\) 的两条对角线为 \(MH\) 与 \(NG\)： \(MH\)：\(y=y_{1}/(x_{1}+2)\cdot (x+2)\)，\(NG\)：\(y=y_{2}/(x_{2}-2)\cdot (x-2)\)，交点横坐标 \(x=2\cdot [y_{1}(x_{2}-2)+y_{2}(x_{1}+2)]/[y_{2}(x_{1}+2)-y_{1}(x_{2}-2)]\)， 由 \(x_{1}+2=my_{1}+6\)、\(x_{2}-2=my_{2}+2\)，分子方括号\(=y_{1}(my_{2}+2)+y_{2}(my_{1}+6)=2my_{1}y_{2}+2y_{1}+6y_{2}\)， 分母方括号\(=y_{2}(my_{1}+6)-y_{1}(my_{2}+2)=6y_{2}-2y_{1}=2(3y_{2}-y_{1})\)， 故 \(x=2\cdot (2my_{1}y_{2}+2y_{1}+6y_{2})/[2(3y_{2}-y_{1})]=2\cdot (my_{1}y_{2}+y_{1}+3y_{2})/(3y_{2}-y_{1})=2\cdot [-(3/2)(y_{1}+y_{2})+y_{1}+3y_{2}]/(3y_{2}-y_{1})=2\cdot [(3/2)y_{2}-(1/2)y_{1}]/(3y_{2}-y_{1})=2\cdot (1/2)=1\)． （分母 \(3y_{2}-y_{1}=0\) 不可能：若 \(y_{1}=3y_{2}\)，由 \(y_{1}+y_{2}=4y_{2}=-24m/(3m^{2}+4)\) 得 \(y_{2}=-6m/(3m^{2}+4)\)，而 \(y_{1}y_{2}=3y_{2}^{2}=36/(3m^{2}+4)\) 得 \(y_{2}^{2}=12/(3m^{2}+4)\)，两式平方相较得 \(36m^{2}=12(3m^{2}+4)\)，即 \(0=48\)，矛盾，故两对角线恒相交、交点总存在 ） 所以对角线交点的横坐标恒为 \(1\)，即四边形 \(MNHG\) 的对角线交点在定直线 \(x=1\) 上．}
\end{ansblock}

\jietitle{2.6.1　双曲线的标准方程（课时13·拓区）}
% 图债注：「如图」指源件配图，印面无图（冻片【注】裁定：去装饰图/条件全在文字/尺寸随注——图债登记汇编报告）
\tuhao{01} \tieside{中}如图，已知动圆 \(M\) 与两个定圆 \(\odot O_{1}\) 和 \(\odot O_{2}\) 分别外切（\(\odot O_{1}\)、\(\odot O_{2}\) 外离），则动圆圆心 \(M\) 的轨迹是什么图形？
【图注】\(\odot O_{1}\)、\(\odot O_{2}\) 半径可不等，两圆外离（圆心距大于两半径之和）。
\liubai[16mm]{解答书写区}
\begin{ansblock}[2章-拓-课时13-01]
% ans:2章-拓-课时13-01
\ansitem{01}{见详解（按半径大小分三类：等半径\(\rightarrow \)线段 \(O_{1}O_{2}\) 的垂直平分线；不等\(\rightarrow \)以 \(O_{1}\)、\(O_{2}\) 为焦点的双曲线的一支）}
\ansnote{详解}{设 \(\odot O_{1}\)、\(\odot O_{2}\) 的半径为 \(R_{1}\)、\(R_{2}\)，动圆 \(M\) 的半径为 \(r\)． 动圆与两定圆都外切 \(\Rightarrow  |MO_{1}|=r+R_{1}\)、\(|MO_{2}|=r+R_{2}\)，两式相减消去动半径 \(r\)： 若 \(R_{1}=R_{2}\)，则 \(|MO_{1}|=|MO_{2}|\)，\(M\) 到两定点等距，轨迹是线段 \(O_{1}O_{2}\) 的垂直平分线； 若 \(R_{1}>R_{2}\)，则 \(|MO_{1}|-|MO_{2}|=R_{1}-R_{2}>0\)，且 \(R_{1}-R_{2}<R_{1}+R_{2}<|O_{1}O_{2}|\)（外离条件）， 故 \(M\) 的轨迹是以 \(O_{1}\)、\(O_{2}\) 为焦点的双曲线中靠近 \(O_{2}\) 的一支（即远离大圆一侧的支），因差恒为正，只含一支； 若 \(R_{2}>R_{1}\)，同理 \(|MO_{2}|-|MO_{1}|=R_{2}-R_{1}\in (0,|O_{1}O_{2}|)\)，轨迹是以 \(O_{1}\)、\(O_{2}\) 为焦点的双曲线的另一支． （三支并论：后两类合起来即\("\)以 \(O_{1}\)、\(O_{2}\) 为焦点的双曲线\("\)，但每种半径关系下只取其中一支，故本题须按 \(R_{1}\)、\(R_{2}\) 大小分类作答，不能直接答\("\)双曲线\("\)。）}
\end{ansblock}

\tuhao{02} \tieside{中}已知椭圆 \(C_{1}\)：\(x^{2}/a_{1}^{2}+y^{2}/b_{1}^{2}=1\)（\(a_{1}>b_{1}>0\)）与双曲线 \(C_{2}\)：\(x^{2}/a_{2}^{2}-y^{2}/b_{2}^{2}=1\)（\(a_{2}>0\)，\(b_{2}>0\)）有公共焦点 \(F_{1}\)，\(F_{2}\)，且两条曲线在第一象限的交点为 \(P\)．若 \(\triangle PF_{1}F_{2}\) 是以 \(PF_{1}\) 为底边的等腰三角形，曲线 \(C_{1}\)、\(C_{2}\) 的离心率分别为 \(e_{1}\) 和 \(e_{2}\)，则 \(1/e_{1} - 1/e_{2} =\)（　　）
\bindopt {\fontsize{10.5pt}{\qpTJgLeadLiti}\selectfont \makebox[\dimexpr0.25\linewidth-0.25\lxhang-0.25em\relax][l]{\(A\)．\(1\)}\makebox[\dimexpr0.25\linewidth-0.25\lxhang-0.25em\relax][l]{\(B\)．\(2\)}\makebox[\dimexpr0.25\linewidth-0.25\lxhang-0.25em\relax][l]{\(C\)．\(3\)}\makebox[\dimexpr0.25\linewidth-0.25\lxhang-0.25em\relax][l]{\(D\)．\(4\)}\par}
\begin{ansblock}[2章-拓-课时13-02]
% ans:2章-拓-课时13-02
\ansitem{02}{\(B\)}
\ansnote{详解}{设公共半焦距为 \(c\)（焦距 \(2c\)），\(|F_{1}F_{2}|=2c\)． \(\triangle PF_{1}F_{2}\) 以 \(PF_{1}\) 为底边，故两腰相等：\(|PF_{2}|=|F_{1}F_{2}|=2c\)． \(P\) 在椭圆上：\(|PF_{1}|+|PF_{2}|=2a_{1} \Rightarrow  |PF_{1}|=2a_{1}-2c\)； \(P\) 又在第一象限（位于双曲线右支，靠 \(F_{2}\) 一侧）：\(|PF_{1}|-|PF_{2}|=2a_{2} \Rightarrow  |PF_{1}|=2a_{2}+2c\)， 两式相等：\(2a_{1}-2c=2a_{2}+2c \Rightarrow  a_{1}-a_{2}=2c\)， 两边同除 \(c\)：\(a_{1}/c - a_{2}/c=2\)，而 \(a_{1}/c=1/e_{1}\)、\(a_{2}/c=1/e_{2}\)，故 \(1/e_{1} - 1/e_{2}=2\)． （相容性核验：取 \(e_{2}=2\)（双曲线 \(e_{2}>1\) ），则 \(1/e_{1}=2+1/2=5/2\)、\(e_{1}=2/5<1\) ；此时 \(a_{1}=5c/2\)、\(a_{2}=c/2\)，\(|PF_{1}|=2a_{2}+2c=3c\)、\(|PF_{2}|=2c\)、\(|F_{1}F_{2}|=2c\)，三边 \(2c\)、\(2c\)、\(3c\) 满足三角不等式  等腰三角形存在．）故选：\(B\)．}
\end{ansblock}

\tuhao{03} \tieside{中}已知直线 \(l\) 与双曲线 \(x^{2} - y^{2}/2 =1\) 交于 \(A\)，\(B\) 两点，且向量\(AB=\lambda \)向量\(OB\)（\(O\) 为坐标原点），若 \(M\) 是直线 \(x-\sqrt{2}y-3=0\) 上的一个动点，则 \(|MA|^{2}+|MB|^{2}\) 的最小值为（　　）
\bindopt {\fontsize{10.5pt}{\qpTJgLeadLiti}\selectfont \makebox[\dimexpr0.25\linewidth-0.25\lxhang-0.25em\relax][l]{\(A\)．\(12\)}\makebox[\dimexpr0.25\linewidth-0.25\lxhang-0.25em\relax][l]{\(B\)．\(6\)}\makebox[\dimexpr0.25\linewidth-0.25\lxhang-0.25em\relax][l]{\(C\)．\(16\)}\makebox[\dimexpr0.25\linewidth-0.25\lxhang-0.25em\relax][l]{\(D\)．\(8\)}\par}
\begin{ansblock}[2章-拓-课时13-03]
% ans:2章-拓-课时13-03
\ansitem{03}{\(D\)}
\ansnote{详解}{由 向量\(AB=\lambda \)向量\(OB\) 得 向量\(OB-\)向量\(OA=\lambda \)向量\(OB\)，即 向量\(OA=(1-\lambda )\)向量\(OB\)，故 \(O\)、\(A\)、\(B\) 三点共线，直线 \(l\) 过原点 \(O\)． 双曲线 \(x^{2}-y^{2}/2=1\) 关于原点中心对称，过原点的直线与它交于 \(A\)、\(B\) 两点时，\(O\) 恰为弦 \(AB\) 的中点． 对任意点 \(M\)，由中线（阿波罗尼奥斯）恒等式（或平行四边形恒等式）：\(|MA|^{2}+|MB|^{2}=2|MO|^{2}+2|AO|^{2}=2|MO|^{2}+(1/2)|AB|^{2}\)． 两项可各自取最小： ① \(|MO|\) 的最小值为原点 \(O\) 到直线 \(x-\sqrt{2}y-3=0\) 的距离 \(d=|-3|/\sqrt{1+2}=3/\sqrt{3}=\sqrt{3}\)，此时 \(2|MO|^{2}=2\times 3=6\)； ② \(|AB|=2|AO|\)，而 \(A\) 在双曲线上、\(O\) 为对称中心，\(|AO|\) 的最小值为实半轴 \(a=1\)（\(A\) 取顶点 \((\pm 1,0)\)，此时直线 \(l\) 为 \(x\) 轴，与双曲线确实交于 \((\pm 1,0)\) 两点 ），故 \(|AB|min=2\)，\((1/2)|AB|^{2}=(1/2)\times 4=2\)； 两最小可同时达到（\(M\) 取 \(O\) 到已知直线的垂足，\(l\) 取 \(x\) 轴，二者互不牵制）， 所以 \(|MA|^{2}+|MB|^{2}\) 的最小值为 \(6+2=8\)．故选：\(D\)．}
\end{ansblock}

% 图债注：「如图」指源件配图，印面无图（冻片【注】裁定：去装饰图/条件全在文字/尺寸随注——图债登记汇编报告）
\tuhao{04} \tieside{简}如图，双曲线 \(C\)：\(x^{2}/9 - y^{2}/10 =1\) 的左焦点为 \(F_{1}\)，双曲线上的点 \(P_{1}\) 与 \(P_{2}\) 关于 \(y\) 轴对称，则 \(|P_{2}F_{1}|-|P_{1}F_{1}|\) 的值是（　　）
\bindopt {\fontsize{10.5pt}{\qpTJgLeadLiti}\selectfont \makebox[\dimexpr0.25\linewidth-0.25\lxhang-0.25em\relax][l]{\(A\)．\(3\)}\makebox[\dimexpr0.25\linewidth-0.25\lxhang-0.25em\relax][l]{\(B\)．\(4\)}\makebox[\dimexpr0.25\linewidth-0.25\lxhang-0.25em\relax][l]{\(C\)．\(6\)}\makebox[\dimexpr0.25\linewidth-0.25\lxhang-0.25em\relax][l]{\(D\)．\(8\)}\par}
\begin{ansblock}[2章-拓-课时13-04]
% ans:2章-拓-课时13-04
\ansitem{04}{\(C\)}
\ansnote{详解}{\(a^{2}=9 \Rightarrow  a=3\)，设右焦点为 \(F_{2}\)． 双曲线关于 \(y\) 轴对称，且 \(F_{1}\) 与 \(F_{2}\) 关于 \(y\) 轴对称，故把 \(P_{1}\) 关于 \(y\) 轴反射到 \(P_{2}\) 时，其到左焦点的距离变为到右焦点的距离：\(|P_{1}F_{1}|=|P_{2}F_{2}|\)． 于是 \(|P_{2}F_{1}|-|P_{1}F_{1}|=|P_{2}F_{1}|-|P_{2}F_{2}|\)， 由图（源图 \(P_{2}\) 在右支上、\(P_{1}\) 在左支上），\(P_{2}\) 在右支 \(\Rightarrow  |P_{2}F_{1}|-|P_{2}F_{2}|=2a=6\)， 故所求值为 \(6\)．故选：\(C\)． （若图中 \(P_{2}\) 落在左支，则差为 \(-2a=-6\)；本题选项只含正值，且源图 \(P_{2}\) 标在右支，故取 \(6\)。）}
\end{ansblock}

\tuhao{05}\duoxuan\hspace{0.6em} \tieside{中}（多选题）若方程 \(x^{2}/(5-t) + y^{2}/(t-1) =1\) 所表示的曲线为 \(C\)，则下面四个命题中正确的是（　　）
\lxopt{\(A\)．若 \(1<t<5\)，则 \(C\) 为椭圆}
\lxopt{\(B\)．若 \(t<1\)，则 \(C\) 为双曲线}
\lxopt{\(C\)．若 \(C\) 为双曲线，则焦距为 \(4\)}
\lxopt{\(D\)．若 \(C\) 为焦点在 \(y\) 轴上的椭圆，则 \(3<t<5\)}
\begin{ansblock}[2章-拓-课时13-05]
% ans:2章-拓-课时13-05
\ansitem{05}{\(BD\)}
\ansnote{详解}{\(A\)：\(1<t<5\) 时两分母同正，但当 \(t=3\) 时 \(5-t=t-1=2\)，方程化为 \(x^{2}+y^{2}=2\)，是圆不是椭圆，\(A\) 错误． \(B\)：\(t<1\) 时 \(5-t>0\)、\(t-1<0\)，方程即 \(x^{2}/(5-t) - y^{2}/(1-t)=1\)，两分母均为正，表示焦点在 \(x\) 轴上的双曲线，\(B\) 正确． \(C\)：\(C\) 为双曲线 \(\Leftrightarrow  (5-t)(t-1)<0 \Leftrightarrow  t<1\) 或 \(t>5\)；此时实、虚半轴平方为 \(|5-t|\) 与 \(|t-1|\)，半焦距平方 \(c^{2}=|5-t|+|t-1|\)，随 \(t\) 变化：\(t<1\) 时 \(c^{2}=6-2t\)、\(t>5\) 时 \(c^{2}=2t-6\)，均非定值。反例 \(t=0\)：\(x^{2}/5 - y^{2}/1=1\)，\(c^{2}=6\)、焦距 \(2c=2\sqrt{6}\neq 4\)，\(C\) 错误． \(D\)：焦点在 \(y\) 轴上的椭圆 \(\Leftrightarrow  t-1>5-t>0 \Leftrightarrow  t>3\) 且 \(t<5 \Leftrightarrow  3<t<5\)，\(D\) 正确． 故选：\(BD\)．}
\end{ansblock}

\tuhao{06} \tieside{中}已知 \(F_{1}(-3,0)\)，\(F_{2}(3,0)\)，若点 \(P(x,y)\) 满足 \(|PF_{1}|-|PF_{2}|=m (m\geqslant 0)\)，则 \(P\) 点的轨迹是什么，并求点 \(P\) 的轨迹方程．
\liubai[16mm]{解答书写区}
\begin{ansblock}[2章-拓-课时13-06]
% ans:2章-拓-课时13-06
\ansitem{06}{\(m=0\)：直线 \(x=0\)；\(0<m<6\)：双曲线右支，\(x^{2}/(m^{2}/4) - y^{2}/(9-m^{2}/4)=1 (x\geqslant m/2)\)；\(m=6\)：射线 \(y=0 (x\geqslant 3)\)；\(m>6\)：点 \(P\) 不存在}
\ansnote{详解}{\(|F_{1}F_{2}|=2c=6\)． \((1) m=0\)：\(|PF_{1}|=|PF_{2}|\)，\(P\) 在 \(F_{1}F_{2}\) 的垂直平分线上，轨迹是直线，方程 \(x=0\)； \((2) 0<m<6\)：\(|PF_{1}|-|PF_{2}|=m<|F_{1}F_{2}|\) 且差为正，故 \(P\) 到 \(F_{2}\) 较近，轨迹是以 \(F_{1}\)、\(F_{2}\) 为焦点的双曲线中靠近 \(F_{2}\) 的一支，即右支； 实半轴 \(a=m/2\)、\(c=3\)、\(b^{2}=c^{2}-a^{2}=9-m^{2}/4\)（\(0<m<6\) 保证 \(b^{2}>0\) ），轨迹方程 \(x^{2}/(m^{2}/4) - y^{2}/(9-m^{2}/4)=1 (x\geqslant m/2)\)； \((3) m=6=|F_{1}F_{2}|\)：\(\|PF_{1}|-|PF_{2}|\)｜\(=|F_{1}F_{2}|\) 须 \(P\)、\(F_{1}\)、\(F_{2}\) 共线且 \(P\) 不在两焦点之间，再由 \(|PF_{1}|>|PF_{2}|\) 定 \(P\) 在 \(F_{2}\) 外侧，轨迹是射线，方程 \(y=0 (x\geqslant 3)\)； \((4) m>6\)：与 \(\|PF_{1}|-|PF_{2}|\)｜\(\leqslant |F_{1}F_{2}|=6\) 矛盾，满足条件的点 \(P\) 不存在．}
\end{ansblock}

\tailfill
\end{multicols}
\end{document}